% Reaction kinetics — Chemistry Upper Sixth — Cours signé OnBuch+
% Official MINESEC syllabus. Compiler avec : tectonic CHE04-reaction-kinetics.tex
\newcommand{\DOCMATIERE}{Chemistry}
\newcommand{\DOCNIVEAU}{Upper Sixth}
\newcommand{\DOCMODULE}{Upper Sixth — Chemistry}
\newcommand{\DOCLECON}{4}
\newcommand{\DOCTITRE}{Reaction kinetics}
% OnBuch+ — préambule « cours signés » ENRICHI (compilé avec tectonic / XeLaTeX)
% Système visuel « Pop » d'OnBuch+ : fond crème (jamais blanc), bordures encre
% épaisses, ombres « dures » décalées, titres Archivo Black en capitales.
% Version MATHÉMATIQUES : graphes (pgfplots/tikz), boîtes pédagogiques
% (définition, propriété, méthode, attention, le savais-tu, exercice résolu,
% exercice, TP, bilan), page de garde et signature OnBuch+ ancrée en bas.
%
% Macros attendues dans main.tex AVANT \input{preamble} :
%   \DOCMATIERE  \DOCNIVEAU  \DOCMODULE  \DOCLECON (numéro)  \DOCTITRE
\documentclass[11pt]{article}

\usepackage{fontspec}
\usepackage[english]{babel}
\usepackage[a4paper,margin=2cm,top=3.4cm,bottom=2.8cm,headheight=1.4cm,footskip=1.3cm]{geometry}
\usepackage{amsmath,amssymb,amsfonts}
\usepackage{mathtools} % \xmapsto, \coloneqq... souvent utilisés par les modèles
\usepackage{cancel} % simplifications barrées (bilans de forces)
% \abs{z} (module, valeur absolue) et \norm{u} (norme), utilisés par les modèles
\providecommand{\abs}{}\let\abs\relax\DeclarePairedDelimiter{\abs}{\lvert}{\rvert}
\providecommand{\norm}{}\let\norm\relax\DeclarePairedDelimiter{\norm}{\lVert}{\rVert}
\usepackage[table]{xcolor} % [table] : fournit aussi \rowcolors (alternance de lignes)
\usepackage{pagecolor}
\usepackage{tikz}
\usetikzlibrary{arrows.meta,positioning,calc,shapes.geometric,shapes.misc,shapes.arrows,decorations.pathmorphing,decorations.pathreplacing,decorations.markings,patterns,shadows,mindmap,trees,fit,backgrounds,matrix,intersections,angles,quotes}
\usepackage{pgfplots}
\usepackage[version=4]{mhchem} % équations nucléaires \ce{^{235}_{92}U}
\usepackage[european,siunitx]{circuitikz} % schémas électriques
\pgfplotsset{compat=1.18}
\usepgfplotslibrary{fillbetween,groupplots} % aires entre courbes (calcul intégral)
\usepackage{enumitem}
\usepackage{fancyhdr}
% [most] charge toutes les bibliothèques tcolorbox (listings, minted, raster,
% documentation...) et gonfle inutilement la mémoire TeX sur un document long
% (~300 boîtes) : on ne charge que ce qui est réellement utilisé.
\usepackage[skins,breakable]{tcolorbox}
\usepackage{graphicx}
\usepackage{colortbl}
\usepackage{booktabs}
\usepackage{tabularx}
\usepackage{diagbox} % case d'en-tête diagonale des tableaux à double entrée
% Colonnes extensibles centrée / gauche / droite (C, L, R), souvent utilisées par les modèles
\newcolumntype{C}{>{\centering\arraybackslash}X}
\newcolumntype{L}{>{\raggedright\arraybackslash}X}
\newcolumntype{R}{>{\raggedleft\arraybackslash}X}
\usepackage{array}
\usepackage{multicol}
\usepackage{siunitx}
% parse-numbers=false : accepte aussi \SI{2,5 \times 10^{-3}}{...}
\sisetup{locale=UK,output-decimal-marker={.},group-minimum-digits=4,parse-numbers=false}
\DeclareSIUnit\ohm{\ensuremath{\symup{\Omega}}} % U+2126 absent de la police
\DeclareSIUnit\division{div} % réglages d'oscilloscope (ms/div, V/div)
\DeclareSIUnit\tr{tr} % tours (vitesse de rotation, ex. \SI{1500}{\tr\per\minute})
\usepackage{qrcode}
% \degree : commande fréquemment utilisée par les modèles pour un angle ou une
% température en dehors de \SI{}{} (non fournie par les paquets chargés).
\providecommand{\degree}{^{\circ}}
% \qed (fin de démonstration), utilisé par les modèles sans amsthm
\providecommand{\qed}{\hfill$\square$}
% \barre : double barre des valeurs interdites dans un tableau de variations (tkz-tab)
\providecommand{\barre}{\ensuremath{\|}}
% environnement proof (amsthm non chargé) utilisé par les modèles
\ifdefined\proof\else\newenvironment{proof}[1][Proof]{\par\noindent\textit{#1.}\ }{\qed\par}\fi
\usepackage{lastpage}
\usepackage{hyperref}
\hypersetup{hidelinks}

% ============================================================
% Palette « Pop » OnBuch+ — mêmes hex que la classe P (plus_ui.dart)
% ============================================================
\definecolor{poporange}{HTML}{F59321}
\definecolor{popdark}{HTML}{E07A0C}
\definecolor{popcream}{HTML}{FFF6E8}
\definecolor{popink}{HTML}{1C1714}
\definecolor{poppurple}{HTML}{7A5AE0}
\definecolor{popgreen}{HTML}{2E9E6A}
\definecolor{poppink}{HTML}{E0457B}
\definecolor{popblue}{HTML}{2F7DE1}
\definecolor{popgold}{HTML}{C98A12}
\definecolor{popmuted}{HTML}{8A7F76}
\definecolor{popline}{HTML}{EADFD2}
\definecolor{popgray}{HTML}{8A7F76}
\colorlet{popgrey}{popgray}
% Alias tolérants (le modèle nomme parfois les couleurs différemment)
\colorlet{popwhite}{white}
\colorlet{popblack}{popink}
\colorlet{popred}{poppink}
\colorlet{popyellow}{popgold}
% Teintes claires pour fonds de tableaux / zones
\colorlet{popblueL}{popblue!10!white}
\colorlet{poporangeL}{poporange!14!white}
\colorlet{popgreenL}{popgreen!12!white}
\colorlet{poppurpleL}{poppurple!12!white}
\colorlet{poppinkL}{poppink!12!white}
\colorlet{popgoldL}{popgold!14!white}

\pagecolor{popcream}

% ============================================================
% Typographie
% ============================================================
\defaultfontfeatures{Ligatures=TeX}
\setmainfont{PlusJakartaSans-Medium.ttf}[Path=../../../fonts/,BoldFont=PlusJakartaSans-Bold.ttf]
% Maths en sans-serif (Fira Math) avec chiffres et lettres droites de Plus
% Jakarta, pour que les formules se fondent dans le texte.
\usepackage[math-style=ISO,bold-style=ISO]{unicode-math}
\setmathfont{FiraMath-Regular.otf}[Path=../../../fonts/]
\setmathfont{PlusJakartaSans-Medium.ttf}[Path=../../../fonts/,range={up/num,up/{Latin,latin},\mathup}]
% (icomma retiré : décimales avec point en anglais)
% Caractères mathématiques Unicode absents des polices de texte (Plus Jakarta,
% Archivo Black) : rendus par leur équivalent mathématique, en texte comme en
% formule — sinon ils s'affichent en carré vide (ex. « Arithmétique dans ℤ »).
\usepackage{newunicodechar}
\newunicodechar{ℝ}{\ensuremath{\mathbb{R}}}
\newunicodechar{ℤ}{\ensuremath{\mathbb{Z}}}
\newunicodechar{ℂ}{\ensuremath{\mathbb{C}}}
\newunicodechar{ℕ}{\ensuremath{\mathbb{N}}}
\newunicodechar{ℚ}{\ensuremath{\mathbb{Q}}}
\newunicodechar{𝔻}{\ensuremath{\mathbb{D}}}
\newunicodechar{α}{\ensuremath{\alpha}}
\newunicodechar{β}{\ensuremath{\beta}}
\newunicodechar{γ}{\ensuremath{\gamma}}
\newunicodechar{δ}{\ensuremath{\delta}}
\newunicodechar{ε}{\ensuremath{\varepsilon}}
\newunicodechar{θ}{\ensuremath{\theta}}
\newunicodechar{λ}{\ensuremath{\lambda}}
\newunicodechar{μ}{\ensuremath{\mu}}
\newunicodechar{σ}{\ensuremath{\sigma}}
\newunicodechar{τ}{\ensuremath{\tau}}
\newunicodechar{φ}{\ensuremath{\varphi}}
\newunicodechar{ψ}{\ensuremath{\psi}}
\newunicodechar{ω}{\ensuremath{\omega}}
\newunicodechar{Σ}{\ensuremath{\Sigma}}
% majuscules produites par \MakeUppercase dans les titres (α-aminés -> Α)
\newunicodechar{Α}{\ensuremath{\alpha}}
\newunicodechar{Β}{\ensuremath{\beta}}
\newunicodechar{Γ}{\ensuremath{\Gamma}}
\newunicodechar{Δ}{\ensuremath{\Delta}}
\newunicodechar{Ε}{\ensuremath{\varepsilon}}
\newunicodechar{Θ}{\ensuremath{\Theta}}
\newunicodechar{Λ}{\ensuremath{\Lambda}}
\newunicodechar{Μ}{\ensuremath{\mu}}
\newunicodechar{Τ}{\ensuremath{\tau}}
\newunicodechar{Φ}{\ensuremath{\Phi}}
\newunicodechar{Ψ}{\ensuremath{\Psi}}
\newunicodechar{Ω}{\ensuremath{\Omega}}
\newunicodechar{↦}{\ensuremath{\mapsto}}
\newunicodechar{∈}{\ensuremath{\in}}
\newunicodechar{∉}{\ensuremath{\notin}}
\newunicodechar{∧}{\ensuremath{\wedge}}
\newunicodechar{⇔}{\ensuremath{\Leftrightarrow}}
\newunicodechar{⟺}{\ensuremath{\Longleftrightarrow}}
\newunicodechar{⇒}{\ensuremath{\Rightarrow}}
\newunicodechar{⟹}{\ensuremath{\Longrightarrow}}
\newunicodechar{∘}{\ensuremath{\circ}}
\newunicodechar{∪}{\ensuremath{\cup}}
\newunicodechar{∩}{\ensuremath{\cap}}
\newunicodechar{≡}{\ensuremath{\equiv}}
\newunicodechar{⊂}{\ensuremath{\subset}}
\newunicodechar{⊥}{\ensuremath{\perp}}
\newunicodechar{∃}{\ensuremath{\exists}}
\newunicodechar{∀}{\ensuremath{\forall}}
\newunicodechar{∛}{\ensuremath{\sqrt[3]{\,}}}
\newunicodechar{∜}{\ensuremath{\sqrt[4]{\,}}}
\newunicodechar{⃗}{\ensuremath{{}^{\to}}}
\newunicodechar{ⁿ}{\ifmmode^{n}\else\textsuperscript{\ensuremath{n}}\fi}
\newunicodechar{ˣ}{\ifmmode^{x}\else\textsuperscript{\ensuremath{x}}\fi}
\newunicodechar{ᵇ}{\ifmmode^{b}\else\textsuperscript{\ensuremath{b}}\fi}
\newunicodechar{ʳ}{\ifmmode^{r}\else\textsuperscript{\ensuremath{r}}\fi}
\newunicodechar{ᵘ}{\ifmmode^{u}\else\textsuperscript{\ensuremath{u}}\fi}
\newunicodechar{ʸ}{\ifmmode^{y}\else\textsuperscript{\ensuremath{y}}\fi}
\newunicodechar{ᵃ}{\ifmmode^{a}\else\textsuperscript{\ensuremath{a}}\fi}
\newunicodechar{ᵉ}{\ifmmode^{e}\else\textsuperscript{\ensuremath{e}}\fi}
\newunicodechar{ᵏ}{\ifmmode^{k}\else\textsuperscript{\ensuremath{k}}\fi}
\newunicodechar{ᵖ}{\ifmmode^{p}\else\textsuperscript{\ensuremath{p}}\fi}
\newunicodechar{ᴺ}{\ifmmode^{N}\else\textsuperscript{\ensuremath{N}}\fi}
\newunicodechar{ᶜ}{\ifmmode^{c}\else\textsuperscript{\ensuremath{c}}\fi}
\newunicodechar{ʰ}{\ifmmode^{h}\else\textsuperscript{\ensuremath{h}}\fi}
\newunicodechar{ᵠ}{\ifmmode^{q}\else\textsuperscript{\ensuremath{q}}\fi}
\newunicodechar{ₙ}{\ifmmode_{n}\else\textsubscript{\ensuremath{n}}\fi}
\newunicodechar{ₐ}{\ifmmode_{a}\else\textsubscript{\ensuremath{a}}\fi}
\newunicodechar{ᵢ}{\ifmmode_{i}\else\textsubscript{\ensuremath{i}}\fi}
\newunicodechar{ᵣ}{\ifmmode_{r}\else\textsubscript{\ensuremath{r}}\fi}
\newunicodechar{ₖ}{\ifmmode_{k}\else\textsubscript{\ensuremath{k}}\fi}
\newunicodechar{ᵦ}{\ifmmode_{\beta}\else\textsubscript{\ensuremath{\beta}}\fi}
\newunicodechar{ₓ}{\ifmmode_{x}\else\textsubscript{\ensuremath{x}}\fi}
\newfontfamily\popdisplay{ArchivoBlack-Regular.ttf}[Path=../../../fonts/]
\newfontfamily\poplogo{SpaceGrotesk-Bold.ttf}[Path=../../../fonts/]
\newfontfamily\popsemi{PlusJakartaSans-SemiBold.ttf}[Path=../../../fonts/]
\newfontfamily\popxbold{PlusJakartaSans-ExtraBold.ttf}[Path=../../../fonts/]
\newfontfamily\popxboldls{PlusJakartaSans-ExtraBold.ttf}[Path=../../../fonts/,LetterSpace=3.0]

\setlist[enumerate]{leftmargin=1.7em,itemsep=5pt,topsep=4pt}
\setlist[itemize]{leftmargin=1.7em,itemsep=4pt,topsep=4pt,label={\color{poporange}\textbullet}}
\setlength{\parindent}{0pt}
\setlength{\parskip}{5pt}
\renewcommand{\arraystretch}{1.25}

% pgfplots : style Pop par défaut
\pgfplotsset{
  every axis/.append style={axis line style={popink,line width=1pt},tick style={popink},
    grid style={popline,line width=0.5pt},label style={font=\small},tick label style={font=\footnotesize}},
  popcurve/.style={poporange,line width=1.8pt,no markers,smooth},
}
% Même style disponible dans un \draw TikZ simple (hors pgfplots)
\tikzset{popcurve/.style={poporange,line width=1.8pt}}
\tikzset{popgrid/.style={popline,very thin}}
\colorlet{popcurve}{poporange} % le nom du style est parfois utilisé comme couleur

% ============================================================
% Pill / squiggle
% ============================================================
\newcommand{\poppill}[3]{%
  \tikz[baseline=-0.55ex]\node[draw=popink,line width=1.2pt,rounded corners=2.6pt,%
    fill=#2,inner xsep=6.5pt,inner ysep=3pt]{%
    \popxboldls\fontsize{7}{7}\selectfont\color{#3}\MakeUppercase{#1}};%
}
\newcommand{\popsquiggle}[1]{%
  \tikz[baseline=-0.4ex]{
    \draw[#1,line width=2.6pt,line cap=round]
      plot [smooth] coordinates {(0,0.09) (0.34,0.01) (0.68,0.21) (1.02,0.01) (1.36,0.10)};
  }%
}
% Mot-clé surligné (marqueur orange sous le texte)
\newcommand{\cle}[1]{{\popxbold\color{popdark}#1}}

% ============================================================
% En-tête / pied
% ============================================================
\pagestyle{fancy}
\fancyhf{}
\renewcommand{\headrulewidth}{0pt}
\renewcommand{\footrulewidth}{0pt}
\lhead{\raisebox{-0.15ex}{\poplogo\fontsize{13}{13}\selectfont\color{popink}OnBuch\color{poporange}+}}
\rhead{\poppill{\DOCMATIERE\ \textbullet\ \DOCNIVEAU\ \textbullet\ Chapter \DOCLECON}{white}{popink}}
\lfoot{\popxbold\fontsize{7.6}{7.6}\selectfont\color{popmuted}\MakeUppercase{Signed course OnBuch+}\ \textbullet\ {\popsemi\DOCTITRE}}
\rfoot{\tikz[baseline=-0.6ex]\node[rounded corners=8.5pt,fill=popink,minimum height=17pt,minimum width=17pt,inner xsep=5pt,inner sep=0]{\popxbold\fontsize{8}{8}\selectfont\color{popcream}\thepage\,/\,\pageref*{LastPage}};}

% ============================================================
% Titres
% ============================================================
\newcommand{\chaptitre}[2]{%
  \begin{tcolorbox}[enhanced,colback=poporange,colframe=popink,boxrule=2.6pt,arc=11pt,
      drop shadow=popink,left=15pt,right=15pt,top=12pt,bottom=11pt,before skip=3pt,after skip=18pt]
    \poppill{#2}{popink}{popcream}\par\vspace{9pt}
    {\raggedright\hyphenpenalty=10000\exhyphenpenalty=10000\popdisplay\fontsize{22}{24}\selectfont\color{popcream}\MakeUppercase{#1}\par}
  \end{tcolorbox}
}
% Section numérotée : pastille orange avec le numéro + titre Archivo
\newcounter{popsec}
\newcommand{\coursec}[1]{%
  \stepcounter{popsec}\setcounter{popsub}{0}%
  \par\vspace{16pt}\noindent
  \begin{minipage}{\linewidth}
  \tikz[baseline=-0.7ex]\node[draw=popink,line width=1.6pt,fill=poporange,rounded corners=4pt,minimum size=22pt,inner sep=0]{\popdisplay\fontsize{12}{12}\selectfont\color{popcream}\thepopsec};%
  \hspace{8pt}{\popdisplay\fontsize{15}{17}\selectfont\color{popink}\MakeUppercase{#1}}\par
  \vspace{4pt}\noindent{\color{poporange}\rule{36pt}{3.6pt}}
  \end{minipage}\par\nopagebreak\vspace{8pt}
}
\newcounter{popsub}
\newcommand{\courssub}[1]{%
  \stepcounter{popsub}%
  \par\vspace{9pt}\noindent{\popxbold\fontsize{11.5}{13}\selectfont\color{popink}{\color{poporange}\thepopsec.\thepopsub}\hspace{6pt}#1}\par\nopagebreak\vspace{4pt}
}

% ============================================================
% Boîtes pédagogiques — même recette : fond blanc, bordure épaisse colorée,
% ombre dure encre, badge plein. Titre optionnel [#1].
% ============================================================
\tcbset{popbase/.style={breakable,enhanced,colback=white,boxrule=2.3pt,arc=9pt,
  shadow={3pt}{-3pt}{0pt}{popink},left=14pt,right=14pt,top=11pt,bottom=11pt,before skip=11pt,after skip=15pt,
  fontupper=\popsemi\fontsize{10.6}{14.5}\selectfont\color{popink},
  fontlower=\popsemi\fontsize{10.6}{14.5}\selectfont\color{popink}}}
% Coche dessinée (le glyphe ✓ n'existe pas dans les polices du document)
\newcommand{\popcheck}[1]{\tikz[baseline=-0.5ex]\draw[#1,line width=1.6pt,line cap=round,line join=round](-0.13,0.0)--(-0.04,-0.09)--(0.13,0.1);}
\providecommand{\coche}{\popcheck{popgreen}}
\providecommand{\resultat}[1]{\par\smallskip\noindent\fcolorbox{popgreen}{popgreenL}{\parbox{\dimexpr\linewidth-2\fboxsep-2\fboxrule\relax}{\textbf{Result.} #1}}\par\smallskip}
\newcommand{\popboxhead}[3]{\poppill{#1}{#2}{white}\ifx\relax#3\relax\else\hspace{6pt}{\popxbold\fontsize{10.6}{12}\selectfont\color{popink}#3}\fi\par\vspace{7pt}}

\newtcolorbox{aretenir}[1][]{popbase,colframe=popgreen,before upper={\popboxhead{Key points}{popgreen}{#1}}}
\newtcolorbox{exemplebox}[1][]{popbase,colframe=poporange,before upper={\popboxhead{Exemple}{poporange}{#1}}}
\newtcolorbox{definition}[1][]{popbase,colframe=popblue,before upper={\popboxhead{Definition}{popblue}{#1}}}
\newtcolorbox{propriete}[1][]{popbase,colframe=poppurple,before upper={\popboxhead{Property}{poppurple}{#1}}}
\newtcolorbox{methode}[1][]{popbase,colframe=popgold,before upper={\popboxhead{Method}{popgold}{#1}}}
\newtcolorbox{attention}[1][]{popbase,colframe=poppink,before upper={\popboxhead{Warning}{poppink}{#1}}}
\newtcolorbox{experience}[1][]{popbase,colframe=popblue,colback=popblueL,before upper={\popboxhead{Experiment}{popblue}{#1}}}
% « Le savais-tu ? » : carte violette pleine (contexte, culture, Cameroun)
\newtcolorbox{savaistu}[1][]{popbase,colback=poppurple,colframe=popink,
  fontupper=\popsemi\fontsize{10.4}{14.2}\selectfont\color{white},
  before upper={\poppill{Did you know?}{white}{poppurple}\ifx\relax#1\relax\else\hspace{6pt}{\popxbold\color{white}#1}\fi\par\vspace{7pt}}}
% Exercice résolu : énoncé en haut, \tcblower puis la solution sur fond crème
\newcounter{popexo}\newcounter{popexr}
\newtcolorbox{exoresolu}[1][]{popbase,colframe=poporange,colbacklower=popcream,
  segmentation style={popink,line width=1.2pt,dashed},
  before upper={\stepcounter{popexr}\popboxhead{Worked example \thepopexr}{poporange}{#1}},
  before lower={\poppill{Solution}{popink}{popcream}\par\vspace{6pt}}}
% Exercice (non résolu en place) — la correction est dans la partie Corrigés
\newtcolorbox{exercice}[1][]{popbase,colframe=popink,
  before upper={\stepcounter{popexo}\popboxhead{Exercise \thepopexo}{popink}{#1}}}
% Corrigé
\newtcolorbox{corrige}[1][]{popbase,colframe=popgreen,colback=popgreenL,
  before upper={\popboxhead{Solution}{popgreen}{#1}}}
% Objectifs / prérequis / bilan
\newtcolorbox{objectifs}{popbase,colframe=popink,colback=poporangeL,
  before upper={\popboxhead{Chapter objectives}{popink}{}}}
\newtcolorbox{prerequis}{popbase,colframe=popink,colback=popblueL,
  before upper={\popboxhead{Prerequisites}{popblue}{}}}
\newtcolorbox{bilan}{popbase,colframe=popink,colback=white,boxrule=2.8pt,
  before upper={\popboxhead{Fiche bilan}{poporange}{L'essentiel à connaître pour l'examen}}}
% Figure encadrée avec légende
\newtcolorbox{popfigure}[1][]{enhanced,breakable=false,colback=white,colframe=popline,boxrule=1.4pt,arc=8pt,
  left=10pt,right=10pt,top=10pt,bottom=8pt,before skip=10pt,after skip=12pt,halign=center,
  lower separated=false,halign lower=center,
  fontlower=\popsemi\fontsize{9}{11.5}\selectfont\color{popmuted},
  #1}
\newcommand{\legende}[1]{\par\vspace{6pt}{\popsemi\fontsize{9}{11.5}\selectfont\color{popmuted}#1}}

% ============================================================
% Signature OnBuch+
% ============================================================
\newcommand{\poplogoinline}[1]{{\poplogo\fontsize{#1}{#1}\selectfont\color{popink}OnBuch\color{poporange}+}}
\newcommand{\signatureonbuch}{%
  \begin{tcolorbox}[enhanced,colback=popink,colframe=popink,boxrule=2.2pt,arc=10pt,
      drop shadow=poporange,left=14pt,right=14pt,top=10pt,bottom=10pt,before skip=0pt,after skip=0pt]
    \tikz[baseline=-0.7ex]\node[circle,fill=popgreen,draw=popcream,line width=1.3pt,minimum size=22pt,inner sep=0]{\popcheck{white}};%
    \hspace{9pt}\begin{minipage}[c]{0.62\linewidth}
      {\popxbold\fontsize{12}{13}\selectfont\color{popcream}Signed\ \textbullet\ {\poplogo OnBuch\color{poporange}+}}\par\vspace{2pt}
      {\raggedright\popsemi\fontsize{8}{10.5}\selectfont\color{popline}\MakeUppercase{OnBuch+ teaching team}\\\MakeUppercase{Follows the official MINESEC syllabus}\par}
    \end{minipage}\hfill
    {\poplogo\fontsize{22}{22}\selectfont\color{popcream}OnBuch\color{poporange}+}
  \end{tcolorbox}%
}

% ============================================================
% Page de garde
% #1 titre · #2 matière · #3 niveau · #4 module · #5 n° leçon · #6 durée ·
% #7 description / objectifs (texte ou itemize)
% ============================================================
\newcommand{\pagegarde}[7]{%
  \thispagestyle{empty}
  \newgeometry{a4paper,margin=1.7cm,top=1.6cm,bottom=1.4cm}%
  \begin{tikzpicture}[remember picture,overlay]
    \fill[poporange!18] ([xshift=-3.2cm,yshift=-3.6cm]current page.north east) circle (4.6cm);
    \fill[poppurple!14] ([xshift=2.4cm,yshift=7.6cm]current page.south west) circle (3.8cm);
  \end{tikzpicture}%
  {\poplogo\fontsize{30}{30}\selectfont\color{popink}OnBuch\color{poporange}+}\hfill
  \raisebox{0.6ex}{\poppill{Signed course \textbullet\ Premium}{popink}{popcream}}\par
  \vspace{1pt}\popsquiggle{poppurple}\par
  \vspace{8pt}
  \begin{tcolorbox}[enhanced,colback=poporange,colframe=popink,boxrule=2.8pt,arc=13pt,
      drop shadow=popink,left=18pt,right=18pt,top=12pt,bottom=13pt,after skip=10pt]
    \poppill{#2 \textbullet\ #3}{popink}{popcream}\hspace{5pt}\poppill{#4}{white}{popink}\par\vspace{8pt}
    {\popdisplay\fontsize{14}{14}\selectfont\color{popink}CHAPTER #5}\par\vspace{4pt}
    {\raggedright\hyphenpenalty=10000\exhyphenpenalty=10000\popdisplay\fontsize{25}{27}\selectfont\color{popcream}\MakeUppercase{#1}\par}
    \vspace{8pt}
    {\popxbold\fontsize{9.5}{9.5}\selectfont\color{popink}\MakeUppercase{Suggested time: #6}}
  \end{tcolorbox}
  \begin{tcolorbox}[enhanced,colback=white,colframe=popink,boxrule=2.2pt,arc=10pt,
      drop shadow=popink,left=16pt,right=16pt,top=10pt,bottom=10pt,after skip=8pt]
    \poppill{In this chapter}{poppurple}{white}\par\vspace{5pt}
    \popsemi\fontsize{9.3}{12.6}\selectfont\color{popink}#7
  \end{tcolorbox}
  \noindent\begin{minipage}[t]{0.487\linewidth}
    \begin{tcolorbox}[enhanced,colback=popgreen,colframe=popink,boxrule=2.2pt,arc=10pt,
        drop shadow=popink,left=11pt,right=11pt,top=10pt,bottom=10pt,height=3.1cm,valign=center]
      \tcbox[colback=white,colframe=popink,boxrule=1.3pt,arc=5pt,boxsep=0pt,nobeforeafter,
        left=3pt,right=3pt,top=3pt,bottom=3pt,valign=center]{\qrcode[height=1.9cm]{https://play.google.com/store/apps/details?id=com.luvvix.onbuch&hl=en}}%
      \hspace{8pt}\begin{minipage}[c]{0.5\linewidth}
        {\popdisplay\fontsize{11}{12.5}\selectfont\color{white}\MakeUppercase{Download OnBuch}}\par\vspace{4pt}
        {\raggedright\popsemi\fontsize{8.4}{11}\selectfont\color{white}Free \textbullet\ Google Play\\AI tutor, past papers, results\par}
      \end{minipage}
    \end{tcolorbox}
  \end{minipage}\hfill
  \begin{minipage}[t]{0.487\linewidth}
    \begin{tcolorbox}[enhanced,colback=poppurple,colframe=popink,boxrule=2.2pt,arc=10pt,
        drop shadow=popink,left=11pt,right=11pt,top=10pt,bottom=10pt,height=3.1cm,valign=center]
      \tikz[baseline=-0.7ex]\node[circle,fill=white,draw=popink,line width=1.3pt,minimum size=1.6cm,inner sep=0]{\popdisplay\fontsize{24}{24}\selectfont\color{poppurple}+};%
      \hspace{8pt}\begin{minipage}[c]{0.6\linewidth}
        {\popdisplay\fontsize{11}{12.5}\selectfont\color{white}\MakeUppercase{Go OnBuch+}}\par\vspace{4pt}
        {\raggedright\popsemi\fontsize{8.4}{11}\selectfont\color{white}All signed courses, unlimited AI tutor, PDF sheets — OnBuch+ tab in the app\par}
      \end{minipage}
    \end{tcolorbox}
  \end{minipage}
  \vspace{8pt}
  \popcontenu
  \vfill
  \signatureonbuch
  \restoregeometry
  \newpage
}

% Rangée « Ce cours contient » (page de garde)
\newcommand{\poptile}[3]{% #1 couleur #2 gros texte #3 légende
  \begin{tcolorbox}[enhanced,colback=#1,colframe=popink,boxrule=2pt,arc=9pt,shadow={2.5pt}{-2.5pt}{0pt}{popink},
      left=6pt,right=6pt,top=9pt,bottom=9pt,halign=center,valign=center,height=1.75cm,width=\linewidth,nobeforeafter]
    {\popdisplay\fontsize{12.5}{13}\selectfont\color{white}\MakeUppercase{#2}}\par\vspace{3pt}
    {\centering\popsemi\fontsize{7.8}{9.5}\selectfont\color{white}#3\par}
  \end{tcolorbox}}
\newcommand{\popcontenu}{%
  \par\noindent{\popxboldls\fontsize{7.5}{7.5}\selectfont\color{popmuted}THIS SIGNED COURSE CONTAINS}\par\vspace{7pt}
  \noindent\begin{minipage}[t]{0.235\linewidth}\poptile{popblue}{Course}{complete, illustrated, step by step}\end{minipage}\hfill
  \begin{minipage}[t]{0.235\linewidth}\poptile{poporange}{Methods}{and worked examples}\end{minipage}\hfill
  \begin{minipage}[t]{0.235\linewidth}\poptile{poppink}{Exercises}{exam style + solutions}\end{minipage}\hfill
  \begin{minipage}[t]{0.235\linewidth}\poptile{popgreen}{Summary sheet}{the essentials to remember}\end{minipage}\par}

% Sommaire
\newtcolorbox{sommaire}{enhanced,colback=white,colframe=popink,boxrule=2.2pt,arc=10pt,
  drop shadow=popink,left=18pt,right=18pt,top=13pt,bottom=13pt,before skip=6pt,after skip=20pt,
  before upper={\poppill{Contents}{poppurple}{white}\par\vspace{9pt}\popsemi\fontsize{10.6}{14.5}\selectfont\color{popink}}}

% Signature de fin de cours — poussée tout en bas de la dernière page
\newcommand{\findecours}{%
  \par\vfill\nopagebreak
  \begin{center}\popsquiggle{poporange}\end{center}\vspace{4pt}
  \signatureonbuch
}
\AtBeginDocument{\def\llbracket{\mathopen{[\mkern-3.5mu[}}\def\rrbracket{\mathclose{]\mkern-3.5mu]}}} % intervalles d'entiers (glyphe ⟦ absent de la police)
% Glyphes absents de Fira Math : redéfinis à partir de glyphes disponibles
\AtBeginDocument{%
  \def\setminus{\mathbin{\backslash}}\let\smallsetminus\setminus
  \def\vdots{\mathord{\vbox{\baselineskip3.2pt\lineskiplimit0pt\kern4pt\hbox{.}\hbox{.}\hbox{.}}}}%
  \def\ddots{\mathinner{\mkern1mu\raise6.5pt\hbox{.}\mkern2mu\raise3.6pt\hbox{.}\mkern2mu\raise0.7pt\hbox{.}\mkern1mu}}%
  \def\triangle{\mathord{\scalebox{0.9}{$\symup{\Delta}$}}}%
  \def\nabla{\mathord{\raisebox{\depth}{\scalebox{1}[-1]{$\symup{\Delta}$}}}}%
  \def\checkmark{\popcheck{popdark}}%
  \def\star{\mathbin{\ast}}\def\bigstar{\mathord{\ast}}%
  \def\diamond{\mathbin{\scalebox{0.6}{$\Diamond$}}}%
  \def\bigcup{\mathop{\mathchoice{\raisebox{-0.3em}{\scalebox{1.7}{$\cup$}}}{\scalebox{1.25}{$\cup$}}{\cup}{\cup}}\displaylimits}%
  \def\bigcap{\mathop{\mathchoice{\raisebox{-0.3em}{\scalebox{1.7}{$\cap$}}}{\scalebox{1.25}{$\cap$}}{\cap}{\cap}}\displaylimits}%
  \def\lesseqqgtr{\mathrel{\lessgtr}}\def\bigoplus{\mathop{\oplus}\displaylimits}%
}
\providecommand{\ssi}{if and only if} % abréviation fréquente
\AtBeginDocument{\providecommand{\wideparen}{\overparen}} % arc de cercle

%% FIGFIX-BEGIN v3 — correctifs globaux des figures (docs/onbuch-generation/tools/figfix)
\usepackage{adjustbox}\usepackage{etoolbox}
% toute figure TikZ/pgfplots plus large que la ligne est réduite à la largeur disponible
\BeforeBeginEnvironment{tikzpicture}{\begin{adjustbox}{max width=\linewidth}}
\AfterEndEnvironment{tikzpicture}{\end{adjustbox}}
% légendes pgfplots lisibles par-dessus les courbes
\pgfplotsset{every axis/.append style={legend style={fill=white,fill opacity=0.92,text opacity=1,draw=black!15,inner sep=3pt}}}
% étiquettes d'arêtes (syntaxe edge["..."]) avec fond blanc
\tikzset{every edge quotes/.append style={fill=white,inner sep=1.2pt}}
% bulles de cartes mentales : texte à la ligne dans la bulle, bulle un peu plus grande
\tikzset{every concept/.append style={text width=2.0cm,align=center,minimum size=2.3cm,inner sep=2pt}}
%% FIGFIX-END
\begin{document}

\pagegarde{Reaction kinetics}{Chemistry}{Upper Sixth}{Upper Sixth — Chemistry}{4}{12 periods}{This chapter introduces the concepts, experimental techniques, mathematical laws, and theoretical models that allow chemists to quantify and predict the rates of chemical reactions. It equips students with the tools to analyse kinetic data, determine reaction mechanisms, and evaluate the role of catalysts.
\vspace{4pt}
\begin{multicols}{2}\small
\begin{itemize}
\item Define reaction rate, average rate, instantaneous rate, and initial rate, and state their units.
\item Identify and explain the factors that affect reaction rates (concentration, temperature, surface area, catalyst).
\item Design and interpret experiments using chemical (titrimetry) and physical (colorimetry, conductimetry, dilatometry, gas syringe, pressure) methods to follow reaction progress.
\item Write the rate law for a reaction, define rate constant, partial orders, and overall order, and determine their units.
\item Determine reaction order and rate constant using the initial‑rates method and graphical methods (including half‑life for first‑order reactions).
\item Apply collision theory, Maxwell‑Boltzmann distribution, and the Arrhenius equation to explain temperature dependence of rates.
\end{itemize}\end{multicols}}

\begin{sommaire}
\begin{multicols}{2}
\begin{enumerate}
\item 1. Foundations of Reaction Kinetics
\item 2. Chemical Method – Titrimetric Rate Determination
\item 3. Physical Methods of Rate Measurement
\item 4. Rate Law, Order of Reaction, and Determination Methods
\item 5. Theoretical Models: Collision Theory, Arrhenius Equation, and Transition‑State Theory
\item 6. Catalysis: Homogeneous and Heterogeneous
\item Integration activity
\item Methods and worked examples
\item Exercises
\item Solutions to the exercises
\item Summary sheet
\end{enumerate}
\end{multicols}
\end{sommaire}

\begin{objectifs}
\begin{itemize}
\item Define reaction rate, average rate, instantaneous rate, and initial rate, and state their units.
\item Identify and explain the factors that affect reaction rates (concentration, temperature, surface area, catalyst).
\item Design and interpret experiments using chemical (titrimetry) and physical (colorimetry, conductimetry, dilatometry, gas syringe, pressure) methods to follow reaction progress.
\item Write the rate law for a reaction, define rate constant, partial orders, and overall order, and determine their units.
\item Determine reaction order and rate constant using the initial‑rates method and graphical methods (including half‑life for first‑order reactions).
\item Apply collision theory, Maxwell‑Boltzmann distribution, and the Arrhenius equation to explain temperature dependence of rates.
\item Describe transition‑state theory, the concept of the activated complex, and compare it with collision theory.
\item Distinguish homogeneous and heterogeneous catalysis, explain their mechanisms, and evaluate their industrial relevance.
\end{itemize}
\end{objectifs}

\begin{prerequis}
\begin{itemize}
\item Mole concept, stoichiometry, and concentration calculations (Lower Sixth).
\item Basic algebra, logarithms, and graph plotting (mathematics).
\item Understanding of energy changes in reactions (enthalpy, activation energy) from energetics chapter.
\item Familiarity with common laboratory apparatus (burette, spectrophotometer, conductivity meter, gas syringe).
\end{itemize}
\end{prerequis}

\newpage

\coursec{Foundations of Reaction Kinetics}

Imagine you are a farmer in the Northwest Region watching a bag of \ce{NPK} fertilizer dissolve in a water tank for irrigation. The speed at which the nutrients become available determines the health of your crops. Understanding why some reactions are fast and others slow, and how to measure and control them, is the power of reaction kinetics.

\courssub{What Is Reaction Rate?}

Chemical reactions do not happen instantaneously; they proceed at a certain speed. The \cle{rate of reaction} tells us how quickly reactants are consumed or products are formed per unit time. 

Consider a general reaction:
\[
a\ce{A} + b\ce{B} \rightarrow c\ce{C} + d\ce{D}
\]
where $a$, $b$, $c$, $d$ are the stoichiometric coefficients. The rate of reaction $v$ (or $r$) is defined so that it has the same value whether we follow the disappearance of a reactant or the appearance of a product.

\begin{definition}[Rate of Reaction]
For the reaction $a\ce{A} + b\ce{B} \rightarrow c\ce{C} + d\ce{D}$, the rate of reaction $v$ is defined as:
\[
v = -\frac{1}{a}\frac{d[\ce{A}]}{dt} = -\frac{1}{b}\frac{d[\ce{B}]}{dt} = +\frac{1}{c}\frac{d[\ce{C}]}{dt} = +\frac{1}{d}\frac{d[\ce{D}]}{dt}
\]
where $[\ce{X}]$ denotes the concentration of species $\ce{X}$ in $\mathrm{mol\,dm^{-3}}$, and $t$ is time in seconds. The negative sign for reactants ensures $v$ is a positive quantity.
\end{definition}

\begin{propriete}[Rate of Change of Concentration vs. Reaction Rate]
\begin{itemize}
\item For a \textbf{reactant}: rate of disappearance (rate of change of concentration) $= -\frac{d[\text{reactant}]}{dt} = a \times v$.
\item For a \textbf{product}: rate of appearance (rate of change of concentration) $= +\frac{d[\text{product}]}{dt} = c \times v$.
\end{itemize}
The \cle{units of reaction rate} $v$ are always $\mathrm{mol\,dm^{-3}\,s^{-1}}$ (or $\mathrm{M\,s^{-1}}$). The units of $\frac{d[\ce{X}]}{dt}$ are also $\mathrm{mol\,dm^{-3}\,s^{-1}}$.
\end{propriete}

\courssub{Average, Instantaneous, and Initial Rates}

\begin{exemplebox}[Decomposition of Dinitrogen Pentoxide]
The gas-phase decomposition $\ce{2N2O5(g) -> 4NO2(g) + O2(g)}$ is studied. At $t = 0$, $[\ce{N2O5}] = 0.100\ \mathrm{mol\,dm^{-3}}$. After $100\ \mathrm{s}$, $[\ce{N2O5}] = 0.070\ \mathrm{mol\,dm^{-3}}$.
\begin{enumerate}
\item Calculate the average rate of reaction over the first $100\ \mathrm{s}$.
\item Express the rate in terms of $\ce{NO2}$ formation.
\end{enumerate}
\tcblower
\textbf{Solution:}
\begin{enumerate}
\item Change in $[\ce{N2O5}] = 0.070 - 0.100 = -0.030\ \mathrm{mol\,dm^{-3}}$.
Average rate of disappearance of $\ce{N2O5} = \frac{0.030}{100} = 3.0 \times 10^{-4}\ \mathrm{mol\,dm^{-3}\,s^{-1}}$.
Reaction rate $v = -\frac{1}{2}\frac{\Delta[\ce{N2O5}]}{\Delta t} = \frac{1}{2} \times 3.0 \times 10^{-4} = \mathbf{1.5 \times 10^{-4}\ \mathrm{mol\,dm^{-3}\,s^{-1}}}$.
\item $v = +\frac{1}{4}\frac{\Delta[\ce{NO2}]}{\Delta t} \Rightarrow \frac{\Delta[\ce{NO2}]}{\Delta t} = 4v = 6.0 \times 10^{-4}\ \mathrm{mol\,dm^{-3}\,s^{-1}}$.
\end{enumerate}
\end{exemplebox}

The \cle{average rate} over a time interval $\Delta t$ is $\frac{\Delta[\text{species}]}{\Delta t}$ (with stoichiometric factor). However, the rate usually changes as concentrations change. We need the \cle{instantaneous rate} at a specific time $t$, which is the derivative $\frac{d[\text{species}]}{dt}$ at that instant.

\begin{methode}[Finding Instantaneous Rate from a Concentration–Time Graph]
\begin{enumerate}
\item Plot concentration of a reactant or product against time.
\item At the desired time $t$, draw a \textbf{tangent} to the curve.
\item Choose two points on the tangent line, $(t_1, c_1)$ and $(t_2, c_2)$, far apart for accuracy.
\item Calculate the gradient of the tangent: $\text{gradient} = \frac{c_2 - c_1}{t_2 - t_1} = \frac{d[\text{species}]}{dt}$.
\item Apply the stoichiometric factor to obtain the reaction rate $v$:
\begin{itemize}
\item If the curve is for a \textbf{reactant} $\ce{A}$ (coefficient $a$): gradient is negative. $v = -\frac{1}{a} \times \text{gradient}$.
\item If the curve is for a \textbf{product} $\ce{C}$ (coefficient $c$): gradient is positive. $v = +\frac{1}{c} \times \text{gradient}$.
\end{itemize}
\end{enumerate}
\end{methode}

The \cle{initial rate} is the instantaneous rate at $t = 0$. It is especially important because it depends only on the initial concentrations, before any significant product has formed or reactant has been depleted.

\begin{popfigure}
\begin{tikzpicture}[scale=0.9]
\begin{axis}[
    width=12cm, height=8cm,
    xlabel={Time $t$ / s}, ylabel={$[\ce{A}]$ / mol dm$^{-3}$},
    xmin=0, xmax=100, ymin=0, ymax=1.1,
    xtick={0,20,40,60,80,100}, ytick={0,0.2,0.4,0.6,0.8,1.0},
    grid=major, grid style={gray!30},
    axis lines=middle, enlargelimits=false,
    clip=false
]
% Concentration curve: exponential decay [A] = 1.0 * exp(-0.02*t)
\addplot[domain=0:100, samples=200, thick, popblue] {1.0*exp(-0.02*x)};
\addlegendentry{$[\ce{A}] = [\ce{A}]_0 e^{-kt}$}

% Tangent at t=0 (initial rate)
% derivative at 0: -k*[A]_0 = -0.02
\addplot[domain=0:30, samples=2, dashed, poporange, thick] {1.0 - 0.02*x};
\node[poporange, anchor=south west, text width=3cm, align=left] at (axis cs:5,0.95) {Tangent at $t=0$ (initial rate)};

% Tangent at t=40 s
% [A] at 40 = 1*exp(-0.8) = 0.4493
% derivative = -0.02*0.4493 = -0.008986
% tangent line: y - 0.4493 = -0.008986(x - 40) => y = 0.80876 - 0.008986x
\addplot[domain=20:80, samples=2, dashed, popgreen, thick] {0.80876 - 0.008986*x};
\node[popgreen, anchor=north east, text width=3cm, align=right] at (axis cs:70,0.25) {Tangent at $t=40$ s (instantaneous rate)};

% Points on tangents for gradient calculation
\fill[poporange] (axis cs:0,1.0) circle (2pt) node[above left] {$(0,[\ce{A}]_0)$};
\fill[poporange] (axis cs:30,0.4) circle (2pt) node[below right] {$(30,0.4)$};
\fill[popgreen] (axis cs:20,0.629) circle (2pt) node[above left] {$(20,0.63)$};
\fill[popgreen] (axis cs:80,0.090) circle (2pt) node[below right] {$(80,0.09)$};

% Vertical line at t=40
\draw[gray, dashed] (axis cs:40,0) -- (axis cs:40,0.4493);
\node[below] at (axis cs:40,0) {$t=40\ \mathrm{s}$};
\end{axis}
\end{tikzpicture}
\legende{Concentration–time curve for a reactant $\ce{A}$ showing the tangent at $t=0$ (initial rate) and at $t=40\ \mathrm{s}$ (instantaneous rate). The gradient of each tangent gives the rate of change of concentration at that instant.}
\end{popfigure}

\begin{attention}[Average vs. Instantaneous Rate]
\textbf{Common mistake:} Using $\frac{\Delta[\ce{A}]}{\Delta t}$ over a large interval (e.g., $0$ to $100\ \mathrm{s}$) and calling it the rate at $t=50\ \mathrm{s}$. This gives the \emph{average} rate over that interval, not the instantaneous rate at the midpoint. The instantaneous rate is the gradient of the \emph{tangent} at the point. For a curved graph, the average rate over an interval is \textbf{not} equal to the instantaneous rate at the midpoint (except for a straight line, i.e., zero-order reactions).
\end{attention}

\courssub{Factors Affecting the Rate of Reaction}

Why does fertilizer dissolve quickly in hot water but slowly in cold? Why does powdered limestone react faster with acid than a large chip? The main factors are:

\begin{enumerate}
\item \textbf{Concentration (or pressure for gases):} Higher concentration $\rightarrow$ more particles per unit volume $\rightarrow$ more frequent collisions $\rightarrow$ higher rate. For gases, increasing pressure increases concentration ($c = n/V = p/RT$).
\item \textbf{Temperature:} Increasing temperature increases the \emph{kinetic energy} of molecules. More molecules have energy $\ge E_a$ (activation energy), and collisions are more frequent. A rough rule: rate $\approx$ doubles for every $10\ \mathrm{K}$ rise.
\item \textbf{Surface Area (heterogeneous reactions):} For a solid reacting with a liquid or gas, only the surface is available. Powdering increases surface area $\rightarrow$ more collision sites $\rightarrow$ faster rate.
\item \textbf{Catalyst:} A substance that increases the rate without being consumed. It provides an alternative pathway with \textbf{lower activation energy} $E_a$. It does \emph{not} change $\Delta H$ or the equilibrium constant.
\item \textbf{Nature of Reactants:} Ionic reactions (e.g., $\ce{Ag+ + Cl- -> AgCl v}$) are almost instantaneous. Covalent bond breaking (e.g., $\ce{H2 + Cl2 -> 2HCl}$) requires energy input and is slower.
\end{enumerate}

\begin{savaistu}[Cameroon Context: Limestone and Cement]
In Figuil (North Region), limestone ($\ce{CaCO3}$) is heated to make cement. The decomposition $\ce{CaCO3(s) -> CaO(s) + CO2(g)}$ is slow at room temperature but rapid at $900\text{--}1000\ \mathrm{^\circ C}$. The high temperature supplies the activation energy to break strong $\ce{Ca-O}$ and $\ce{C-O}$ bonds. In the lab, we study the reverse: $\ce{CaCO3 + 2HCl -> CaCl2 + H2O + CO2}$. Using powdered limestone instead of chips dramatically increases the rate — a classic demonstration of surface area effect.
\end{savaistu}

\courssub{Why Study Reaction Kinetics?}

\begin{itemize}
\item \textbf{Industrial Optimisation:} Haber process ($\ce{N2 + 3H2 <=> 2NH3}$) — kinetics tells us the best $T$, $p$, and catalyst (Fe) to maximise yield per unit time. In Cameroon, \emph{SONARA} (National Oil Refinery) uses kinetic data to optimise cracking and reforming units.
\item \textbf{Environmental Monitoring:} Rate of $\ce{SO2}$ oxidation to $\ce{SO3}$ (acid rain), $\ce{NO_x}$ formation in engines, $\ce{CFC}$ decomposition in stratosphere. Predicting pollutant lifetimes requires rate constants.
\item \textbf{Biological Processes:} Enzyme kinetics (Michaelis–Menten) governs digestion, respiration, drug metabolism. Understanding rates helps design medicines (e.g., antiretrovirals for HIV).
\item \textbf{Food Preservation:} Rates of spoilage reactions (oxidation, microbial growth) determine shelf life. Refrigeration lowers $T$ $\rightarrow$ slows rates.
\item \textbf{Safety:} Explosions are extremely fast reactions. Kinetics helps design controlled detonations (mining in East Region) and prevent accidental runaway reactions.
\end{itemize}

\courssub{Worked Example: Rates from a Graph}

\begin{exoresolu}[Determining Rates from a Concentration–Time Graph]
For the reaction $\ce{2A -> B + C}$, the concentration of $\ce{A}$ was measured over time:
\begin{center}
\begin{tabular}{|c|cccccc|}\hline
$t / \mathrm{s}$ & 0 & 20 & 40 & 60 & 80 & 100 \\ \hline
$[\ce{A}] / \mathrm{mol\,dm^{-3}}$ & 1.00 & 0.67 & 0.45 & 0.30 & 0.20 & 0.13 \\ \hline
\end{tabular}
\end{center}
\begin{enumerate}
\item Plot $[\ce{A}]$ vs $t$ and draw a smooth curve.
\item Determine the initial rate of reaction.
\item Determine the instantaneous rate at $t = 40\ \mathrm{s}$.
\item Calculate the average rate of formation of $\ce{B}$ between $t = 20\ \mathrm{s}$ and $t = 60\ \mathrm{s}$.
\end{enumerate}
\tcblower
\textbf{Solution:}
\begin{enumerate}
\item \textit{(Graph plotted — see figure above for similar shape.)}
\item \textbf{Initial rate:} Draw tangent at $t=0$. Choose points $(0, 1.00)$ and $(20, 0.60)$ on tangent.
Gradient $= \frac{0.60 - 1.00}{20 - 0} = -0.020\ \mathrm{mol\,dm^{-3}\,s^{-1}}$.
Rate of disappearance of $\ce{A} = 0.020\ \mathrm{mol\,dm^{-3}\,s^{-1}}$.
Reaction rate $v = -\frac{1}{2}\frac{d[\ce{A}]}{dt} = \frac{1}{2} \times 0.020 = \mathbf{0.010\ \mathrm{mol\,dm^{-3}\,s^{-1}}}$.
\item \textbf{At $t = 40\ \mathrm{s}$:} $[\ce{A}] \approx 0.45$. Tangent gradient $\approx \frac{0.30 - 0.60}{60 - 20} = -0.0075\ \mathrm{mol\,dm^{-3}\,s^{-1}}$.
$v = \frac{1}{2} \times 0.0075 = \mathbf{3.75 \times 10^{-3}\ \mathrm{mol\,dm^{-3}\,s^{-1}}}$.
\item \textbf{Average rate of formation of $\ce{B}$:}
$\Delta[\ce{A}] = 0.30 - 0.67 = -0.37\ \mathrm{mol\,dm^{-3}}$ over $\Delta t = 40\ \mathrm{s}$.
Average rate of disappearance of $\ce{A} = \frac{0.37}{40} = 9.25 \times 10^{-3}\ \mathrm{mol\,dm^{-3}\,s^{-1}}$.
Average reaction rate $v_{\text{av}} = \frac{1}{2} \times 9.25 \times 10^{-3} = 4.625 \times 10^{-3}\ \mathrm{mol\,dm^{-3}\,s^{-1}}$.
Rate of formation of $\ce{B} = +\frac{d[\ce{B}]}{dt} = 1 \times v_{\text{av}} = \mathbf{4.6 \times 10^{-3}\ \mathrm{mol\,dm^{-3}\,s^{-1}}}$ (2 s.f.).
\end{enumerate}
\end{exoresolu}

\courssub{Summary of Key Points}

\begin{aretenir}[Foundations of Reaction Kinetics]
\begin{itemize}
\item Rate of reaction $v = \frac{1}{\nu_i}\frac{d[\ce{C_i}]}{dt}$ (positive for products, negative for reactants). Units: $\mathrm{mol\,dm^{-3}\,s^{-1}}$.
\item \textbf{Average rate} over $\Delta t$: $\frac{\Delta[\text{species}]}{\Delta t}$ (with stoichiometric factor).
\item \textbf{Instantaneous rate} at $t$: gradient of tangent to concentration–time curve at $t$.
\item \textbf{Initial rate}: instantaneous rate at $t=0$; depends only on initial concentrations.
\item Factors increasing rate: higher concentration/pressure, higher temperature, larger surface area (heterogeneous), catalyst (lowers $E_a$).
\item Applications: industrial synthesis, environmental chemistry, biochemistry, food technology, safety.
\end{itemize}
\end{aretenir}

\begin{exercice}[Practice Questions]
\begin{enumerate}
\item For the reaction $\ce{N2(g) + 3H2(g) -> 2NH3(g)}$, if the rate of formation of $\ce{NH3}$ is $2.0 \times 10^{-4}\ \mathrm{mol\,dm^{-3}\,s^{-1}}$, what are the rates of disappearance of $\ce{N2}$ and $\ce{H2}$?
\item A student measures the volume of $\ce{CO2}$ evolved when $\ce{CaCO3}$ reacts with excess $\ce{HCl}$. Explain why the rate of reaction decreases with time even though the temperature is constant.
\item Sketch a concentration–time curve for a reactant in a first-order reaction. On the same axes, show how the initial rate and the rate at half-life would be determined graphically.
\item In the Haber process, the reaction is carried out at $450\ \mathrm{^\circ C}$ and $200\ \mathrm{atm}$ with an iron catalyst. Explain, using kinetic theory, the effect of each of these three conditions on the rate of reaction.
\end{enumerate}
\end{exercice}

\begin{corrige}[Exercise 1]
\begin{enumerate}
\item $v = +\frac{1}{2}\frac{d[\ce{NH3}]}{dt} = 1.0 \times 10^{-4}\ \mathrm{mol\,dm^{-3}\,s^{-1}}$.
Rate of disappearance of $\ce{N2} = -\frac{d[\ce{N2}]}{dt} = 1 \times v = \mathbf{1.0 \times 10^{-4}\ \mathrm{mol\,dm^{-3}\,s^{-1}}}$.
Rate of disappearance of $\ce{H2} = -\frac{d[\ce{H2}]}{dt} = 3 \times v = \mathbf{3.0 \times 10^{-4}\ \mathrm{mol\,dm^{-3}\,s^{-1}}}$.
\item As $\ce{HCl}$ is consumed, its concentration decreases. Fewer $\ce{H+}$ ions per unit volume $\rightarrow$ fewer effective collisions per second with the $\ce{CaCO3}$ surface $\rightarrow$ rate decreases.
\item Curve: exponential decay. Tangent at $t=0$ (steepest). Tangent at $t=t_{1/2}$ where $[\ce{A}] = [\ce{A}]_0/2$ (less steep).
\item $450\ \mathrm{^\circ C}$: high $T$ $\rightarrow$ more molecules with $E \ge E_a$ $\rightarrow$ higher rate (but compromise with equilibrium yield). $200\ \mathrm{atm}$: high pressure $\rightarrow$ high concentration of gases $\rightarrow$ more frequent collisions $\rightarrow$ higher rate. Fe catalyst: provides alternative pathway with lower $E_a$ $\rightarrow$ higher rate at same $T$.
\end{enumerate}
\end{corrige}

\coursec{Chemical Method – Titrimetric Rate Determination}

In the previous section, we established the foundations of reaction kinetics: we defined the rate of reaction, explored the factors that influence it (concentration, temperature, catalysts, surface area), and learned how to express rate mathematically. We also saw that to determine a rate law experimentally, we need reliable concentration–time data. But how do we actually \emph{obtain} that data for a reaction happening in solution?

Imagine you are studying the hydrolysis of an ester, \ce{CH3COOCH2CH3 + H2O -> CH3COOH + CH3CH2OH}, catalysed by acid. The reaction proceeds slowly enough to be followed over minutes or hours, but you cannot "see" the concentration changing directly. You need a way to "freeze" the reaction at specific instants, analyse the composition, and then plot concentration versus time. This is the essence of the \textbf{chemical (titrimetric) method} — a classic, robust technique that remains a cornerstone of kinetic studies in our laboratories.

\courssub{Principle: Quenching and Aliquot Analysis}

The titrimetric method relies on a simple but powerful idea: \textbf{remove a small portion (aliquot) of the reaction mixture at a known time, stop the reaction instantly in that portion, and then determine the concentration of a reactant or product by titration.}

\begin{definition}[Quenching]
\textbf{Quenching} is the process of rapidly stopping (or drastically slowing down) a chemical reaction in a sampled aliquot so that its composition remains essentially unchanged from the moment of sampling until analysis. Common quenching techniques include:
\begin{itemize}
    \item \textbf{Cooling:} Plunging the aliquot into an ice–water bath ($0\,^\circ\mathrm{C}$) to reduce the rate constant $k$ to near zero.
    \item \textbf{Dilution:} Adding the aliquot to a large volume of cold water or solvent, reducing concentrations and thus the rate.
    \item \textbf{Neutralisation / Reagent addition:} For acid-catalysed reactions, adding excess standard base to neutralise the catalyst; for reactions involving a specific reagent, adding a scavenger that consumes a reactant instantly.
\end{itemize}
The key requirement: \textbf{quenching must be effectively instantaneous} compared to the reaction half-life. If the reaction continues during quenching, the measured concentration will not correspond to the true concentration at the sampling time $t$.
\end{definition}

\begin{popfigure}
\begin{tikzpicture}[scale=0.9, every node/.style={font=\small}]
    % Reaction Flask
    \draw[thick] (0,0) -- (0,5) -- (1,6) -- (3,6) -- (4,5) -- (4,0);
    \draw[thick] (1,6) -- (1,6.5) -- (3,6.5) -- (3,6); % neck
    \draw[fill=popblue!20] (0,0) rectangle (4,2.5); % reaction mixture
    \node at (2,1.25) {\textbf{Reaction Mixture}};
    \node at (2,4) {\ce{A + B -> Products}};
    \draw[->, thick, poporange] (4.5,3) -- (5.5,3) node[midway, above] {Aliquot (pipette)};
    
    % Ice Bath / Quenching Vessel
    \draw[thick] (6,0) -- (6,4) -- (7,4.5) -- (9,4.5) -- (10,4) -- (10,0);
    \draw[pattern=north east lines, pattern color=popblue] (6,0) rectangle (10,1.5); % ice
    \draw[fill=popgreen!20] (6,1.5) rectangle (10,3.5); % quenched aliquot
    \node at (8,2.5) {\textbf{Quenched Aliquot}};
    \node at (8,0.7) {\textbf{Ice Bath ($0\,^\circ\mathrm{C}$)}};
    \draw[->, thick, poporange] (10.5,3) -- (11.5,3) node[midway, above] {To Titration};
    
    % Titration Setup (Burette + Conical Flask)
    \draw[thick] (12.5,0) -- (12.5,6) -- (13,6.5) -- (15,6.5) -- (15.5,6) -- (15.5,0); % Burette stand
    \draw[thick] (13.5,6.5) -- (13.5,1); % Burette tube
    \draw[thick] (13.2,1) -- (13.8,1); % Tip
    \draw[fill=poppurple!30] (13.5,1) rectangle (13.5,4); % Burette solution (simplified)
    \node at (14.5,5) {\textbf{Burette: Standard Soln}};
    \node at (14.5,4.5) {\ce{(e.g. NaOH, Thiosulfate)}};
    
    \draw[thick] (12, -0.5) ellipse (1.5 and 0.4); % Conical flask top
    \draw[thick] (10.5, -0.5) -- (10.5, -2.5) -- (13.5, -2.5) -- (13.5, -0.5); % Conical flask body
    \draw[fill=poppink!30] (10.5, -0.5) rectangle (13.5, -1.8); % Titration mixture
    \node at (12, -1.15) {\textbf{Quenched Aliquot + Indicator}};
    \node at (12, -3.2) {\textbf{White Tile}};
    
    % Indicator drop
    \draw[fill=popgold] (11.5, 0) circle (0.1);
    \node[right] at (11.6, 0) {Ind.};
\end{tikzpicture}
\legende{Schematic of the quenching–titration workflow: (1) Reaction flask maintained at constant temperature; (2) Aliquot withdrawn by pipette and immediately quenched in ice bath; (3) Quenched aliquot titrated against a standard solution with a suitable indicator.}
\end{popfigure}

\courssub{Experimental Procedure: Step by Step}

A typical experiment follows this rigorous protocol:

\begin{methode}[Titrimetric Rate Determination Protocol]
\begin{enumerate}
    \item \textbf{Prepare reactants:} Make up solutions of known initial concentrations. If temperature dependence is studied, thermostat the reactants separately at the chosen temperature (e.g., $25.0\,^\circ\mathrm{C}$) in a water bath.
    \item \textbf{Initiate reaction:} Mix reactants in the reaction flask (start the clock, $t=0$). Swirl gently to ensure homogeneity. Maintain temperature control.
    \item \textbf{Sampling (Aliquot withdrawal):} At predetermined times $t_1, t_2, t_3, \dots$, use a clean, dry pipette (often $10.0\,\mathrm{cm}^3$ or $25.0\,\mathrm{cm}^3$) to withdraw an aliquot. \textbf{Immediately} transfer it to the quenching vessel.
    \item \textbf{Quenching:} Ensure instantaneous quenching. For an ice bath, the quenching vessel should be pre-cooled and contain enough ice-water slurry to bring the aliquot to $0\,^\circ\mathrm{C}$ within seconds.
    \item \textbf{Titration:} Transfer the quenched aliquot (or titrate directly in the quenching vessel if volume permits) to a conical flask. Add 2–3 drops of a suitable indicator. Titrate against a standard solution of known concentration $C_{\text{std}}$ until the endpoint.
    \item \textbf{Record data:} Note the time $t$ and the titre volume $V_{\text{std}}$ for each aliquot. Also perform a "$t=\infty$" titration (after reaction is complete, e.g., 10 half-lives) to find the final concentration.
    \item \textbf{Repeat:} Carry out at least two runs for reproducibility.
\end{enumerate}
\end{methode}

\courssub{Stoichiometric Calculations: From Titre to Concentration}

The raw data are times $t$ and titre volumes $V_{\text{std}}$. We must convert these into concentrations $[\ce{A}]_t$ of the species being monitored (reactant or product).

\begin{methode}[Calculating Concentration from Titration Data]
Let the reaction be:
\[ a\ce{A} + b\ce{B} \rightarrow \text{products} \]
Suppose we titrate species \ce{A} (reactant) with standard titrant \ce{T} of concentration $C_{\text{std}}$.
The titration reaction is: \ce{p A + q T -> products}.
\begin{enumerate}
    \item Moles of titrant used at time $t$: $n_{\ce{T}} = C_{\text{std}} \times V_{\text{std}}(t)$ (in $\mathrm{dm}^3$).
    \item Moles of \ce{A} in the aliquot: $n_{\ce{A}}(t) = \frac{p}{q} \times n_{\ce{T}} = \frac{p}{q} C_{\text{std}} V_{\text{std}}(t)$.
    \item Concentration of \ce{A} in the \emph{original reaction mixture} at time $t$:
    \[ [\ce{A}]_t = \frac{n_{\ce{A}}(t)}{V_{\text{aliquot}}} = \frac{p}{q} \frac{C_{\text{std}} V_{\text{std}}(t)}{V_{\text{aliquot}}} \]
    where $V_{\text{aliquot}}$ is the pipette volume (in $\mathrm{dm}^3$).
\end{enumerate}
If titrating a \textbf{product} \ce{P} formed with stoichiometry $a\ce{A} \rightarrow r\ce{P}$, then $[\ce{A}]_t = [\ce{A}]_0 - \frac{a}{r}[\ce{P}]_t$.
\end{methode}

\begin{propriete}[Concentration–Time Table Construction]
After calculating $[\ce{A}]_t$ for each time point, construct a table:
\begin{center}
\begin{tabular}{|c|c|c|c|}\hline
\rowcolor{popblueL}
\textbf{Time $t$ / s} & \textbf{Titre $V_{\text{std}}$ / $\mathrm{cm}^3$} & \textbf{$[\ce{A}]_t$ / $\mathrm{mol\,dm^{-3}}$} & \textbf{$\ln[\ce{A}]_t$ or $1/[\ce{A}]_t$} \\ \hline
0 & $V_0$ & $[\ce{A}]_0$ & $\ln[\ce{A}]_0$ \\ \hline
$t_1$ & $V_1$ & $[\ce{A}]_1$ & $\ln[\ce{A}]_1$ \\ \hline
$t_2$ & $V_2$ & $[\ce{A}]_2$ & $\ln[\ce{A}]_2$ \\ \hline
$\vdots$ & $\vdots$ & $\vdots$ & $\vdots$ \\ \hline
$\infty$ & $V_\infty$ & $[\ce{A}]_\infty$ (often 0) & -- \\ \hline
\end{tabular}
\end{center}
This table is the raw material for graphical determination of order (Section 4).
\end{propriete}

\courssub{Choice of Indicator and Standard Solution}

The success of the method hinges on a sharp, reproducible endpoint.

\begin{itemize}
    \item \textbf{Acid–base titrations:} Phenolphthalein (pH 8.2–10, colourless to pink) for strong base vs weak/strong acid; Methyl orange (pH 3.1–4.4, red to yellow) for strong acid vs weak base. \textbf{Warning:} In kinetic runs where pH changes during reaction (e.g., acid-catalysed ester hydrolysis producing acid), the indicator must be chosen for the \emph{final} pH of the quenched aliquot.
    \item \textbf{Redox titrations:} Starch indicator (blue-black complex with \ce{I2}) for thiosulfate–iodine titrations (\ce{I2 + 2S2O3^{2-} -> 2I- + S4O6^{2-}}). Add starch \emph{near} the endpoint (pale yellow) to avoid adsorption errors.
    \item \textbf{Standard solutions:} Must be primary standards (stable, high purity, non-hygroscopic) or standardised against one. Examples: \ce{Na2CO3} (acid–base), \ce{K2Cr2O7} or \ce{Na2S2O3} (redox). Concentration known to $\pm 0.1\%$.
\end{itemize}

\courssub{Worked Example: Acid-Catalysed Hydrolysis of Ethyl Ethanoate}

\begin{exoresolu}[Determining Rate Constant from Titration Data]
\textbf{Statement:} The acid-catalysed hydrolysis of ethyl ethanoate is studied at $25\,^\circ\mathrm{C}$:
\ce{CH3COOCH2CH3(l) + H2O(l) ->[H+] CH3COOH(aq) + CH3CH2OH(aq)}.
A reaction mixture is prepared by mixing $50.0\,\mathrm{cm}^3$ of $0.200\,\mathrm{mol\,dm^{-3}}$ ester with $50.0\,\mathrm{cm}^3$ of $0.100\,\mathrm{mol\,dm^{-3}}$ \ce{HCl} (catalyst). At various times, $10.0\,\mathrm{cm}^3$ aliquots are withdrawn, quenched in ice, and titrated against $0.100\,\mathrm{mol\,dm^{-3}}$ \ce{NaOH} using phenolphthalein. The titres are:

\begin{center}
\begin{tabular}{|c|c|c|c|c|c|}\hline
$t$ / min & 0 & 10 & 20 & 30 & $\infty$ \\ \hline
$V_{\ce{NaOH}}$ / $\mathrm{cm}^3$ & 15.0 & 20.5 & 24.8 & 28.2 & 40.0 \\ \hline
\end{tabular}
\end{center}

\begin{enumerate}
    \item Explain why the $t=0$ titre is not zero.
    \item Calculate $[\ce{CH3COOH}]_t$ at each time.
    \item Show that the reaction is first order with respect to ester and determine the rate constant $k$ (in $\mathrm{s^{-1}}$).
\end{enumerate}

\tcblower
\textbf{Detailed Solution:}

\textbf{1. Why $V_0 \neq 0$?}
At $t=0$, the mixture already contains the catalyst \ce{HCl} ($0.100\,\mathrm{mol\,dm^{-3}}$ before mixing, $0.050\,\mathrm{mol\,dm^{-3}}$ after mixing $50+50\,\mathrm{cm}^3$). The titration with \ce{NaOH} neutralises \textbf{all free acid} present: the catalyst \ce{HCl} \emph{plus} any acetic acid formed. So $V_0$ corresponds to the initial catalyst concentration.

\textbf{2. Standard Approach for Ester Hydrolysis.}
The titration measures \textbf{total acid} (\ce{HCl} + \ce{CH3COOH}). Both react 1:1 with \ce{NaOH}.
Let $V_{\text{catalyst}} = V_0$ be the volume of \ce{NaOH} needed to neutralise \ce{HCl} alone in the $10.0\,\mathrm{cm}^3$ aliquot.
Let $V_\infty$ be the titre at completion (all ester hydrolysed).
Then:
\begin{itemize}
    \item Volume of \ce{NaOH} proportional to acetic acid formed at time $t$ = $V_t - V_0$.
    \item Volume of \ce{NaOH} proportional to ester \textbf{remaining} at time $t$ = $V_\infty - V_t$.
\end{itemize}
Since $[\text{Ester}]_t \propto (V_\infty - V_t)$, we can use $(V_\infty - V_t)$ directly in the first-order integrated rate law.

\begin{center}
\begin{tabular}{|c|c|c|c|}\hline
\rowcolor{popblueL}
$t$ / min & $V_t$ / $\mathrm{cm}^3$ & $V_\infty - V_t$ (prop. to [Ester]) & $\ln(V_\infty - V_t)$ \\ \hline
0 & 15.0 & 25.0 & 3.219 \\ \hline
10 & 20.5 & 19.5 & 2.970 \\ \hline
20 & 24.8 & 15.2 & 2.721 \\ \hline
30 & 28.2 & 11.8 & 2.468 \\ \hline
\end{tabular}
\end{center}

\textbf{3. Graphical Test for First Order.}
For a first-order reaction: $\ln[\ce{A}]_t = \ln[\ce{A}]_0 - kt$.
Since $[\ce{A}]_t \propto (V_\infty - V_t)$, a plot of $\ln(V_\infty - V_t)$ vs $t$ should be linear with slope $-k$.

\begin{center}
\begin{tikzpicture}[scale=0.7]
\begin{axis}[
    width=10cm, height=6cm,
    xlabel={Time $t$ / min}, ylabel={$\ln(V_\infty - V_t)$},
    grid=major, grid style={popmuted},
    xmin=0, xmax=35, ymin=2.3, ymax=3.4,
    scatter, only marks, mark=*, mark size=2, popblue
]
\addplot coordinates {(0,3.219) (10,2.970) (20,2.721) (30,2.468)};
\addplot[domain=0:30, samples=2, thick, poporange, dashed] {3.22 - 0.0253*x}; % approximate fit
\end{axis}
\end{tikzpicture}
\end{center}
The points fall on a straight line $\rightarrow$ \textbf{First order confirmed}.

\textbf{Slope calculation:}
Using $t=0$ and $t=30$ min:
Slope $= \frac{2.468 - 3.219}{30 - 0} = \frac{-0.751}{30} = -0.02503\,\mathrm{min^{-1}}$.
$k = 0.0250\,\mathrm{min^{-1}}$.
Convert to $\mathrm{s^{-1}}$: $k = \frac{0.0250}{60} = \mathbf{4.17 \times 10^{-4}\,\mathrm{s^{-1}}}$.
(Using linear regression on all points gives $k \approx 4.2 \times 10^{-4}\,\mathrm{s^{-1}}$).
\end{exoresolu}

\courssub{Sources of Error and Precautions}

Even with careful technique, the titrimetric method has inherent limitations. Understanding these is crucial for the \emph{Planning} and \emph{Analysis} papers.

\begin{attention}[Common Pitfalls in Titrimetric Kinetics]
\begin{enumerate}
    \item \textbf{Incomplete / Slow Quenching:} If the aliquot takes 30 seconds to cool to $0\,^\circ\mathrm{C}$, and the half-life is 2 minutes, significant reaction occurs during quenching. The measured $[A]_t$ will be \textbf{lower} than the true value (for a reactant). \textbf{Remedy:} Use a large ice bath, small aliquot volume, or chemical quenching (e.g., adding excess base to acid-catalysed reaction).
    \item \textbf{Sampling Volume Error:} Using a $25\,\mathrm{cm}^3$ pipette removes a significant fraction of the reaction mixture, changing the total volume and concentrations over the course of the experiment. \textbf{Remedy:} Use a $10\,\mathrm{cm}^3$ pipette; keep total volume large (e.g., $250\,\mathrm{cm}^3$); or correct mathematically for volume loss.
    \item \textbf{Endpoint Detection:} Phenolphthalein endpoint is subtle (colourless to faint pink). Over-titration adds excess \ce{NaOH}, overestimating acid concentration. \textbf{Remedy:} Use a white tile, good lighting, and titrate dropwise near the end. Carry out a "blank" titration of the indicator in water to correct for indicator blank.
    \item \textbf{Indicator Interference:} In redox kinetics (e.g., \ce{S2O8^{2-}/I-} with thiosulfate), starch must be added \emph{only} when the solution is pale yellow (low \ce{I2}). Adding it too early forms a stable starch–iodine complex that is hard to titrate.
    \item \textbf{Standard Solution Stability:} \ce{NaOH} absorbs \ce{CO2} from air, forming \ce{Na2CO3}, which changes its concentration and gives a double endpoint with phenolphthalein/methyl orange. \textbf{Remedy:} Standardise \ce{NaOH} frequently against \ce{KHP} (potassium hydrogen phthalate); store in a bottle with a \ce{CO2} guard tube (soda lime).
\end{enumerate}
\end{attention}

\courssub{Advantages and Limitations}

\begin{center}
\begin{tabularx}{\linewidth}{|l|X|X|}\hline
\rowcolor{popblueL}
\textbf{Aspect} & \textbf{Advantages} & \textbf{Limitations} \\ \hline
\textbf{Equipment} & Simple, inexpensive glassware (burette, pipette, flask). Available in every school lab. & Labour-intensive; manual sampling limits time resolution (cannot follow very fast reactions $< 1$ min). \\ \hline
\textbf{Versatility} & Works for any reaction where a reactant or product can be titrated (acid–base, redox, complexometric, precipitation). & Requires a suitable titration reaction; not universal (e.g., no simple titration for \ce{N2O5} decomposition). \\ \hline
\textbf{Accuracy} & High precision possible ($\pm 0.1\%$) with good technique and primary standards. & Cumulative errors: sampling + quenching + titration + indicator. \\ \hline
\textbf{Data Quality} & Gives absolute concentrations directly ($\mathrm{mol\,dm^{-3}}$), not just relative signals. & Discrete data points only; cannot capture continuous concentration profile. \\ \hline
\end{tabularx}
\end{center}

\courssub{Connection to Rate Law Determination}

The concentration–time data $(t, [\ce{A}]_t)$ obtained here are the \textbf{raw material} for the next section. We will use them to:
\begin{itemize}
    \item Plot $[\ce{A}]_t$ vs $t$ $\rightarrow$ tangent method for instantaneous rates.
    \item Plot $\ln[\ce{A}]_t$ vs $t$ and $1/[\ce{A}]_t$ vs $t$ $\rightarrow$ graphical determination of order (Section 4, Lesson 97).
    \item Use initial rates from different initial concentrations $\rightarrow$ initial rate method (Section 4, Lesson 96).
\end{itemize}

\begin{aretenir}[Key Points for Examination]
\begin{itemize}
    \item \textbf{Quenching} must be instantaneous (ice bath, dilution, neutralisation). Incomplete quenching $\rightarrow$ systematic error (rate appears slower).
    \item \textbf{Concentration calculation:} $[\ce{A}]_t = \frac{\text{stoich factor} \times C_{\text{std}} \times V_{\text{std}}(t)}{V_{\text{aliquot}}}$. For ester hydrolysis: $[\text{Ester}]_t \propto (V_\infty - V_t)$.
    \item \textbf{First-order test:} $\ln(V_\infty - V_t)$ vs $t$ is linear; slope $= -k$.
    \item \textbf{Major errors:} Volume change due to sampling, indicator blank, \ce{CO2} absorption by \ce{NaOH}, slow quenching.
    \item \textbf{Units of $k$:} First order $\rightarrow$ $\mathrm{s^{-1}}$; Second order $\rightarrow$ $\mathrm{dm^3\,mol^{-1}\,s^{-1}}$.
\end{itemize}
\end{aretenir}

\begin{savaistu}[Titrimetric Kinetics in Cameroon's Industries]
In the \textbf{CDC (Cameroon Development Corporation)} rubber and palm oil processing plants, the \textbf{saponification value} of oils (a measure of average molecular weight) is determined by refluxing oil with excess ethanolic \ce{KOH} and back-titrating the remaining \ce{KOH} with \ce{HCl}. While this is an analytical method, the \emph{kinetics} of saponification (rate of hydrolysis of triglycerides) is crucial for designing the reactors: too fast $\rightarrow$ runaway temperature; too slow $\rightarrow$ low throughput. Chemical engineers use pilot-plant titration data (exactly as described here) to size industrial saponification reactors for soap production in Douala and Bassa.
\end{savaistu}

\begin{exercice}[Practice Problem]
The decomposition of hydrogen peroxide catalysed by \ce{MnO2} is studied: \ce{2H2O2(aq) -> 2H2O(l) + O2(g)}.
At fixed intervals, $10.0\,\mathrm{cm}^3$ aliquots are withdrawn, the \ce{MnO2} is filtered off quickly, and the filtrate is titrated with $0.0200\,\mathrm{mol\,dm^{-3}}$ \ce{KMnO4} in acidic medium: \ce{2MnO4- + 5H2O2 + 6H+ -> 2Mn^{2+} + 5O2 + 8H2O}.
The following titres are obtained:
\begin{center}
\begin{tabular}{|c|c|c|c|c|c|}\hline
$t$ / min & 0 & 5 & 10 & 15 & 20 \\ \hline
$V_{\ce{KMnO4}}$ / $\mathrm{cm}^3$ & 25.0 & 18.4 & 13.6 & 10.0 & 7.4 \\ \hline
\end{tabular}
\end{center}
\begin{enumerate}
    \item Calculate $[\ce{H2O2}]_t$ at each time.
    \item Plot $\ln[\ce{H2O2}]_t$ vs $t$ and determine the order and rate constant $k$.
    \item Suggest a better quenching method than filtration for this heterogeneous catalysis.
\end{enumerate}
\end{exercice}

\begin{corrige}[Exercise 2]
\begin{enumerate}
    \item Stoichiometry: 2 \ce{MnO4-} $\equiv$ 5 \ce{H2O2}. Moles \ce{H2O2} = $\frac{5}{2} \times C_{\ce{KMnO4}} \times V_{\ce{KMnO4}}$.
    $[\ce{H2O2}]_t = \frac{5}{2} \times \frac{0.0200 \times V_{\ce{KMnO4}}(\text{in dm}^3)}{0.0100} = 5.00 \times V_{\ce{KMnO4}}(\text{in dm}^3)$ $\mathrm{mol\,dm^{-3}}$.
    $V$ in $\mathrm{cm}^3$ $\rightarrow$ $[\ce{H2O2}] = 0.00500 \times V_{\ce{KMnO4}}$.
    $t=0$: 0.125 M; $t=5$: 0.0920 M; $t=10$: 0.0680 M; $t=15$: 0.0500 M; $t=20$: 0.0370 M.
    \item $\ln[\ce{H2O2}]$: -2.08, -2.38, -2.69, -2.99, -3.29. Plot is linear $\rightarrow$ First order.
    Slope $\approx \frac{-3.29 - (-2.08)}{20} = -0.0605\,\mathrm{min^{-1}} = 1.01 \times 10^{-3}\,\mathrm{s^{-1}}$.
    \item Filtration is slow. Better: Add a \textbf{complexing agent} (e.g., \ce{EDTA} or phosphate buffer) that instantly binds \ce{Mn^{2+}}/poisons the catalyst surface, or dilute massively into cold acidified solution containing a \ce{MnO2} inhibitor.
\end{enumerate}
\end{corrige}

\coursec{Physical Methods of Rate Measurement}

In the previous section we followed a reaction by the \emph{chemical} method: we withdrew samples at known times, quenched them, and titrated. That method is reliable but slow, destructive (each sample is lost), and impossible for very fast reactions. \cle{Physical methods} overcome these drawbacks: instead of removing samples, we continuously monitor some \emph{physical property} of the reaction mixture that changes in proportion to the concentration of a reactant or product --- its colour, its electrical conductivity, its volume, or the pressure/volume of a gas it produces. The mixture is never disturbed, readings can be taken every few seconds, and a complete rate curve is obtained from a \emph{single} experiment.

The general strategy is always the same:
\begin{enumerate}
\item choose a measurable physical quantity $X$ that is a \emph{known, simple function} of a concentration (ideally proportional to it);
\item calibrate, i.e.\ establish the relation $X = f(c)$ using standard solutions;
\item record $X$ against time during the reaction;
\item convert $X(t)$ into $c(t)$, and hence obtain the rate as the gradient of the concentration--time curve, exactly as in Section~2.
\end{enumerate}

\courssub{Colorimetry}

\textbf{Intuition.}\par
Many reactions involve a coloured species. When purple acidified permanganate oxidises ethanedioic acid,
\[
\ce{2MnO4^{-}(aq) + 5C2O4^{2-}(aq) + 16H^{+}(aq) -> 2Mn^{2+}(aq) + 10CO2(g) + 8H2O(l)},
\]
the solution slowly turns from purple to colourless: the intensity of the colour is a direct, visible measure of $[\ce{MnO4^-}]$. A \cle{colorimeter} (or spectrophotometer) simply makes this judgement quantitative: it shines light of a chosen wavelength through the solution and measures how much is absorbed.

\begin{definition}[Beer--Lambert law]
When monochromatic light passes through a solution of an absorbing species, the \cle{absorbance} $A$ is proportional to the concentration $c$ of the absorbing species and to the path length $l$ of the cell:
\[
A = \varepsilon\, l\, c
\]
where $\varepsilon$ is the \cle{molar absorption coefficient} (units $\mathrm{dm^3\,mol^{-1}\,cm^{-1}}$), $l$ is in $\mathrm{cm}$ and $c$ in $\mathrm{mol\,dm^{-3}}$. Absorbance itself is dimensionless: $A = \log_{10}(I_0/I)$, where $I_0$ and $I$ are the intensities of the incident and transmitted light.
\end{definition}

\begin{propriete}[Consequence for kinetics]
Since $\varepsilon$ and $l$ are constant for a given experiment, $A \propto c$. Therefore a plot of $A$ against time has \emph{exactly the same shape} as the concentration--time curve, and
\[
\text{rate} = -\frac{dc}{dt} = -\frac{1}{\varepsilon l}\frac{dA}{dt}.
\]
In practice one converts $A$ to $c$ with a \cle{calibration curve}: the absorbances of several standard solutions of known concentration are measured, and $A$ is plotted against $c$. The result should be a straight line through the origin, from which any measured $A$ can be read off as a concentration.
\end{propriete}

\begin{attention}[Choosing the wavelength]
Colorimetry requires that \emph{only one species} absorbs significantly at the chosen wavelength. If reactant and product absorb at the same wavelength (overlapping spectra), the measured absorbance no longer tracks a single concentration and the kinetic analysis is wrong. Always choose the wavelength of maximum absorption of the species of interest ($\lambda_{\max}$), and check that other species are transparent there. Also remember that the Beer--Lambert law fails at high concentrations (typically $c > 0.01\ \mathrm{mol\,dm^{-3}}$): work in the linear range of the calibration curve.
\end{attention}

\begin{popfigure}
\begin{center}
\begin{tikzpicture}
\begin{axis}[
  width=0.62\linewidth, height=5.2cm,
  xlabel={$t$ / s}, ylabel={Absorbance $A$},
  xmin=0, xmax=300, ymin=0, ymax=1.1,
  axis lines=left, tick label style={font=\small},
  label style={font=\small},
  title style={font=\small},
  title={Decay of a coloured reactant},
  every axis plot/.append style={thick}]
\addplot[popblue, domain=0:300, samples=80] {exp(-x/90)};
\addplot[poporange, dashed, domain=0:60] {1 - x/90};
\node[font=\small, text=popdark, align=center] at (axis cs:200,0.55) {$A = A_0\,e^{-kt}$\\ gradient at $t=0$\\ $= -kA_0$};
\end{axis}
\end{tikzpicture}\hfill
\begin{tikzpicture}
\begin{axis}[
  width=0.34\linewidth, height=5.2cm,
  xlabel={$c$ / $10^{-4}\ \mathrm{mol\,dm^{-3}}$}, ylabel={$A$},
  xmin=0, xmax=5, ymin=0, ymax=1.1,
  axis lines=left, tick label style={font=\small},
  label style={font=\small},
  title style={font=\small},
  title={Calibration line},
  every axis plot/.append style={thick}]
\addplot[popgreen, domain=0:5] {0.22*x};
\addplot[only marks, mark=*, mark size=1.5pt, popdark] coordinates {(1,0.22) (2,0.44) (3,0.66) (4,0.88)};
\end{axis}
\end{tikzpicture}
\end{center}
\legende{Colorimetric following of a reaction. Left: absorbance of the coloured reactant falls exponentially with time; the initial rate is obtained from the tangent at $t = 0$. Right: the calibration line $A = \varepsilon l c$ used to convert absorbance into concentration.}
\end{popfigure}

\begin{exoresolu}[Using a calibration curve]
The hydrolysis of a coloured dye is followed colorimetrically at its $\lambda_{\max}$. A calibration experiment gives $\varepsilon l = 1.1\times 10^{4}\ \mathrm{dm^3\,mol^{-1}}$. During the reaction the absorbance falls from $A_0 = 0.88$ at $t=0$ to $A = 0.44$ at $t = 120\ \mathrm{s}$. Find (a) the initial concentration of the dye, (b) the average rate of disappearance over the first $120\ \mathrm{s}$.
\tcblower
(a) From $A = \varepsilon l c$:
\[
c_0 = \frac{A_0}{\varepsilon l} = \frac{0.88}{1.1\times10^{4}} = 8.0\times10^{-5}\ \mathrm{mol\,dm^{-3}}.
\]
(b) At $t = 120\ \mathrm{s}$: $c = 0.44/(1.1\times10^{4}) = 4.0\times10^{-5}\ \mathrm{mol\,dm^{-3}}$.
\[
\text{average rate} = -\frac{\Delta c}{\Delta t} = -\frac{(4.0-8.0)\times10^{-5}}{120} = 3.3\times10^{-7}\ \mathrm{mol\,dm^{-3}\,s^{-1}}.
\]
Note that the absorbance halved in $120\ \mathrm{s}$; because $A\propto c$, the concentration also halved --- a hint that the reaction may be first order (constant half-life), to be confirmed in Section~4.
\end{exoresolu}

\courssub{Conductimetry}

\textbf{Intuition.}\par
A solution conducts electricity through its ions. If a reaction \emph{produces} or \emph{consumes} ions, the conductivity of the mixture changes as the reaction proceeds, and two inert electrodes (usually platinum) dipping into the flask give us a continuous, non-destructive concentration probe. A classic example is the alkaline hydrolysis of an ester:
\[
\ce{CH3COOC2H5(aq) + OH^{-}(aq) -> CH3COO^{-}(aq) + C2H5OH(aq)}.
\]
Here the highly mobile \ce{OH-} ion is replaced by the much less mobile \ce{CH3COO-} ion, so the conductivity \emph{falls} steadily during the reaction.

\begin{definition}[Conductivity and molar ionic conductivities]
The \cle{conductivity} $\kappa$ (unit $\mathrm{S\,cm^{-1}}$ or $\mathrm{S\,m^{-1}}$) of an ionic solution is the sum of the contributions of all ions present:
\[
\kappa = \sum_i \lambda_i\, c_i
\]
where $c_i$ is the concentration of ion $i$ and $\lambda_i$ its \cle{molar ionic conductivity} (unit $\mathrm{S\,cm^2\,mol^{-1}}$). Ions such as \ce{H+} ($\lambda \approx 350\ \mathrm{S\,cm^2\,mol^{-1}}$) and \ce{OH-} ($\lambda \approx 198\ \mathrm{S\,cm^2\,mol^{-1}}$) have exceptionally large $\lambda$ values, so reactions involving them give large, easily measured conductivity changes.
\end{definition}

\begin{methode}[From conductivity to concentration]
To convert a conductivity reading $\kappa_t$ into a concentration:
\begin{enumerate}
\item Identify all ions present at time $t$ and write $\kappa_t = \sum_i \lambda_i c_i$.
\item Use the stoichiometry of the reaction to express every $c_i$ in terms of a single variable (e.g.\ the extent of reaction $x$ or the concentration of the limiting reactant).
\item Use the initial and final readings $\kappa_0$ and $\kappa_\infty$ as anchors. If the conductivity change is due solely to the consumption of one ionic reactant (others constant), the fraction of that reactant remaining is
\[
\frac{[\text{reactant}]_t}{[\text{reactant}]_0} = \frac{\kappa_t - \kappa_\infty}{\kappa_0 - \kappa_\infty}.
\]
\item Plot concentration against time and take gradients to obtain rates.
\end{enumerate}
\end{methode}

\begin{exemplebox}[Ester hydrolysis followed conductimetrically]
For $\ce{CH3COOC2H5 + OH- -> CH3COO- + C2H5OH}$ with the ester in large excess (pseudo-first-order conditions): at $t=0$ the conductivity is due essentially to \ce{Na+} (from \ce{NaOH}) and \ce{OH-}; at $t=\infty$ all \ce{OH-} has been replaced by \ce{CH3COO-}. Since $\lambda(\ce{OH-}) > \lambda(\ce{CH3COO-})$, $\kappa$ decreases linearly with the extent of reaction, and $[\ce{OH-}]_t \propto (\kappa_t - \kappa_\infty)$.
\end{exemplebox}

\begin{popfigure}
\begin{center}
\begin{tikzpicture}
\begin{axis}[
  width=0.7\linewidth, height=5.5cm,
  xlabel={$t$ / min}, ylabel={Conductivity $\kappa$ / $\mathrm{mS\,cm^{-1}}$},
  xmin=0, xmax=60, ymin=0, ymax=12,
  axis lines=left, tick label style={font=\small}, label style={font=\small},
  every axis plot/.append style={thick}]
\addplot[poppurple, domain=0:60, samples=80] {2 + 8*exp(-x/18)};
\addplot[popmuted, dashed] coordinates {(0,2) (60,2)};
\node[font=\small, text=popmuted] at (axis cs:45,2.8) {$\kappa_\infty$};
\node[font=\small, text=popdark] at (axis cs:38,8.5) {$\kappa_t = \kappa_\infty + (\kappa_0-\kappa_\infty)e^{-kt}$};
\end{axis}
\end{tikzpicture}
\end{center}
\legende{Conductivity--time trace for the alkaline hydrolysis of ethyl ethanoate: the fast \ce{OH-} ions are replaced by slower \ce{CH3COO-} ions, so $\kappa$ falls towards its final value $\kappa_\infty$.}
\end{popfigure}

\begin{attention}[Conditions for valid conductimetry]
\begin{itemize}
\item Keep the temperature constant (conductivity increases roughly $2\%$ per degree Celsius).
\item Use an AC supply (typically $1\ \mathrm{kHz}$) to avoid electrolysis and polarisation at the electrodes.
\item Ensure the conductivity change is dominated by \emph{one} ionic transformation; reactions between neutral molecules with no ionic change cannot be followed this way.
\item The cell constant must be known or calibrated with a standard solution (e.g.\ $0.01\ \mathrm{M}\ \ce{KCl}$).
\end{itemize}
\end{attention}

\courssub{Dilatometry}

\textbf{Intuition.}\par
Some liquid-phase reactions occur with a small but measurable change in total volume, because reactants and products have different molar volumes. A \cle{dilatometer} --- a reaction bulb fitted with a narrow graduated capillary tube, rather like a giant thermometer --- converts this tiny volume change into a large, readable movement of the meniscus in the capillary.

\begin{propriete}[Principle of dilatometry]
If the volume of the reaction mixture changes linearly with the extent of reaction $x$, i.e.\ $V_t = V_0 + \alpha x$, then the height $h$ of the liquid in the capillary (of uniform cross-section $A$) is a linear measure of $x$:
\[
\frac{x_t}{x_\infty} = \frac{h_t - h_0}{h_\infty - h_0}.
\]
The method has been classically applied to reactions such as the hydrolysis of esters and the inversion of sucrose, where the products occupy a different volume from the reactants.
\end{propriete}

The great advantage is simplicity: no reagent, no sampling, just a ruler reading against time. The limitations are equally clear: the volume change must be large enough to detect (typically $>0.1\ \mathrm{cm^3}$), and temperature must be controlled to a fraction of a degree (e.g.\ $\pm 0.01\ ^{\circ}\mathrm{C}$), because thermal expansion of the solvent would swamp the chemical volume change.

\courssub{Gas Syringe and Pressure Measurement}

\textbf{Intuition.}\par
When a reaction produces a gas, the gas itself is the most natural thing to measure. Think of adding dilute hydrochloric acid to marble chips (a familiar sight in every Cameroonian laboratory):
\[
\ce{CaCO3(s) + 2HCl(aq) -> CaCl2(aq) + H2O(l) + CO2(g)}.
\]
Every molecule of \ce{CO2} released corresponds to one formula unit of \ce{CaCO3} consumed. Two experimental arrangements are used.

\textbf{(a) Gas syringe (constant pressure, variable volume).}\par
The reaction flask is connected to a graduated gas syringe. As gas is evolved, the plunger is pushed out and the \emph{volume} of gas is read at recorded times. The experiment is done at atmospheric pressure and room temperature. This method is ideal for reactions of moderate speed.

\textbf{(b) Pressure measurement (constant volume, variable pressure).}\par
The reaction is carried out in a \emph{sealed} vessel of fixed volume connected to a manometer or pressure sensor. As gas accumulates, the pressure rises, and the pressure--time trace replaces the volume--time trace. This method is more sensitive and can follow very fast gas-evolving reactions.

\begin{propriete}[From gas measurements to moles]
In both cases the amount of gas is obtained from the \cle{ideal gas law}:
\[
n = \frac{PV}{RT}, \qquad R = 8.314\ \mathrm{J\,K^{-1}\,mol^{-1}}.
\]
Gas syringe: $P$ and $T$ are constant, so $n \propto V$. Pressure method: $V$ and $T$ are constant, so $n \propto P$. Plotting $n$ (or directly $V$ or $P$) against time gives the rate curve; the gradient at any point is the rate of gas production, and the \emph{initial} gradient gives the initial rate.
\end{propriete}

\begin{popfigure}
\begin{center}
\begin{tikzpicture}[scale=1.0, every node/.style={font=\small}]
% Flask
\draw[thick, popdark] (0,0) -- (0,1.6);
\draw[thick, popdark] (0,0) .. controls (0.1,-0.9) and (1.9,-0.9) .. (2,0) -- (2,1.6);
\draw[thick, popdark] (0,1.6) -- (0,2.0);
\draw[thick, popdark] (2,1.6) -- (2,2.0);
\draw[thick, popdark] (0,2.0) -- (2,2.0);
% Liquid
\fill[popblueL] (0.06,-0.35) .. controls (0.15,-0.8) and (1.85,-0.8) .. (1.94,-0.35) -- cycle;
\draw[popblue, thick] (0.06,-0.35) .. controls (0.7,-0.25) and (1.3,-0.45) .. (1.94,-0.35);
\node[text=popdark] at (1,-0.55) {\ce{CaCO3} + \ce{HCl}};
% Bubbles
\foreach \p in {(0.5,0.3),(1.0,0.7),(1.5,0.4),(0.8,1.1),(1.3,1.3)} {\draw[popblue] \p circle (0.06);}
% Delivery tube
\draw[thick, popdark] (1,2.0) -- (1,2.6) -- (3.0,2.6);
% Syringe barrel
\draw[thick, popdark] (3.0,2.3) rectangle (6.2,2.9);
% Plunger
\draw[thick, poporange, fill=popgold!40] (5.2,2.28) rectangle (5.45,2.92);
\draw[thick, poporange] (5.45,2.6) -- (6.6,2.6);
% Graduations
\foreach \x in {3.2,3.6,4.0,4.4,4.8} {\draw[popdark] (\x,2.9) -- (\x,2.75);}
\node[text=popdark] at (4.1,3.25) {graduated gas syringe};
\node[text=popdark] at (4.0,2.05) {gas collected: $V$ / $\mathrm{cm^3}$};
\node[text=popdark] at (1,2.9) {delivery tube};
\node[text=popdark] at (1,-1.25) {conical flask (reaction mixture)};
\end{tikzpicture}
\end{center}
\legende{Gas syringe method: the \ce{CO2} evolved pushes the plunger outwards; the volume is read on the graduated barrel at measured times. Plotting $V$ (or $n = PV/RT$) against time gives the rate curve.}
\end{popfigure}

\begin{exoresolu}[Pressure method and the ideal gas law]
A $0.500\ \mathrm{g}$ sample of magnesium ribbon reacts with excess dilute hydrochloric acid in a sealed flask of volume $250\ \mathrm{cm^3}$ at $25\ ^{\circ}\mathrm{C}$:
\[
\ce{Mg(s) + 2HCl(aq) -> MgCl2(aq) + H2(g)}.
\]
(a) Calculate the final pressure rise due to the \ce{H2} produced. (b) The pressure rises by $50\ \mathrm{kPa}$ after $40\ \mathrm{s}$; find the average rate of \ce{H2} production in $\mathrm{mol\,s^{-1}}$ over this interval. ($M(\ce{Mg}) = 24.3\ \mathrm{g\,mol^{-1}}$; gas volume of the flask available: $250\ \mathrm{cm^3}$.)
\tcblower
(a) Moles of Mg: $n = 0.500/24.3 = 2.06\times10^{-2}\ \mathrm{mol}$, so $n(\ce{H2}) = 2.06\times10^{-2}\ \mathrm{mol}$ (1:1 stoichiometry, acid in excess). Then
\[
P = \frac{nRT}{V} = \frac{2.06\times10^{-2}\times 8.314\times 298}{250\times10^{-6}} = 2.04\times10^{5}\ \mathrm{Pa} \approx 204\ \mathrm{kPa}.
\]
(b) A rise of $50\ \mathrm{kPa}$ corresponds to
\[
n(\ce{H2}) = \frac{PV}{RT} = \frac{50\times10^{3}\times250\times10^{-6}}{8.314\times298} = 5.04\times10^{-3}\ \mathrm{mol},
\]
so the average rate is $5.04\times10^{-3}/40 = 1.3\times10^{-4}\ \mathrm{mol\,s^{-1}}$.
\end{exoresolu}

\begin{attention}[Pitfalls of gas methods]
\begin{itemize}
\item Check for leaks: a poorly fitting bung or syringe gives low, erratic readings.
\item The gas must be (nearly) insoluble in the reaction mixture; \ce{CO2} dissolves appreciably in water, so early readings may underestimate the rate. Use a saturated solution or correct for solubility.
\item Record temperature and atmospheric pressure --- they are needed for $n = PV/RT$.
\item The syringe plunger must move freely; friction delays the reading.
\item Start timing at the moment of mixing, and mix quickly and thoroughly.
\end{itemize}
\end{attention}

\courssub{Choosing the Right Physical Method}

\begin{center}
\begin{tabularx}{\linewidth}{|l|X|X|}
\hline
\rowcolor{popblueL} \textbf{Method} & \textbf{Quantity measured} & \textbf{Best suited for} \\
\hline
Colorimetry & Absorbance $A=\varepsilon l c$ & Reactions with one coloured species \\
\hline
Conductimetry & $\kappa = \sum \lambda_i c_i$ & Reactions producing/consuming ions \\
\hline
Dilatometry & Liquid level in capillary & Liquid reactions with volume change \\
\hline
Gas syringe & Gas volume at constant $P$ & Gas-evolving reactions (moderate speed) \\
\hline
Pressure sensor & Gas pressure at constant $V$ & Gas-evolving reactions (fast or slow) \\
\hline
\end{tabularx}
\end{center}

\begin{aretenir}[Key points of Section 3]
\begin{itemize}
\item Physical methods monitor a property proportional to concentration \emph{continuously} and without disturbing the mixture.
\item Colorimetry: $A = \varepsilon l c$; use a calibration line and a wavelength where only one species absorbs ($\lambda_{\max}$).
\item Conductimetry: $\kappa = \sum_i \lambda_i c_i$; large effects when \ce{H+} or \ce{OH-} are involved; use AC current and thermostatting.
\item Dilatometry: capillary rise measures extent of reaction via volume change; needs excellent temperature control ($\pm 0.01\ ^{\circ}\mathrm{C}$).
\item Gas methods: $n = PV/RT$; plot $n$, $V$ or $P$ against time and take gradients for rates.
\item Whatever the method, the final step is identical: convert the measured quantity to concentration, plot $c$ vs $t$, and measure the gradient --- which prepares us for the rate laws of Section~4.
\end{itemize}
\end{aretenir}

\begin{exercice}[Choosing and exploiting a method]
For each reaction, state the most appropriate physical method and justify your choice:
\begin{enumerate}
\item $\ce{H2O2(aq) -> H2O(l) + 1/2 O2(g)}$ (catalysed by \ce{MnO2});
\item the reaction between acidified \ce{KMnO4} and \ce{H2C2O4};
\item the hydrolysis of \ce{CH3COOCH3} by dilute \ce{NaOH}.
\end{enumerate}
\end{exercice}

\begin{corrige}[Exercise 1]
\begin{enumerate}
\item Gas syringe (or pressure sensor): \ce{O2} gas is evolved and is insoluble enough; $n(\ce{O2}) = PV/RT$.
\item Colorimetry: \ce{MnO4-} is intensely purple and the products are colourless; follow $A$ at $\lambda_{\max}$ of \ce{MnO4-} and use $A = \varepsilon l c$.
\item Conductimetry: the fast \ce{OH-} ions are replaced by slower \ce{CH3COO-} ions, so $\kappa$ decreases measurably; $[\ce{OH-}]_t \propto (\kappa_t - \kappa_\infty)$.
\end{enumerate}
\end{corrige}

\coursec{Rate Law, Order of Reaction, and Determination Methods}

In the previous section we learned how to \emph{measure} how fast a reaction goes: by titrating samples, by following colour, conductivity, volume or pressure. But a list of measurements is not yet a law of nature. The real power of kinetics appears when we compress all those measurements into a single mathematical sentence: the \cle{rate law}. This section is the heart of the chapter. Once you can write $\text{rate} = k[\ce{A}]^m[\ce{B}]^n$ and determine $k$, $m$ and $n$ from experimental data, you can predict the rate of a reaction under \emph{any} conditions --- exactly what a chemical engineer designing a plant in Douala or Limbe needs.

\courssub{The Rate Law and the Rate Constant}

\textbf{From observation to law.}\par
Consider the decomposition of hydrogen peroxide, a reaction you can follow easily in the school laboratory with a gas syringe:
\[
\ce{2H2O2(aq) -> 2H2O(l) + O2(g)}
\]
If you double the concentration of \ce{H2O2} (keeping temperature and catalyst fixed), the initial rate doubles. If you triple it, the rate triples. The rate is \emph{directly proportional} to $[\ce{H2O2}]$. For other reactions, doubling a concentration may quadruple the rate, or may change nothing at all. Experiment --- not the balanced equation --- tells us how the rate depends on each concentration.

\begin{definition}[Rate of reaction and rate law]
For a reaction $\ce{aA + bB -> products}$, the \cle{rate of reaction} $v$ is defined as
\[
v = -\frac{1}{a}\frac{d[\ce{A}]}{dt} = -\frac{1}{b}\frac{d[\ce{B}]}{dt} = \frac{1}{c}\frac{d[\ce{C}]}{dt} = \dots
\]
with units $\si{mol.dm^{-3}.s^{-1}}$. The \cle{rate law} (or \cle{rate equation}) is the experimentally determined relationship between this rate and the concentrations of the species that influence it:
\[
v = k[\ce{A}]^{m}[\ce{B}]^{n}
\]
where:
\begin{itemize}
\item $k$ is the \cle{rate constant} (or rate coefficient);
\item $m$ is the \cle{partial order} of reaction with respect to \ce{A};
\item $n$ is the partial order with respect to \ce{B};
\item $m + n$ is the \cle{overall order} of the reaction.
\end{itemize}
The exponents $m$ and $n$ are found \textbf{only by experiment}; they are \emph{not} the stoichiometric coefficients $a$ and $b$.
\end{definition}

\begin{attention}[Orders are not coefficients!]
For $\ce{2N2O5(g) -> 4NO2(g) + O2(g)}$, experiment gives $v = k[\ce{N2O5}]$, first order --- the coefficient 2 does \emph{not} appear as an exponent. For $\ce{NO2(g) + CO(g) -> NO(g) + CO2(g)}$, experiment gives $v = k[\ce{NO2}]^2$: the reaction is \emph{zero order in CO} even though CO is a reactant. Never read orders from the balanced equation.
\end{attention}

\begin{definition}[The rate constant $k$]
The \cle{rate constant} $k$ is the proportionality factor in the rate law. Numerically, it equals the rate of reaction when all concentration terms equal $\SI{1}{mol.dm^{-3}}$. Its value:
\begin{itemize}
\item is \textbf{constant} for a given reaction at a given temperature;
\item \textbf{increases with temperature} (its temperature dependence is described by the Arrhenius equation, Section 5);
\item does \textbf{not} depend on concentrations;
\item changes if a catalyst is used, because the catalyst provides a different reaction pathway.
\end{itemize}
\end{definition}

\textbf{Meaning of the orders.}\par
\begin{itemize}
\item \textbf{Zero order} ($m=0$): the rate is independent of $[\ce{A}]$; $v = k$. Typical of reactions occurring on a saturated catalyst surface.
\item \textbf{First order} ($m=1$): doubling $[\ce{A}]$ doubles the rate.
\item \textbf{Second order} ($m=2$): doubling $[\ce{A}]$ multiplies the rate by four.
\end{itemize}

\begin{propriete}[Units of the rate constant]
Since the rate always has units of $\si{mol.dm^{-3}.s^{-1}}$, rearranging $k = v/[\ce{A}]^{m}[\ce{B}]^{n}$ gives the general rule:
\[
\text{units of } k = (\si{mol.dm^{-3}})^{1-\text{overall order}}\,\si{s^{-1}}
\]
\begin{center}
\begin{tabularx}{\linewidth}{|l|X|X|}
\hline
\rowcolor{popblueL} \textbf{Overall order} & \textbf{Example rate law} & \textbf{Units of $k$} \\
\hline
0 & $v = k$ & $\si{mol.dm^{-3}.s^{-1}}$ \\
\hline
1 & $v = k[\ce{A}]$ & $\si{s^{-1}}$ \\
\hline
2 & $v = k[\ce{A}]^2$ or $k[\ce{A}][\ce{B}]$ & $\si{dm^3.mol^{-1}.s^{-1}}$ \\
\hline
3 & $v = k[\ce{A}]^2[\ce{B}]$ & $\si{dm^6.mol^{-2}.s^{-1}}$ \\
\hline
\end{tabularx}
\end{center}
A quick way to check a proposed rate law in the examination: work out the units of $k$ --- if they do not match the data, the law is wrong.
\end{propriete}

\courssub{The Initial-Rates Method}

\textbf{The idea.}\par
The cleanest way to find the orders is to measure the \emph{initial} rate of the reaction (the gradient of the concentration--time curve at $t=0$, Section 2) for several experiments in which the initial concentrations are changed \emph{one at a time}. At $t=0$ no products have accumulated, so no reverse or side reactions complicate the measurement.

\begin{methode}[Finding orders by the initial-rates (ratio) method]
\begin{enumerate}
\item Tabulate the experiments: initial concentrations $[\ce{A}]_0$, $[\ce{B}]_0$ and measured initial rate $v_0$ for each run.
\item Choose two experiments where $[\ce{B}]$ is constant and only $[\ce{A}]$ changes. Write the ratio:
\[
\frac{v_1}{v_2} = \left(\frac{[\ce{A}]_1}{[\ce{A}]_2}\right)^{m}
\]
and solve for $m$ (if the concentration doubles and the rate quadruples, $2^m = 4$, so $m=2$).
\item Repeat with a pair of experiments where $[\ce{A}]$ is constant to find $n$.
\item Write the complete rate law, then substitute the data of \emph{any one} experiment to calculate $k$, with its units.
\end{enumerate}
\end{methode}

\begin{exoresolu}[Initial rates: the reaction between NO and H$_2$]
The reaction $\ce{2NO(g) + 2H2(g) -> N2(g) + 2H2O(g)}$ was studied at $\SI{800}{\ensuremath{{}^{\circ}\mathrm{C}}}$:
\begin{center}
\begin{tabular}{|c|c|c|c|}
\hline
\rowcolor{popblueL} Expt & $[\ce{NO}]_0$ / $\si{mol.dm^{-3}}$ & $[\ce{H2}]_0$ / $\si{mol.dm^{-3}}$ & Initial rate / $\si{mol.dm^{-3}.s^{-1}}$ \\
\hline
1 & $0.10$ & $0.10$ & $1.2\times10^{-3}$ \\
\hline
2 & $0.20$ & $0.10$ & $4.8\times10^{-3}$ \\
\hline
3 & $0.20$ & $0.20$ & $9.6\times10^{-3}$ \\
\hline
\end{tabular}
\end{center}
Determine the rate law, the overall order and the value of $k$.
\tcblower
\textbf{Order in NO:} compare experiments 1 and 2 ($[\ce{H2}]$ constant). $[\ce{NO}]$ doubles ($0.10 \to 0.20$) and the rate is multiplied by $4.8/1.2 = 4$. Hence $2^m = 4$, so $m = 2$.

\textbf{Order in H$_2$:} compare experiments 2 and 3 ($[\ce{NO}]$ constant). $[\ce{H2}]$ doubles and the rate doubles ($9.6/4.8 = 2$). Hence $2^n = 2$, so $n = 1$.

\textbf{Rate law:} $v = k[\ce{NO}]^2[\ce{H2}]$; overall order $= 2 + 1 = 3$.

\textbf{Rate constant} (using experiment 1):
\[
k = \frac{v}{[\ce{NO}]^2[\ce{H2}]} = \frac{1.2\times10^{-3}}{(0.10)^2(0.10)} = \SI{1.2}{dm^6.mol^{-2}.s^{-1}}
\]
Check the units: $(\si{mol.dm^{-3}})^{1-3}\si{s^{-1}} = \si{dm^6.mol^{-2}.s^{-1}}$. $\surd$
\end{exoresolu}

\textbf{The log--log alternative.}\par
When only one reactant influences the rate, $v = k[\ce{A}]^m$ can be linearised by taking logarithms:
\[
\log(v_0) = \log k + m\log[\ce{A}]_0
\]
A plot of $\log(\text{initial rate})$ against $\log[\ce{A}]_0$ is a straight line whose \textbf{gradient is the order $m$} and whose intercept is $\log k$. This is especially useful when the order is not an obvious integer.

\begin{popfigure}
\begin{center}
\begin{tikzpicture}
\begin{axis}[
  width=0.72\linewidth, height=6cm,
  xlabel={$\log([\ce{A}]_0\ /\ \si{mol.dm^{-3}})$},
  ylabel={$\log(v_0)$},
  xmin=-2.2, xmax=-0.4, ymin=-5.5, ymax=-1.5,
  grid=both, grid style={popline!50},
  legend style={at={(0.03,0.97)}, anchor=north west, font=\small},
]
\addplot[domain=-2.1:-0.5, thick, popblue] {2*x - 1.0};
\addlegendentry{gradient $= m = 2$}
\addplot[only marks, mark=*, mark size=1.8pt, poporange] coordinates {(-2.0,-5.0) (-1.5,-4.0) (-1.0,-3.0) (-0.6,-2.2)};
\end{axis}
\end{tikzpicture}
\end{center}
\legende{Log--log plot of initial rate against initial concentration. The slope gives the order directly: here $m = 2$, and the intercept gives $\log k$.}
\end{popfigure}

\courssub{Graphical Method: Integrated Rate Laws}

\textbf{Why integrate?}\par
The initial-rates method needs several experiments. If instead we follow \emph{one} reaction mixture over time (using any method of Section 3), we can find the order from the \emph{shape} of the concentration--time data, using the integrated forms of the rate laws. For reactions with more than one reactant, we use the \cle{isolation method} (or \cle{pseudo-order method}): one reactant is in large excess so its concentration stays effectively constant, and the rate law reduces to $v = k'[\ce{A}]^m$ where $k' = k[\ce{B}]_0^n$ is a pseudo-rate constant.

\begin{propriete}[Integrated rate laws (statements and linear plots)]
For a reactant \ce{A} of concentration $[\ce{A}]_t$ at time $t$ and $[\ce{A}]_0$ at $t=0$:
\begin{itemize}
\item \textbf{Zero order} ($-d[\ce{A}]/dt = k$): $\;[\ce{A}]_t = [\ce{A}]_0 - kt$. A plot of $[\ce{A}]$ vs $t$ is linear, gradient $= -k$.
\item \textbf{First order} ($-d[\ce{A}]/dt = k[\ce{A}]$): $\;\ln[\ce{A}]_t = \ln[\ce{A}]_0 - kt$. A plot of $\ln[\ce{A}]$ vs $t$ is linear, gradient $= -k$.
\item \textbf{Second order} ($-d[\ce{A}]/dt = k[\ce{A}]^2$): $\;\dfrac{1}{[\ce{A}]_t} = \dfrac{1}{[\ce{A}]_0} + kt$. A plot of $1/[\ce{A}]$ vs $t$ is linear, gradient $= +k$.
\end{itemize}
The first-order derivation is examinable: separating variables,
\[
\int_{[\ce{A}]_0}^{[\ce{A}]_t} \frac{d[\ce{A}]}{[\ce{A}]} = -k\int_0^t dt
\quad\Longrightarrow\quad
\ln[\ce{A}]_t - \ln[\ce{A}]_0 = -kt.
\]
\end{propriete}

\begin{methode}[Identifying the order graphically]
\begin{enumerate}
\item From the experimental data, compute $[\ce{A}]$, $\ln[\ce{A}]$ and $1/[\ce{A}]$ at each time.
\item Plot the three graphs: $[\ce{A}]$ vs $t$, $\ln[\ce{A}]$ vs $t$, $1/[\ce{A}]$ vs $t$.
\item \textbf{Only one} of the three is a straight line: that plot identifies the order.
\item Read the gradient of the straight line to obtain $k$ (with the correct units).
\end{enumerate}
\end{methode}

\begin{popfigure}
\begin{center}
\begin{tikzpicture}
\begin{axis}[
  name=p1, width=0.34\linewidth, height=4.6cm,
  title={$[\ce{A}]$ vs $t$}, xlabel={$t$ / s}, ylabel={$[\ce{A}]$},
  xmin=0, xmax=100, ymin=0, ymax=1.1, tick label style={font=\scriptsize}, label style={font=\scriptsize}, title style={font=\small},
]
\addplot[domain=0:100, thick, popblue, samples=60] {exp(-0.025*x)};
\addplot[only marks, mark=*, mark size=1pt, poporange] coordinates {(0,1) (20,0.607) (40,0.368) (60,0.223) (80,0.135) (100,0.082)};
\end{axis}
\begin{axis}[
  name=p2, at={($(p1.east)+(1.1cm,0)$)}, anchor=west,
  width=0.34\linewidth, height=4.6cm,
  title={$\ln[\ce{A}]$ vs $t$}, xlabel={$t$ / s}, ylabel={$\ln[\ce{A}]$},
  xmin=0, xmax=100, ymin=-2.7, ymax=0.2, tick label style={font=\scriptsize}, label style={font=\scriptsize}, title style={font=\small},
]
\addplot[domain=0:100, thick, popgreen] {-0.025*x};
\addplot[only marks, mark=*, mark size=1pt, poporange] coordinates {(0,0) (20,-0.5) (40,-1.0) (60,-1.5) (80,-2.0) (100,-2.5)};
\end{axis}
\begin{axis}[
  name=p3, at={($(p2.east)+(1.1cm,0)$)}, anchor=west,
  width=0.34\linewidth, height=4.6cm,
  title={$1/[\ce{A}]$ vs $t$}, xlabel={$t$ / s}, ylabel={$1/[\ce{A}]$},
  xmin=0, xmax=100, ymin=0, ymax=13, tick label style={font=\scriptsize}, label style={font=\scriptsize}, title style={font=\small},
]
\addplot[domain=0:100, thick, poppurple, samples=60] {exp(0.025*x)};
\addplot[only marks, mark=*, mark size=1pt, poporange] coordinates {(0,1) (20,1.65) (40,2.72) (60,4.48) (80,7.39) (100,12.2)};
\end{axis}
\end{tikzpicture}
\end{center}
\legende{The same first-order data set tested three ways. Only the $\ln[\ce{A}]$ vs $t$ plot is linear (gradient $=-k$): the reaction is first order. The curved $[\ce{A}]$ vs $t$ and $1/[\ce{A}]$ vs $t$ plots confirm it is neither zero nor second order.}
\end{popfigure}

\begin{exoresolu}[Graphical test for the hydrolysis of an ester]
The hydrolysis of methyl ethanoate in excess dilute acid was followed by titration (Section 2). The concentration of ester remaining was:
\begin{center}
\begin{tabular}{|c|c|c|c|c|c|}
\hline
\rowcolor{popblueL} $t$ / min & 0 & 20 & 40 & 60 & 80 \\
\hline
$[\text{ester}]$ / $\si{mol.dm^{-3}}$ & 0.500 & 0.389 & 0.303 & 0.236 & 0.184 \\
\hline
$\ln[\text{ester}]$ & $-0.693$ & $-0.944$ & $-1.194$ & $-1.444$ & $-1.693$ \\
\hline
\end{tabular}
\end{center}
Show the reaction is first order in ester and find $k$.
\tcblower
The values of $\ln[\text{ester}]$ fall by almost exactly $0.250$ every 20 minutes: the plot of $\ln[\text{ester}]$ against $t$ is a straight line, so the reaction is \textbf{first order} in ester (the acid, in large excess, does not appear --- its concentration is effectively constant, giving a \emph{pseudo-first-order} rate constant $k'$).

Gradient of the line:
\[
-k' = \frac{-1.693 - (-0.693)}{80 - 0} = \frac{-1.000}{80} = -0.0125\ \si{min^{-1}}
\]
Hence $k' = \SI{0.0125}{min^{-1}} = \SI{2.08e-4}{s^{-1}}$. (First order, so the unit is $\si{s^{-1}}$.)
\end{exoresolu}

\courssub{Half-Life of a First-Order Reaction}

\begin{definition}[Half-life]
The \cle{half-life} $t_{1/2}$ of a reaction is the time taken for the concentration of a reactant to fall to half of its initial value.
\end{definition}

For a first-order reaction, set $[\ce{A}]_t = [\ce{A}]_0/2$ in the integrated law:
\[
\ln\!\left(\frac{[\ce{A}]_0}{[\ce{A}]_0/2}\right) = kt_{1/2}
\quad\Longrightarrow\quad
\ln 2 = kt_{1/2}
\quad\Longrightarrow\quad
\boxed{\,t_{1/2} = \frac{\ln 2}{k} = \frac{0.693}{k}\,}
\]

\begin{propriete}[First-order half-life is constant]
For a first-order reaction, $t_{1/2} = \ln 2 / k$ depends \textbf{only on $k$}, not on the initial concentration. The concentration therefore falls in the pattern $[\ce{A}]_0 \to [\ce{A}]_0/2 \to [\ce{A}]_0/4 \to [\ce{A}]_0/8$ after $1, 2, 3$ half-lives. This constant half-life is a quick experimental \emph{test} for first order: if the time to halve is the same wherever you start on the curve, the reaction is first order. 

For comparison:
\begin{itemize}
\item Zero order: $t_{1/2} = \dfrac{[\ce{A}]_0}{2k}$ (decreases as reaction proceeds).
\item Second order: $t_{1/2} = \dfrac{1}{k[\ce{A}]_0}$ (increases as reaction proceeds).
\end{itemize}
\end{propriete}

\begin{exemplebox}[Half-life of N$_2$O$_5$ decomposition]
For $\ce{2N2O5(g) -> 4NO2(g) + O2(g)}$, $k = \SI{6.2e-4}{s^{-1}}$ at $\SI{45}{\ensuremath{{}^{\circ}\mathrm{C}}}$. Then
\[
t_{1/2} = \frac{0.693}{6.2\times10^{-4}} \approx \SI{1.1e3}{s} \approx \SI{19}{min}.
\]
Whatever the starting pressure of \ce{N2O5}, half of it disappears every 19 minutes at this temperature.
\end{exemplebox}

\begin{attention}[Fractional and negative orders exist]
Do not force the data to give order 0, 1 or 2. Many real reactions have \textbf{fractional orders} (e.g.\ $v = k[\ce{CH3CHO}]^{3/2}$ for the thermal decomposition of ethanal) or even negative orders, because they proceed through multi-step mechanisms. If your ratio calculation gives $2^m = 2.8$, take logarithms: $m = \log 2.8 / \log 2 \approx 1.5$. Also remember: the order is defined only with respect to concentration; $k$ itself changes with temperature.
\end{attention}

\begin{exercice}[Practice: orders and half-life]
\begin{enumerate}
\item For the reaction $\ce{A + B -> C}$, the rate doubles when $[\ce{A}]$ doubles (at fixed $[\ce{B}]$) and is unchanged when $[\ce{B}]$ doubles. Write the rate law, state the overall order, and give the units of $k$.
\item A first-order reaction has $t_{1/2} = \SI{346}{s}$. Calculate $k$, and the time needed for the concentration to fall to one eighth of its initial value.
\item Using the data of the worked example on ester hydrolysis, verify that the half-life is constant by computing the time to go from 0.500 to 0.250 and from 0.400 to 0.200 $\si{mol.dm^{-3}}$.
\end{enumerate}
\end{exercice}

\begin{corrige}[Exercise 4.4]
\begin{enumerate}
\item Rate $v = k[\ce{A}]$; first order in \ce{A}, zero order in \ce{B}, overall order 1; units of $k$: $\si{s^{-1}}$.
\item $k = 0.693/346 = \SI{2.00e-3}{s^{-1}}$. One eighth means three half-lives: $t = 3 \times 346 = \SI{1038}{s}$.
\item From the integrated law, $t_{1/2} = \ln 2 / k' = 0.693 / 0.0125 = \SI{55.4}{min}$.
  \begin{itemize}
  \item From 0.500 to 0.250: $\ln(0.250) = -1.386$. Interpolating between $t=40$ ($\ln=-1.194$) and $t=60$ ($\ln=-1.444$): $t \approx 40 + \frac{1.386-1.194}{1.444-1.194}\times20 = \SI{55.4}{min}$.
  \item From 0.400 to 0.200: $\ln(0.400)=-0.916$, $\ln(0.200)=-1.609$. Interpolating: $t_{0.400}\approx 17.8$ min, $t_{0.200}\approx 73.2$ min. Difference $= \SI{55.4}{min}$.
  \end{itemize}
  Equal half-lives confirm first order. $\surd$
\end{enumerate}
\end{corrige}

\begin{aretenir}[Key points of Section 4]
\begin{itemize}
\item Rate law: $v = k[\ce{A}]^m[\ce{B}]^n$; orders $m$, $n$ come from \textbf{experiment only}, never from stoichiometric coefficients.
\item Units of $k$: $(\si{mol.dm^{-3}})^{1-\text{overall order}}\si{s^{-1}}$.
\item Initial-rates method: vary one concentration at a time and use $v_1/v_2 = ([\ce{A}]_1/[\ce{A}]_2)^m$; a log--log plot gives the order as its gradient.
\item Graphical method: the linear plot among $[\ce{A}]$, $\ln[\ce{A}]$, $1/[\ce{A}]$ vs $t$ identifies zero, first or second order; the gradient gives $k$. For multi-reactant systems, use the isolation (pseudo-order) method.
\item First-order half-life: $t_{1/2} = \ln 2 / k$, independent of initial concentration --- a diagnostic test for first order.
\item Orders may be fractional or zero; let the data speak.
\end{itemize}
\end{aretenir}

With the rate law firmly in hand, one deep question remains: \emph{why} does the rate constant depend so strongly on temperature, and what actually happens when two molecules meet? That is the subject of the next section: collision theory, the Maxwell--Boltzmann distribution, the Arrhenius equation and transition-state theory.

\coursec{Theoretical Models: Collision Theory, Arrhenius Equation, and Transition-State Theory}

In the previous section we learned how to express the rate of a reaction as a rate law, and how to determine the order of reaction and the rate constant $k$ from experimental data. We now ask a deeper question: \emph{why} does the rate constant have the value it does, and \emph{how} does it change with temperature? The answer lies in the microscopic picture of reacting particles. Three complementary models --- Collision Theory, the Arrhenius Equation, and Transition-State Theory --- provide that picture. They are the theoretical backbone of chemical kinetics and are examined in detail in the GCE Advanced Level syllabus (Lessons 98--99).

\courssub{Collision Theory and Effective Collisions}

Collision theory starts from a simple, intuitive idea: for a reaction to occur, reactant particles must \textbf{collide}. However, not every collision leads to products. Two conditions must be satisfied simultaneously:

\begin{enumerate}
  \item \textbf{Energy requirement:} the colliding particles must together possess kinetic energy at least equal to the \cle{activation energy} $E_{\mathrm{a}}$. This is the minimum energy needed to start breaking the existing bonds so that new ones can form.
  \item \textbf{Orientation requirement:} the colliding particles must be oriented in a way that allows the relevant atoms to interact. For example, in the reaction $\ce{NO + O3 -> NO2 + O2}$, the nitrogen atom of \ce{NO} must strike a terminal oxygen atom of \ce{O3}; a collision in which the oxygen end of \ce{NO} hits \ce{O3} is ineffective.
\end{enumerate}

Only collisions fulfilling \emph{both} criteria are called \cle{effective collisions}. The rate of reaction is therefore proportional to the frequency of effective collisions.

\begin{definition}[Activation Energy $E_{\mathrm{a}}$]
The minimum energy that reacting particles must possess for a collision between them to result in a chemical reaction. It is the energy barrier separating reactants from products.
\end{definition}

The fraction of molecules that have energy greater than or equal to $E_{\mathrm{a}}$ at a given temperature $T$ is obtained from the \cle{Maxwell--Boltzmann distribution} of molecular energies. As temperature rises, the distribution broadens and its maximum shifts to higher energies, while the area under the curve (the total number of molecules) stays constant. The fraction of molecules with energy at least $E_{\mathrm{a}}$ therefore increases sharply with temperature.

\begin{popfigure}
\begin{tikzpicture}
\begin{axis}[
  width=0.9\linewidth,
  height=7cm,
  domain=0:6,
  samples=200,
  xlabel={Molecular energy $E$},
  ylabel={Fraction of molecules},
  xmin=0, xmax=6,
  ymin=0, ymax=0.55,
  legend pos=north east,
  axis lines=left,
  tick style={draw=none},
  xtick=\empty, ytick=\empty,
  label style={font=\small},
  declare function={mb(\x,\t)=(2/sqrt(max(0,pi)))*\t^(-1.5)*sqrt(max(0,\x))*exp(-\x/\t);},
]
% Lower temperature curve
\addplot[popblue, thick] {mb(x,1.2)};
% Higher temperature curve
\addplot[poporange, thick] {mb(x,2.2)};
% Activation energy line
\addplot[popdark, dashed, thick] coordinates {(3,0) (3,0.52)};
% Shaded area for lower T
\addplot[popblueL, fill opacity=0.4] plot[domain=3:6] (\x,{mb(\x,1.2)}) -- (6,0) -- (3,0) -- cycle;
% Shaded area for higher T
\addplot[poporange, fill opacity=0.25] plot[domain=3:6] (\x,{mb(\x,2.2)}) -- (6,0) -- (3,0) -- cycle;
\legend{$T_1$ (lower), $T_2$ (higher), $E_{\mathrm{a}}$}
\node[popblue, font=\small, align=center] at (axis cs:4.6,0.06) {small area at $T_1$};
\node[poporange, font=\small, align=center] at (axis cs:4.6,0.16) {larger area at $T_2$};
\node[popdark, font=\small, rotate=90, anchor=south] at (axis cs:3,0.30) {$E_{\mathrm{a}}$};
\end{axis}
\end{tikzpicture}
\legende{Maxwell--Boltzmann energy distributions at two temperatures. The shaded area to the right of $E_{\mathrm{a}}$ represents the fraction of molecules with sufficient energy to react. Raising the temperature from $T_1$ to $T_2$ increases this fraction considerably, which is why rates increase so rapidly with temperature.}
\end{popfigure}

For a bimolecular gas-phase reaction between molecules A and B, the collision frequency $Z$ (number of collisions per unit volume per unit time) is proportional to the product of the concentrations $[\ce{A}][\ce{B}]$ and to $\sqrt{T}$. Only a fraction $f = e^{-E_{\mathrm{a}}/RT}$ of these collisions have enough energy, and only a fraction $p$ (the \cle{steric} or \cle{orientation factor}, usually $0 < p \le 1$) have the correct orientation. The rate constant predicted by simple collision theory is therefore
\[
k = p\, Z_0\, e^{-E_{\mathrm{a}}/RT}
\]
where $Z_0$ is the collision frequency at unit concentrations. This already resembles the Arrhenius form, with the pre-exponential factor $A \approx pZ_0$ (with a weak additional dependence on $\sqrt{T}$).

\begin{attention}[The Steric Factor]
For many reactions $p < 1$ because only a fraction of collisions have the correct geometry. Do not state, however, that $p$ is \emph{always} less than 1: for some reactions (for example ``harpoon'' reactions between alkali metals and halogens, where an electron is transferred at long range) the effective collision cross-section is larger than the geometric one and $p$ can exceed 1.
\end{attention}

\courssub{The Arrhenius Equation}

Svante Arrhenius (1889) proposed an empirical relationship that fits the temperature dependence of the vast majority of rate constants:
\[
k = A\, e^{-E_{\mathrm{a}}/RT}
\]
where
\begin{itemize}
  \item $k$ = rate constant (units depend on the overall order),
  \item $A$ = \cle{pre-exponential factor} (or frequency factor), with the same units as $k$,
  \item $E_{\mathrm{a}}$ = activation energy ($\mathrm{J\,mol^{-1}}$ or $\mathrm{kJ\,mol^{-1}}$),
  \item $R$ = gas constant ($8.314\ \mathrm{J\,mol^{-1}\,K^{-1}}$),
  \item $T$ = absolute temperature ($\mathrm{K}$).
\end{itemize}

Taking natural logarithms gives the linear form:
\[
\ln k = \ln A - \frac{E_{\mathrm{a}}}{R}\,\frac{1}{T}
\]
A plot of $\ln k$ against $1/T$ (an \cle{Arrhenius plot}) is a straight line of slope $-E_{\mathrm{a}}/R$ and intercept $\ln A$. This is the standard experimental method for measuring activation energies.

\begin{propriete}[Meaning of the Pre-exponential Factor $A$]
$A$ is related to the frequency of collisions and to the probability that a collision has the correct orientation. In collision theory $A \approx pZ_0$; in transition-state theory $A = (k_{\mathrm{B}}T/h)\, e^{\Delta S^{\ddagger}/R}$. $A$ depends weakly on temperature, but over modest temperature ranges it is treated as constant. (The identification with the Eyring parameters is examinable as a comparison, not as a formal proof.)
\end{propriete}

\begin{methode}[Determining $E_{\mathrm{a}}$ and $A$ from an Arrhenius Plot]
\begin{enumerate}
  \item Measure the rate constant $k$ at several different temperatures (at least three, preferably five or six).
  \item Convert each temperature to kelvin and compute $1/T$ ($\mathrm{K^{-1}}$) and $\ln k$.
  \item Plot $\ln k$ (vertical axis) against $1/T$ (horizontal axis).
  \item Draw the best-fit straight line and determine its slope $m$ and intercept $c$.
  \item Calculate $E_{\mathrm{a}} = -mR$ (in $\mathrm{J\,mol^{-1}}$, then divide by $1000$ for $\mathrm{kJ\,mol^{-1}}$) and $A = e^{c}$ (same units as $k$).
\end{enumerate}
\end{methode}

\begin{popfigure}
\begin{tikzpicture}
\begin{axis}[
  width=0.9\linewidth,
  height=7cm,
  xlabel={$1/T\ (\mathrm{K^{-1}})$},
  ylabel={$\ln k$},
  xmin=0.0028, xmax=0.0034,
  ymin=-5, ymax=0,
  grid=major,
  legend pos=south west,
  tick label style={font=\small},
  label style={font=\small},
]
% Simulated data points
\addplot[only marks, mark=*, popblue] coordinates {
  (0.00333, -4.2)
  (0.00323, -3.5)
  (0.00313, -2.8)
  (0.00303, -2.1)
  (0.00294, -1.5)
  (0.00286, -0.9)
};
% Best-fit line: slope -1.2e4, intercept 35.8
\addplot[poporange, thick, domain=0.0028:0.0034] {-12000*x + 35.8};
\legend{Experimental points, Best-fit line}
\node[poporange, font=\small, anchor=south west, align=left] at (axis cs:0.00305,-1.2) {slope $m = -E_{\mathrm{a}}/R$};
\node[popdark, font=\small, anchor=west, align=left] at (axis cs:0.00281,-4.6) {intercept $c = \ln A$};
\end{axis}
\end{tikzpicture}
\legende{Typical Arrhenius plot. The straight line confirms Arrhenius behaviour. From the slope $m = -1.20\times10^{4}\ \mathrm{K}$ we obtain $E_{\mathrm{a}} = -mR = (1.20\times10^{4})(8.314) = 99.8\ \mathrm{kJ\,mol^{-1}}$.}
\end{popfigure}

\begin{attention}[Apparent Activation Energy]
For complex (multi-step) reactions the observed rate constant may be a combination of several elementary rate constants. The $E_{\mathrm{a}}$ obtained from an Arrhenius plot is then an \emph{apparent} activation energy, which can vary with temperature if the rate-determining step changes. Always check the linearity of the Arrhenius plot; curvature signals a change of mechanism.
\end{attention}

\begin{exoresolu}[Determining $E_{\mathrm{a}}$ from Two Temperatures]
The rate constant for the decomposition of \ce{N2O5} in \ce{CCl4} is $k_1 = 3.4\times10^{-5}\ \mathrm{s^{-1}}$ at $T_1 = 298\ \mathrm{K}$ and $k_2 = 1.2\times10^{-3}\ \mathrm{s^{-1}}$ at $T_2 = 318\ \mathrm{K}$. Assuming Arrhenius behaviour, calculate the activation energy $E_{\mathrm{a}}$.
\tcblower
\textbf{Solution:} Dividing the Arrhenius equation at $T_2$ by that at $T_1$ eliminates $A$:
\[
\ln\frac{k_2}{k_1} = -\frac{E_{\mathrm{a}}}{R}\left(\frac{1}{T_2} - \frac{1}{T_1}\right)
\]
\[
\ln\frac{1.2\times10^{-3}}{3.4\times10^{-5}} = \ln(35.3) = 3.56
\]
\[
\frac{1}{T_2} - \frac{1}{T_1} = \frac{1}{318} - \frac{1}{298} = (3.145 - 3.356)\times10^{-3} = -2.11\times10^{-4}\ \mathrm{K^{-1}}
\]
\[
E_{\mathrm{a}} = -\frac{R \ln(k_2/k_1)}{1/T_2 - 1/T_1}
= -\frac{8.314 \times 3.56}{-2.11\times10^{-4}}
= 1.40\times10^{5}\ \mathrm{J\,mol^{-1}} = \mathbf{140\ \mathrm{kJ\,mol^{-1}}}
\]
\end{exoresolu}

\courssub{Transition-State Theory (Activated Complex Theory)}

Transition-State Theory (TST), developed by Eyring, Evans and Polanyi in the 1930s, gives a more detailed microscopic picture. It postulates that reactants pass through a high-energy, transient arrangement of atoms called the \cle{activated complex} or \cle{transition state} (denoted $\ddagger$), which is assumed to be in quasi-equilibrium with the reactants:
\[
\ensuremath{\mathrm{A + B \rightleftharpoons  AB^{\ddagger} \rightarrow  Products}}
\]
The rate of reaction is the rate at which activated complexes pass over the energy barrier to products. Assuming a universal frequency of passage $\nu = k_{\mathrm{B}}T/h$ (where $k_{\mathrm{B}}$ is Boltzmann's constant and $h$ is Planck's constant), one obtains the \cle{Eyring equation}:
\[
k = \frac{k_{\mathrm{B}}T}{h}\, e^{-\Delta G^{\ddagger}/RT}
= \frac{k_{\mathrm{B}}T}{h}\, e^{\Delta S^{\ddagger}/R}\, e^{-\Delta H^{\ddagger}/RT}
\]
where $\Delta G^{\ddagger} = \Delta H^{\ddagger} - T\Delta S^{\ddagger}$ is the Gibbs free energy of activation, $\Delta H^{\ddagger}$ the enthalpy of activation and $\Delta S^{\ddagger}$ the entropy of activation.

\begin{definition}[Activated Complex ($\ddagger$)]
A transient, high-energy arrangement of atoms corresponding to the maximum of the potential energy profile along the reaction coordinate. It has a very short lifetime (of the order of $10^{-13}\ \mathrm{s}$) and must not be confused with a stable reaction intermediate, which sits in an energy \emph{minimum} and can in principle be detected.
\end{definition}

\begin{popfigure}
\begin{tikzpicture}
\begin{axis}[
  width=0.9\linewidth,
  height=7cm,
  axis lines=middle,
  xlabel={Reaction coordinate},
  ylabel={Potential energy},
  xtick=\empty, ytick=\empty,
  xmin=-0.8, xmax=5.8,
  ymin=-0.5, ymax=5.4,
  clip=false,
  label style={font=\small},
]
% Reactants plateau
\addplot[popblue, thick, domain=-0.5:1] {0.5};
\node[popblue, font=\small, anchor=north] at (axis cs:0.25,0.4) {Reactants};
% Barrier
\addplot[poporange, thick, domain=1:3, samples=100] {0.5 + 4*sin(deg((x-1)*pi/2))^2};
\node[poporange, font=\small, anchor=south] at (axis cs:2,4.55) {$\ddagger$ (activated complex)};
% Products plateau
\addplot[popgreen, thick, domain=3:5.5] {1.5};
\node[popgreen, font=\small, anchor=north] at (axis cs:4.25,1.4) {Products};
% Ea forward arrow
\draw[popdark, <->, thick] (axis cs:1.15,0.5) -- (axis cs:1.15,4.5) node[midway, right, font=\small] {$E_{\mathrm{a}}$ (forward)};
% Ea reverse arrow
\draw[popmuted, <->, thick] (axis cs:2.85,1.5) -- (axis cs:2.85,4.5) node[midway, left, font=\small] {$E_{\mathrm{a}}$ (reverse)};
% Delta H reaction
\draw[poppurple, <->, thick] (axis cs:-0.3,0.5) -- (axis cs:-0.3,1.5) node[midway, left, font=\small] {$\Delta H_{\mathrm{rxn}}$};
\end{axis}
\end{tikzpicture}
\legende{Potential energy profile for an exothermic reaction. The activated complex ($\ddagger$) sits at the energy maximum. For an exothermic reaction $E_{\mathrm{a}}(\text{forward}) < E_{\mathrm{a}}(\text{reverse})$, and $\Delta H_{\mathrm{rxn}} = E_{\mathrm{a}}(\text{forward}) - E_{\mathrm{a}}(\text{reverse}) < 0$.}
\end{popfigure}

Comparing the Eyring equation with the Arrhenius equation,
\[
k = \underbrace{\frac{k_{\mathrm{B}}T}{h}\, e^{\Delta S^{\ddagger}/R}}_{A}\; e^{-\Delta H^{\ddagger}/RT}
\]
we identify, for a gas-phase bimolecular reaction,
\[
E_{\mathrm{a}} = \Delta H^{\ddagger} + RT \qquad\text{and}\qquad A = \frac{k_{\mathrm{B}}T}{h}\, e^{\Delta S^{\ddagger}/R}
\]
The entropy of activation $\Delta S^{\ddagger}$ measures the change in order on forming the activated complex: bringing two separate molecules together into one complex usually gives $\Delta S^{\ddagger} < 0$, which lowers $A$ and slows the reaction --- a quantitative explanation of the steric factor $p$ of collision theory.

\courssub{Comparison of Collision Theory and Transition-State Theory}

Both theories explain the temperature dependence of reaction rates, but they differ in scope and detail.

\begin{center}
\begin{tabularx}{\linewidth}{|l|X|X|}
\hline
\rowcolor{popblueL}
\textbf{Aspect} & \textbf{Collision Theory} & \textbf{Transition-State Theory} \\ \hline
\textbf{Basic idea} & Reactants must collide with sufficient energy and correct orientation. & Reactants form an activated complex in quasi-equilibrium; passage of $\ddagger$ over the barrier gives products. \\ \hline
\textbf{Key parameters} & $E_{\mathrm{a}}$, steric factor $p$, collision frequency $Z$. & $\Delta H^{\ddagger}$, $\Delta S^{\ddagger}$, $\Delta G^{\ddagger}$. \\ \hline
\textbf{Rate constant} & $k = p Z_0 e^{-E_{\mathrm{a}}/RT}$ (approx.) & $k = \dfrac{k_{\mathrm{B}}T}{h}\, e^{-\Delta G^{\ddagger}/RT}$ \\ \hline
\textbf{Entropy effects} & Not explicit; hidden in the empirical factor $p$. & Explicit via $\Delta S^{\ddagger}$; explains why some reactions with low $E_{\mathrm{a}}$ are still slow. \\ \hline
\textbf{Applicability} & Mainly gas-phase bimolecular reactions. & Gas phase, solution and surfaces; unimolecular and bimolecular. \\ \hline
\textbf{Link with Arrhenius} & $A \approx p Z_0$ (weak $T$ dependence). & $A = \dfrac{k_{\mathrm{B}}T}{h}\, e^{\Delta S^{\ddagger}/R}$; $E_{\mathrm{a}} = \Delta H^{\ddagger} + RT$. \\ \hline
\end{tabularx}
\end{center}

\begin{aretenir}[Key Points for the Examination]
\begin{itemize}
  \item Collision theory: an effective collision requires both $E \ge E_{\mathrm{a}}$ and correct orientation. The Maxwell--Boltzmann distribution shows how the fraction of molecules with $E \ge E_{\mathrm{a}}$ increases with temperature.
  \item Arrhenius equation: $k = A e^{-E_{\mathrm{a}}/RT}$; linear form $\ln k = \ln A - E_{\mathrm{a}}/(RT)$. Plot $\ln k$ against $1/T$: slope $= -E_{\mathrm{a}}/R$, intercept $= \ln A$.
  \item Two-temperature form: $\ln(k_2/k_1) = -\dfrac{E_{\mathrm{a}}}{R}\left(\dfrac{1}{T_2}-\dfrac{1}{T_1}\right)$.
  \item Transition-state theory: activated complex $\ddagger$ in quasi-equilibrium with reactants; Eyring equation $k = (k_{\mathrm{B}}T/h)\, e^{-\Delta G^{\ddagger}/RT}$.
  \item $\Delta G^{\ddagger} = \Delta H^{\ddagger} - T\Delta S^{\ddagger}$; for gas-phase reactions $E_{\mathrm{a}} = \Delta H^{\ddagger} + RT$.
  \item TST includes the entropy of activation and so explains steric effects quantitatively; collision theory lumps them into the empirical factor $p$.
  \item The apparent $E_{\mathrm{a}}$ from an Arrhenius plot may not be constant for complex mechanisms.
\end{itemize}
\end{aretenir}

\begin{exoresolu}[From Arrhenius to Eyring Parameters]
For a certain gas-phase reaction the Arrhenius parameters are $A = 2.5\times10^{13}\ \mathrm{dm^3\,mol^{-1}\,s^{-1}}$ and $E_{\mathrm{a}} = 75.0\ \mathrm{kJ\,mol^{-1}}$ at $T = 298\ \mathrm{K}$. Calculate $\Delta H^{\ddagger}$, $\Delta S^{\ddagger}$ and $\Delta G^{\ddagger}$. (Take $k_{\mathrm{B}} = 1.381\times10^{-23}\ \mathrm{J\,K^{-1}}$, $h = 6.626\times10^{-34}\ \mathrm{J\,s}$.)
\tcblower
\textbf{Solution:}
\begin{align*}
\Delta H^{\ddagger} &= E_{\mathrm{a}} - RT \\
&= 75.0\ \mathrm{kJ\,mol^{-1}} - (8.314\times10^{-3}\ \mathrm{kJ\,mol^{-1}\,K^{-1}} \times 298\ \mathrm{K}) \\
&= 75.0 - 2.48 = \mathbf{72.5\ \mathrm{kJ\,mol^{-1}}} \\[4pt]
A &= \frac{k_{\mathrm{B}}T}{h}\, e^{\Delta S^{\ddagger}/R}
\quad\Rightarrow\quad
\Delta S^{\ddagger} = R \ln\!\left(\frac{A h}{k_{\mathrm{B}}T}\right) \\[2pt]
\frac{A h}{k_{\mathrm{B}}T} &= \frac{(2.5\times10^{13})(6.626\times10^{-34})}
{(1.381\times10^{-23})(298)}
= \frac{1.657\times10^{-20}}{4.115\times10^{-21}} = 4.03 \\[2pt]
\Delta S^{\ddagger} &= 8.314 \times \ln(4.03) = 8.314 \times 1.39 = \mathbf{+11.6\ \mathrm{J\,mol^{-1}\,K^{-1}}} \\[4pt]
\Delta G^{\ddagger} &= \Delta H^{\ddagger} - T\Delta S^{\ddagger} \\
&= 72.5\ \mathrm{kJ\,mol^{-1}} - (298\ \mathrm{K})(11.6\times10^{-3}\ \mathrm{kJ\,mol^{-1}\,K^{-1}}) \\
&= 72.5 - 3.46 = \mathbf{69.0\ \mathrm{kJ\,mol^{-1}}}
\end{align*}
\emph{Remark:} strictly, $A$ and $k_{\mathrm{B}}T/h$ have different units for a bimolecular reaction, so the ratio $Ah/(k_{\mathrm{B}}T)$ implicitly uses the standard-state concentration ($1\ \mathrm{mol\,dm^{-3}}$); this convention is assumed in the GCE treatment.
\end{exoresolu}

\begin{exercice}[Practice: Arrhenius Plot and TST]
The rate constant for the gas-phase reaction $\ce{2HI(g) -> H2(g) + I2(g)}$ was measured at three temperatures:
\begin{center}
\begin{tabular}{cc}
\hline
$T/\mathrm{K}$ & $k/\mathrm{dm^3\,mol^{-1}\,s^{-1}}$ \\ \hline
500 & $1.2\times10^{-3}$ \\
550 & $4.8\times10^{-3}$ \\
600 & $1.6\times10^{-2}$ \\ \hline
\end{tabular}
\end{center}
\begin{enumerate}
  \item Construct an Arrhenius plot and determine $E_{\mathrm{a}}$ and $A$.
  \item Using the values obtained, calculate $\Delta H^{\ddagger}$, $\Delta S^{\ddagger}$ and $\Delta G^{\ddagger}$ at $550\ \mathrm{K}$.
  \item Comment on the sign of $\Delta S^{\ddagger}$ in relation to the molecularity of the reaction.
\end{enumerate}
\end{exercice}

\begin{corrige}[Exercise 5.1]
\begin{enumerate}
  \item Compute $1/T$ and $\ln k$:
  \begin{center}
  \begin{tabular}{cc}
  \hline
  $1/T\ (\mathrm{K^{-1}})$ & $\ln k$ \\ \hline
  $2.000\times10^{-3}$ & $-6.73$ \\
  $1.818\times10^{-3}$ & $-5.34$ \\
  $1.667\times10^{-3}$ & $-4.14$ \\ \hline
  \end{tabular}
  \end{center}
  The points are collinear. Using the extreme points:
  \[
  m = \frac{-4.14-(-6.73)}{1.667\times10^{-3}-2.000\times10^{-3}}
  = \frac{2.59}{-3.33\times10^{-4}} = -7.77\times10^{3}\ \mathrm{K}
  \]
  \[
  E_{\mathrm{a}} = -mR = (7.77\times10^{3})(8.314) = 6.46\times10^{4}\ \mathrm{J\,mol^{-1}} = \mathbf{64.6\ \mathrm{kJ\,mol^{-1}}}
  \]
  Intercept: $c = \ln k - m(1/T) = -6.73 + (7.77\times10^{3})(2.000\times10^{-3}) = 8.81$, so
  \[
  A = e^{8.81} = \mathbf{6.7\times10^{3}\ \mathrm{dm^3\,mol^{-1}\,s^{-1}}}
  \]
  \item At $T = 550\ \mathrm{K}$:
  \[
  \Delta H^{\ddagger} = E_{\mathrm{a}} - RT = 64.6 - (8.314\times10^{-3}\times550) = 64.6 - 4.57 = \mathbf{60.0\ \mathrm{kJ\,mol^{-1}}}
  \]
  \[
  \Delta S^{\ddagger} = R\ln\!\left(\frac{Ah}{k_{\mathrm{B}}T}\right)
  = 8.314\ln\!\left(\frac{(6.7\times10^{3})(6.626\times10^{-34})}{(1.381\times10^{-23})(550)}\right)
  \]
  \[
  = 8.314\ln(5.8\times10^{-10}) = 8.314\times(-21.3) = \mathbf{-177\ \mathrm{J\,mol^{-1}\,K^{-1}}}
  \]
  \[
  \Delta G^{\ddagger} = \Delta H^{\ddagger} - T\Delta S^{\ddagger}
  = 60.0 - (550)(-0.177) = 60.0 + 97.2 = \mathbf{157\ \mathrm{kJ\,mol^{-1}}}
  \]
  \item The reaction is bimolecular: two \ce{HI} molecules must come together to form a single activated complex. This association reduces the freedom of motion (translational degrees of freedom) of the particles, so the system becomes more ordered and $\Delta S^{\ddagger}$ is expected to be \emph{negative} --- exactly what is found. The large negative $\Delta S^{\ddagger}$ makes $\Delta G^{\ddagger}$ much larger than $\Delta H^{\ddagger}$ and corresponds to a small steric factor $p$ in collision theory: only a small fraction of sufficiently energetic collisions have the correct orientation.
\end{enumerate}
\end{corrige}

\begin{savaistu}[Did You Know? --- The Arrhenius Equation in Everyday Life]
The Arrhenius equation explains why food spoils faster at room temperature than in a refrigerator. Taking a typical activation energy of about $80\ \mathrm{kJ\,mol^{-1}}$ for the chemical and microbial processes of spoilage, lowering the temperature from $25^{\circ}\mathrm{C}$ ($298\ \mathrm{K}$) to $4^{\circ}\mathrm{C}$ ($277\ \mathrm{K}$) changes the rate constant by the factor
\[
\frac{k_{277}}{k_{298}} = \exp\!\left[-\frac{E_{\mathrm{a}}}{R}\left(\frac{1}{277}-\frac{1}{298}\right)\right]
= \exp(-2.45) \approx 0.086
\]
i.e. spoilage becomes roughly \emph{12 times slower}. This is why cold storage is so valuable in Cameroon's tropical climate for preserving tomatoes, plantains, fish and meat between the market and the kitchen.
\end{savaistu}

\begin{attention}[Common Examination Traps]
\begin{itemize}
  \item Forgetting to convert $E_{\mathrm{a}}$ to $\mathrm{J\,mol^{-1}}$ when using $R = 8.314\ \mathrm{J\,mol^{-1}\,K^{-1}}$.
  \item Using $\log_{10}$ instead of $\ln$ in the Arrhenius plot; with $\log_{10}$ the slope is $-E_{\mathrm{a}}/(2.303R)$.
  \item Confusing $E_{\mathrm{a}}$ (Arrhenius) with $\Delta H^{\ddagger}$ (Eyring); remember $E_{\mathrm{a}} = \Delta H^{\ddagger} + RT$ for gas-phase reactions.
  \item Confusing the \emph{activated complex} (an energy maximum, undetectable) with a reaction \emph{intermediate} (an energy minimum, detectable).
  \item Stating that the steric factor $p$ is always less than 1; long-range attractions can make $p > 1$.
  \item Claiming that a catalyst lowers $E_{\mathrm{a}}$ ``by giving molecules more energy'': a catalyst provides an \emph{alternative pathway} of lower activation energy; the temperature, and hence the molecular energies, are unchanged.
\end{itemize}
\end{attention}

\coursec{Catalysis: Homogeneous and Heterogeneous}

In the previous section we saw that the rate of a reaction is controlled by its \cle{activation energy} $E_a$: the Arrhenius equation $k = A e^{-E_a/RT}$ tells us that even a modest decrease in $E_a$ produces an enormous increase in the rate constant $k$, because $E_a$ appears in an exponential. This raises a practical question of immense industrial importance: \emph{can we lower the activation energy of a reaction without changing the reactants or the temperature?} The answer is yes — this is precisely what a \cle{catalyst} does. From the manufacture of ammonia that feeds fertiliser factories supplying Cameroonian agriculture, to the catalytic converters fitted in the exhaust pipes of taxis in Douala, catalysis is one of the most economically important applications of chemical kinetics.

\courssub{What is a Catalyst?}

\begin{definition}[Catalyst]
A \cle{catalyst} is a substance that \textbf{increases the rate of a chemical reaction} without itself being \textbf{permanently consumed} in the reaction. It is recovered, chemically unchanged, at the end of the reaction, although it may take part in intermediate steps.
\end{definition}

The key to understanding catalysis lies in the energy profile. A catalyst does \emph{not} simply ``push'' the reactants harder; it provides an \textbf{alternative reaction pathway} (a different mechanism, usually with several steps) whose highest energy barrier is \textbf{lower} than the activation energy $E_a$ of the uncatalysed pathway. Since a much larger fraction of molecules now possesses enough energy to cross the barrier (recall the Maxwell--Boltzmann distribution), the rate increases dramatically.

\begin{popfigure}
\begin{center}
\begin{tikzpicture}[scale=1.0]
\draw[->, thick, popmuted] (0,0) -- (10,0) node[below left, font=\small] {Reaction coordinate};
\draw[->, thick, popmuted] (0,0) -- (0,5.2) node[above, font=\small] {Energy};
% reactants and products levels
\draw[thick, popdark] (0.4,2.2) -- (1.8,2.2) node[right, font=\small, popdark] {Reactants};
\draw[thick, popdark] (8.2,1.0) -- (9.6,1.0) node[left, font=\small, popdark] {Products};
% uncatalysed curve
\draw[very thick, poporange] (1.8,2.2) .. controls (3.6,5.0) and (5.0,5.0) .. (8.2,1.0);
% catalysed curve (two humps)
\draw[very thick, popgreen] (1.8,2.2) .. controls (3.0,3.4) and (3.6,3.4) .. (4.4,2.6)
   .. controls (5.2,1.9) and (5.8,3.2) .. (8.2,1.0);
% Ea arrows
\draw[<->, thick, poporange] (2.6,2.2) -- (2.6,4.55) node[midway, left, font=\small, poporange] {$E_a$ (uncat.)};
\draw[<->, thick, popgreen] (3.55,2.2) -- (3.55,3.25) node[midway, right, font=\small, popgreen] {$E_a$ (cat.)};
% Delta H
\draw[<->, thick, popblue] (9.0,1.0) -- (9.0,2.2) node[midway, right, font=\small, popblue] {$\Delta H$};
\end{tikzpicture}
\end{center}
\legende{Energy profile for an exothermic reaction. The catalysed pathway (green) proceeds through an intermediate (the dip between the two humps) with a lower activation energy than the uncatalysed pathway (orange). Note that $\Delta H$ is unchanged.}
\end{popfigure}

\begin{propriete}[What a catalyst does and does not do]
\begin{itemize}
\item A catalyst \textbf{lowers the activation energy} by providing an alternative mechanism, so the rate constant increases: $k = A e^{-E_a/RT}$ grows as $E_a$ falls.
\item A catalyst \textbf{does not change} the enthalpy change $\Delta H$, the free energy change $\Delta G$, or the \textbf{equilibrium constant} $K$ of the reaction.
\item Because it lowers the barrier in \emph{both} directions, a catalyst accelerates the \textbf{forward and reverse reactions equally}. It therefore helps the system to \emph{reach equilibrium faster}, but the position of equilibrium is unchanged.
\item Only a very small amount of catalyst is needed, since it is regenerated during the catalytic cycle.
\end{itemize}
(The thermodynamic statement — that catalysis cannot alter $K$ because $K$ depends only on $\Delta G^{\circ}$ — is examinable as a short-answer justification.)
\end{propriete}

\begin{attention}[A catalyst never makes an impossible reaction happen]
A common examination error is to write that a catalyst ``makes a reaction occur''. If a reaction is thermodynamically impossible ($\Delta G > 0$), no catalyst can make it go. A catalyst only speeds up a reaction that is already feasible but kinetically slow. Likewise, a catalyst does not appear in the overall balanced equation — if you find it as a reactant on the left and never regenerated, it is not a catalyst but a reagent. Finally, do not confuse an \textbf{intermediate} (a real, sometimes isolable species sitting in an energy \emph{valley} between two humps) with the \textbf{transition state} (the unstable arrangement at the \emph{top} of a hump).
\end{attention}

Catalysts are classified into two great families according to the \emph{phase} they occupy relative to the reactants.

\begin{definition}[Homogeneous and heterogeneous catalysts]
A \cle{homogeneous catalyst} is in the \textbf{same phase} as the reactants (typically all in aqueous solution or all gases). A \cle{heterogeneous catalyst} is in a \textbf{different phase} from the reactants — most often a \emph{solid} catalyst in contact with gaseous or liquid reactants.
\end{definition}

\courssub{Homogeneous Catalysis}

In homogeneous catalysis, the catalyst mixes freely with the reactants at the molecular level, so every catalyst molecule is accessible. The mechanism proceeds through the formation of a reactive \cle{intermediate}: the catalyst is consumed in one step and regenerated in a later one.

\textbf{Example 1: acid-catalysed ester hydrolysis.} The hydrolysis of an ester such as ethyl ethanoate is slow with pure water but rapid in the presence of dilute \ce{HCl} or \ce{H2SO4}, which supply \ce{H+} ions:

\[
\ce{CH3COOC2H5(aq) + H2O(l) ->[H+] CH3COOH(aq) + C2H5OH(aq)}
\]

The \ce{H+} ion first protonates the carbonyl oxygen of the ester, making the carbonyl carbon much more susceptible to nucleophilic attack by water. After a series of proton transfers, the products are released and \ce{H+} is regenerated. Each step of this catalysed route has a lower activation energy than the direct attack of water on the neutral ester.

\textbf{Example 2: the persulfate--iodide reaction.} The oxidation of iodide by peroxodisulfate,

\[
\ce{S2O8^{2-}(aq) + 2I^{-}(aq) -> 2SO4^{2-}(aq) + I2(aq)},
\]

is slow because it brings two negative ions together (strong electrostatic repulsion). It is catalysed by \ce{Fe^{2+}} or \ce{Fe^{3+}} ions, which provide a two-step route in which each step pairs ions of \emph{opposite} charge:

\begin{align*}
\ce{S2O8^{2-} + 2Fe^{2+} &-> 2SO4^{2-} + 2Fe^{3+}}\\
\ce{2Fe^{3+} + 2I^{-} &-> 2Fe^{2+} + I2}
\end{align*}

The iron ions are consumed in the first step and regenerated in the second — the signature of a true catalyst. Adding the two steps regenerates the overall equation, with no iron species appearing in it.

\textbf{Example 3: enzyme catalysis.} Enzymes are biological catalysts of extraordinary efficiency and specificity. Amylase in saliva hydrolyses starch; catalase decomposes hydrogen peroxide into water and oxygen almost instantaneously. Enzymes work by binding the reactant (the \emph{substrate}) at a specifically shaped \emph{active site}, holding it in the ideal geometry for reaction — a beautiful molecular illustration of transition-state stabilisation.

\begin{savaistu}[Catalysis and life]
A single molecule of catalase can decompose millions of hydrogen peroxide molecules per second. Without enzymes, the biochemical reactions that keep us alive would be far too slow at body temperature — life itself is, in a sense, a triumph of catalysis. The fermentation of palm wine and of cassava dough for \emph{garri}, familiar processes in Cameroon, are also driven by enzyme catalysis.
\end{savaistu}

\courssub{Heterogeneous Catalysis}

In heterogeneous catalysis the reaction occurs at the \textbf{surface} of the catalyst. The mechanism is usually described in four stages:

\begin{enumerate}
\item \textbf{Diffusion} of reactant molecules to the surface.
\item \textbf{Adsorption}: reactant molecules bind to \cle{active sites} on the surface. Two types are distinguished: \emph{physisorption} (weak van der Waals attraction, easily reversed) and \emph{chemisorption} (formation of genuine chemical bonds with surface atoms — the step that activates the molecule, often by weakening or breaking its internal bonds).
\item \textbf{Surface reaction}: adsorbed species react together (or rearrange) on the surface; neighbouring molecules are held close and correctly oriented, and bonds are already partially broken, so the activation energy is low.
\item \textbf{Desorption}: product molecules leave the surface, freeing the active sites for new reactant molecules.
\end{enumerate}

\begin{popfigure}
\begin{center}
\begin{tikzpicture}[scale=0.9]
% metal surface
\fill[popblueL] (0,0) rectangle (12.6,-0.9);
\foreach \x in {0.45,1.35,...,12.15} {\fill[popblue] (\x,-0.45) circle (0.28);}
\node[font=\small, popdark] at (6.3,-1.35) {Metal catalyst surface (active sites)};
% stage 1: adsorption of A2 and B2
\draw[poporange, fill=poporange] (1.2,1.6) circle (0.22); \draw[poporange, fill=poporange] (1.7,1.6) circle (0.22);
\node[font=\small, poporange] at (1.45,2.15) {\ce{A2}};
\draw[popgreen, fill=popgreen] (2.6,1.9) circle (0.18); \draw[popgreen, fill=popgreen] (3.0,1.9) circle (0.18);
\node[font=\small, popgreen] at (2.8,2.4) {\ce{B2}};
\draw[->, thick, popmuted] (1.45,1.3) -- (1.45,0.35);
\draw[->, thick, popmuted] (2.8,1.6) -- (2.8,0.35);
\node[font=\small, popdark, align=center] at (2.1,3.0) {1. Adsorption};
% stage 2: chemisorbed atoms
\draw[poporange, fill=poporange] (5.0,0.25) circle (0.22); \draw[poporange, fill=poporange] (5.8,0.25) circle (0.22);
\draw[popgreen, fill=popgreen] (6.6,0.25) circle (0.18); \draw[popgreen, fill=popgreen] (7.3,0.25) circle (0.18);
\draw[thick, popdark] (5.0,0.05) -- (5.0,-0.15); \draw[thick, popdark] (5.8,0.05) -- (5.8,-0.15);
\draw[thick, popdark] (6.6,0.05) -- (6.6,-0.15); \draw[thick, popdark] (7.3,0.05) -- (7.3,-0.15);
\node[font=\small, popdark, align=center] at (6.15,1.2) {2. Bonds weaken\\(chemisorption)};
% stage 3: surface reaction
\draw[poporange, fill=poporange] (9.0,0.3) circle (0.22); \draw[popgreen, fill=popgreen] (9.35,0.3) circle (0.18);
\draw[poporange, fill=poporange] (10.1,0.3) circle (0.22); \draw[popgreen, fill=popgreen] (10.45,0.3) circle (0.18);
\node[font=\small, popdark, align=center] at (9.7,1.2) {3. Surface reaction\\(\ce{AB} forms)};
% stage 4: desorption
\draw[poporange, fill=poporange] (11.6,1.5) circle (0.22); \draw[popgreen, fill=popgreen] (11.95,1.5) circle (0.18);
\draw[poporange, fill=poporange] (12.3,2.1) circle (0.22); \draw[popgreen, fill=popgreen] (12.65,2.1) circle (0.18);
\draw[->, thick, popmuted] (11.2,0.5) -- (11.6,1.2);
\node[font=\small, popdark, align=center] at (11.9,2.9) {4. Desorption};
\end{tikzpicture}
\end{center}
\legende{Schematic of heterogeneous catalysis: gaseous reactants adsorb onto active sites of a solid metal surface, their bonds weaken, they react on the surface, and the products desorb, regenerating the sites.}
\end{popfigure}

\begin{attention}[Catalyst poisoning]
Impurities in the reactant stream can \textbf{poison} a heterogeneous catalyst by adsorbing strongly and irreversibly onto the active sites, blocking them. For example, sulfur compounds poison the iron catalyst of the Haber process, and lead compounds destroy the platinum of catalytic converters (which is why leaded petrol had to be phased out). Poisoned catalysts must be \textbf{regenerated} (e.g.\ by heating in a controlled atmosphere) or replaced — a major industrial cost. \textbf{Promoters}, by contrast, are substances (e.g.\ \ce{K2O} and \ce{Al2O3} added to iron in the Haber process) that are not catalysts themselves but \emph{increase the activity or lifetime} of the catalyst, for instance by improving its surface structure.
\end{attention}

\courssub{Industrial Applications}

\begin{center}
\begin{tabularx}{\linewidth}{|l|l|X|}
\hline
\rowcolor{popblueL} \textbf{Process} & \textbf{Catalyst} & \textbf{Reaction} \\
\hline
Haber process & Fe (with \ce{K2O}, \ce{Al2O3} promoters) & $\ce{N2(g) + 3H2(g) <=> 2NH3(g)}$ — ammonia for fertilisers \\
\hline
Contact process & \ce{V2O5} (solid) & $\ce{2SO2(g) + O2(g) <=> 2SO3(g)}$ — key step in \ce{H2SO4} manufacture \\
\hline
Catalytic converter & Pt / Rh on ceramic honeycomb & $\ce{2CO(g) + 2NO(g) -> N2(g) + 2CO2(g)}$; oxidation of unburnt hydrocarbons \\
\hline
Hydrogenation of oils & Ni (solid) & $\ce{C=C}$ bonds reduced by \ce{H2} — margarine manufacture \\
\hline
\end{tabularx}
\end{center}

In the Haber process, \ce{N2} and \ce{H2} chemisorb onto the iron surface; the strong \ce{N#N} triple bond is weakened, allowing hydrogenation step by step before \ce{NH3} desorbs. In the catalytic converter, the catalyst also shows \cle{selectivity}: platinum and rhodium accelerate the \emph{desired} reactions (oxidation of \ce{CO}, reduction of \ce{NO}) far more than competing ones. Selectivity — steering reactants towards one product among several possibilities — is one of the most valuable properties a catalyst can have.

\courssub{Measuring Catalytic Activity}

\begin{methode}[Determining catalytic activity by the rate enhancement factor]
\begin{enumerate}
\item Measure the rate constant of the \textbf{uncatalysed} reaction, $k_{\text{uncat}}$, under controlled conditions (fixed temperature, concentrations).
\item Repeat the experiment under \textbf{identical conditions} with a known amount of catalyst and measure $k_{\text{cat}}$ (e.g.\ from the gradient of a concentration--time curve, or from the half-life for a first-order reaction, $k = \ln 2 / t_{1/2}$).
\item Compute the \cle{rate enhancement factor} $\dfrac{k_{\text{cat}}}{k_{\text{uncat}}}$.
\item Interpret with the Arrhenius equation: $\ln\!\left(\dfrac{k_{\text{cat}}}{k_{\text{uncat}}}\right) = \dfrac{E_a(\text{uncat}) - E_a(\text{cat})}{RT}$, which gives the \textbf{reduction in activation energy} produced by the catalyst (assuming the pre-exponential factor $A$ is approximately the same for both pathways).
\end{enumerate}
\end{methode}

\begin{exoresolu}[Rate enhancement and lowering of $E_a$]
At \SI{298}{K}, a first-order reaction has $k_{\text{uncat}} = 3.0\times10^{-6}\ \mathrm{s^{-1}}$. In the presence of a catalyst, $k_{\text{cat}} = 1.5\times10^{-2}\ \mathrm{s^{-1}}$ at the same temperature. (a) Calculate the rate enhancement factor. (b) Calculate the reduction in activation energy. Take $R = 8.31\ \mathrm{J\,K^{-1}\,mol^{-1}}$.
\tcblower
\textbf{(a)} The enhancement factor is
\[
\frac{k_{\text{cat}}}{k_{\text{uncat}}} = \frac{1.5\times10^{-2}}{3.0\times10^{-6}} = 5.0\times10^{3}.
\]
The catalyst makes the reaction $5000$ times faster.

\textbf{(b)} From the Arrhenius equation, assuming the pre-exponential factor $A$ is similar for both pathways:
\begin{align*}
\ln\!\left(\frac{k_{\text{cat}}}{k_{\text{uncat}}}\right) &= \frac{E_a(\text{uncat}) - E_a(\text{cat})}{RT}\\[2pt]
\Delta E_a &= RT\ln(5.0\times10^{3}) = 8.31 \times 298 \times \ln(5000)\\
&= 8.31 \times 298 \times 8.52 \approx 2.11\times10^{4}\ \mathrm{J\,mol^{-1}}.
\end{align*}
\[
\boxed{\Delta E_a \approx 21\ \mathrm{kJ\,mol^{-1}}}
\]
A reduction of only about \SI{21}{kJ.mol^{-1}} in activation energy is enough to accelerate the reaction by a factor of $5000$ — a striking demonstration of the exponential sensitivity of rate to $E_a$.
\end{exoresolu}

\begin{exercice}[Catalysis in the laboratory and industry]
\begin{enumerate}
\item Explain, using an energy profile, why a catalyst increases the rate of the decomposition of hydrogen peroxide, $\ce{2H2O2(aq) -> 2H2O(l) + O2(g)}$, but does not change the total volume of oxygen finally obtained.
\item Classify each of the following as homogeneous or heterogeneous catalysis: (a) \ce{MnO2} powder added to \ce{H2O2} solution; (b) \ce{H+} in ester hydrolysis; (c) iron in the Haber process; (d) \ce{Fe^{2+}} in the \ce{S2O8^{2-}}/\ce{I-} reaction.
\item State two reasons why the catalyst in a catalytic converter is spread as a thin coating over a ceramic honeycomb of very large surface area.
\item A catalyst lowers $E_a$ from $75$ to $50\ \mathrm{kJ\,mol^{-1}}$ at \SI{300}{K}. Estimate the rate enhancement factor. ($R = 8.31\ \mathrm{J\,K^{-1}\,mol^{-1}}$.)
\end{enumerate}
\end{exercice}

\begin{corrige}[Exercise 1]
\textbf{1.} The catalyst (e.g.\ \ce{MnO2}) provides an alternative pathway with lower $E_a$, so more molecules cross the barrier per second and \ce{O2} is evolved faster. The catalyst changes neither $\Delta G$ nor the equilibrium position, so the \emph{final} amount of oxygen — fixed by the initial amount of \ce{H2O2} — is unchanged; only the \emph{time} to reach it is shorter.

\textbf{2.} (a) Heterogeneous (solid \ce{MnO2}, liquid reactant). (b) Homogeneous (all in aqueous solution). (c) Heterogeneous (solid Fe, gaseous reactants). (d) Homogeneous (all aqueous).

\textbf{3.} (i) Heterogeneous catalysis occurs at the surface, so a large surface area provides many more active sites and a higher rate; (ii) a thin coating uses very little of the expensive platinum and rhodium, reducing cost while the honeycomb still allows exhaust gases to flow freely.

\textbf{4.} $\dfrac{k_{\text{cat}}}{k_{\text{uncat}}} = \exp\!\left(\dfrac{\Delta E_a}{RT}\right) = \exp\!\left(\dfrac{25000}{8.31\times300}\right) = e^{10.03} \approx 2.3\times10^{4}$. The reaction is accelerated by a factor of about $23\,000$.
\end{corrige}

\begin{aretenir}[Key points on catalysis]
\begin{itemize}
\item A \cle{catalyst} speeds up a reaction by providing an alternative pathway of \textbf{lower activation energy}; it is regenerated and not consumed overall.
\item It does \textbf{not} change $\Delta H$, $\Delta G$ or $K$; it accelerates forward and reverse reactions equally, so equilibrium is reached faster but its position is unchanged.
\item \textbf{Homogeneous} catalysis (same phase) works through intermediate formation; \textbf{heterogeneous} catalysis (different phase) works by adsorption at active sites, surface reaction, then desorption.
\item Industrial cornerstones: Fe in the Haber process, \ce{V2O5} in the Contact process, Pt/Rh in catalytic converters.
\item Catalysts can be \textbf{poisoned} by impurities blocking active sites, assisted by \textbf{promoters}, and valued for their \textbf{selectivity}.
\item Catalytic activity is quantified by the enhancement factor $k_{\text{cat}}/k_{\text{uncat}}$, linked to the lowering of $E_a$ through the Arrhenius equation.
\end{itemize}
\end{aretenir}

\newpage

\coursec{Integration activity}

\coursec{Integration Activity: Kinetic Study of the Catalytic Decomposition of Hydrogen Peroxide}

\courssub{Learning Objectives}
By the end of this integration activity, you should be able to:
\begin{itemize}
    \item \cle{Design} a complete kinetic experiment: choose an appropriate monitoring method, control variables, and plan the variation of concentrations and temperature;
    \item \cle{Measure} and \cle{process} raw data (gas volumes versus time) to obtain initial rates;
    \item \cle{Determine} partial orders, the overall order and the rate constant of a reaction from initial-rate data and from graphical methods (half-life);
    \item \cle{Construct} an Arrhenius plot ($\ln k$ vs $1/T$) and extract the activation energy $E_a$;
    \item \cle{Interpret} results at the molecular level using collision theory and transition-state theory, and \cle{compare} homogeneous and heterogeneous catalysts;
    \item \cle{Communicate} scientific findings in a short written report, as expected in the GCE Advanced Level practical paper.
\end{itemize}

\begin{experience}[Context: The School Dispensary Problem]
The science club of Government Bilingual High School Nkolbisson, Yaoundé, has been asked to help the school dispensary. The dispensary stocks hydrogen peroxide solution, \ce{H2O2(aq)}, used as an antiseptic, but the pharmacist complains that bottles opened during the hot season ``go flat'' within weeks: the peroxide decomposes and no longer releases oxygen on wounds. The decomposition reaction is:
\[
\ce{2H2O2(aq) -> 2H2O(l) + O2(g)}
\]
The club decides to carry out a full kinetic study of this decomposition, catalysed in two ways:
\begin{itemize}
    \item by manganese(IV) oxide, \ce{MnO2(s)} (a \emph{heterogeneous} catalyst, insoluble in water);
    \item by iron(III) chloride solution, \ce{FeCl3(aq)} providing \ce{Fe^{3+}(aq)} ions (a \emph{homogeneous} catalyst).
\end{itemize}
The aims are to establish the rate law, measure the activation energy with each catalyst, and advise the pharmacist on storage conditions.

\textbf{Laboratory equipment provided:} a \SI{100}{cm^3} gas syringe connected to a conical flask, a stopwatch, a thermostatted water bath (\SIrange{20}{60}{\ensuremath{{}^{\circ}\mathrm{C}}}), standard \SI{1.0}{mol.L^{-1}} \ce{H2O2} solution, \ce{MnO2} powder, \SI{0.10}{mol.L^{-1}} \ce{FeCl3} solution, distilled water and the usual glassware.
\end{experience}

\courssub{Experimental Programme}

\begin{experience}[Experiment 1 -- Choice of Monitoring Method]
The club rejects titrimetric sampling (the reaction is too fast at room temperature to quench and titrate conveniently) and chooses to follow the volume of \ce{O2(g)} evolved with the gas syringe.
\end{experience}

\begin{experience}[Experiment 2 -- Initial Rates at \SI{25}{\ensuremath{{}^{\circ}\mathrm{C}}} with \ce{MnO2}]
A fixed mass of \SI{0.20}{g} of \ce{MnO2} is used in each run; total volume is kept at \SI{50}{cm^3} by adding distilled water. Initial rates are obtained from the tangent at $t = 0$ on the volume--time curves.

\begin{center}
\begin{tabularx}{\linewidth}{|c|X|X|X|}
\hline
\rowcolor{popblueL} Run & $[\ce{H2O2}]$ / \si{mol.L^{-1}} & mass \ce{MnO2} / g & Initial rate / \si{cm^3.min^{-1}} \\
\hline
1 & 0.20 & 0.20 & 12.0 \\
\hline
2 & 0.40 & 0.20 & 24.0 \\
\hline
3 & 0.60 & 0.20 & 36.0 \\
\hline
4 & 0.40 & 0.40 & 48.0 \\
\hline
\end{tabularx}
\end{center}
\end{experience}

\begin{experience}[Experiment 3 -- Temperature Dependence of the Rate Constant]
Using $[\ce{H2O2}] = \SI{0.40}{mol.L^{-1}}$, the initial rate is measured at three temperatures with each catalyst. The reaction is first order in \ce{H2O2} in both cases, so the rate constant $k$ (units \si{s^{-1}}) is deduced from the initial rate $v_0 = k[\ce{H2O2}]_0$.

\begin{center}
\begin{tabularx}{\linewidth}{|c|X|X|}
\hline
\rowcolor{popblueL} $T$ / K & $k$ (\ce{MnO2}) / \si{s^{-1}} & $k$ (\ce{Fe^{3+}}) / \si{s^{-1}} \\
\hline
298 & $1.0 \times 10^{-3}$ & $2.0 \times 10^{-4}$ \\
\hline
308 & $2.2 \times 10^{-3}$ & $5.0 \times 10^{-4}$ \\
\hline
318 & $4.6 \times 10^{-3}$ & $1.2 \times 10^{-3}$ \\
\hline
\end{tabularx}
\end{center}
\end{experience}

\begin{experience}[Experiment 4 -- Half-Life Check with \ce{Fe^{3+}} at \SI{298}{K}]
Starting from $[\ce{H2O2}]_0 = \SI{0.40}{mol.L^{-1}}$, the club finds that the concentration falls to \SI{0.20}{mol.L^{-1}} after \SI{58}{min} and to \SI{0.10}{mol.L^{-1}} after a further \SI{58}{min}.

\textbf{Data:} $R = \SI{8.314}{J.K^{-1}.mol^{-1}}$; molar volume of gas at \SI{25}{\ensuremath{{}^{\circ}\mathrm{C}}}: $V_m = \SI{24.0}{dm^3.mol^{-1}}$.
\end{experience}

\courssub{Tasks}

\begin{enumerate}
    \item Explain why the gas syringe method is appropriate here, state two precautions needed to obtain reliable volumes, and identify two variables that must be kept constant in Experiment 2.
    \item Convert the initial rate of Run 1 into \si{mol.L^{-1}.s^{-1}} of \ce{H2O2} decomposed.
    \item Using Runs 1--3, determine the partial order with respect to \ce{H2O2}. Using Runs 2 and 4, determine the partial order with respect to the \ce{MnO2} catalyst. Write the rate equation and state the overall order.
    \item Calculate the rate constant $k$ at \SI{298}{K} for the \ce{Fe^{3+}}-catalysed reaction from the half-life data of Experiment 4, and check its consistency with the table of Experiment 3.
    \item Plot (or describe the construction of) the Arrhenius graph $\ln k$ versus $1/T$ for each catalyst, and calculate the activation energy $E_a$ for each.
    \item Using transition-state theory, explain why \ce{MnO2} gives a larger rate constant than \ce{Fe^{3+}} at the same temperature, and state what can be said about $\Delta H^{\ddagger}$ for the two catalysed pathways.
    \item Classify each catalyst as homogeneous or heterogeneous, give one advantage and one disadvantage of each type in this context, and write a short recommendation to the pharmacist on how to store the \ce{H2O2} stock.
\end{enumerate}

\courssub{Model Answers with Examiner's Commentary}

\begin{exoresolu}[Task 1: Choice of Method, Precautions and Controlled Variables]
\textbf{Statement:} Explain why the gas syringe method is appropriate here, state two precautions needed to obtain reliable volumes, and identify two variables that must be kept constant in Experiment 2.

\tcblower
\textbf{Detailed Solution:}

\textbf{Why the gas syringe is appropriate:}
The reaction produces a gaseous product, \ce{O2}, whose volume is directly proportional to the extent of reaction (amount of \ce{H2O2} decomposed). The gas syringe provides a \emph{continuous, non-invasive, rapid} measurement without the need for sampling and quenching, which is essential for a reaction that is relatively fast at room temperature (especially when catalysed). It directly yields $V(t)$ curves from which initial rates (tangent at $t=0$) are easily obtained.

\textbf{Two essential precautions for reliable volumes:}
\begin{enumerate}
    \item \textbf{Gas-tight apparatus:} Check all connections (bung, tubing, syringe barrel) for leaks before each run. A leak leads to systematically low volumes and underestimation of the rate.
    \item \textbf{Temperature equilibration and reading:} The gas in the syringe must be at the same temperature as the water bath (or room temperature) when the volume is read. Allow the syringe to equilibrate; read the volume at eye level, at the bottom of the meniscus (or the inner edge of the plunger), and record the temperature. If the gas is warmer than the surroundings, the volume will be overestimated (Charles's law).
\end{enumerate}
\textit{Additional good practice:} Swirl the flask gently and continuously so that the \ce{MnO2} powder remains suspended and its surface area stays constant; start the stopwatch \emph{exactly} at the moment of mixing (adding the catalyst to the peroxide solution).

\textbf{Two variables held constant in Experiment 2 (Runs 1--4):}
\begin{enumerate}
    \item Total volume of the reaction mixture (\SI{50}{cm^3}) -- ensured by making up with distilled water.
    \item Temperature (\SI{25}{\ensuremath{{}^{\circ}\mathrm{C}}} = \SI{298}{K}) -- maintained by the water bath.
\end{enumerate}
(For Runs 1--3, the mass of \ce{MnO2} is also constant at \SI{0.20}{g}; for Runs 2 and 4, $[\ce{H2O2}]$ is constant at \SI{0.40}{mol.L^{-1}}.)
\end{exoresolu}

\begin{exoresolu}[Task 2: Conversion of Initial Rate to Concentration Units]
\textbf{Statement:} Convert the initial rate of Run 1 into \si{mol.L^{-1}.s^{-1}} of \ce{H2O2} decomposed.

\tcblower
\textbf{Detailed Solution:}

Run 1 gives an initial rate of \ce{O2} evolution: $v_{\ce{O2}} = \SI{12.0}{cm^3.min^{-1}}$.

\textbf{Step 1: Convert volume rate to molar rate of \ce{O2}.}
Molar volume at \SI{25}{\ensuremath{{}^{\circ}\mathrm{C}}}: $V_m = \SI{24.0}{dm^3.mol^{-1}} = \SI{24000}{cm^3.mol^{-1}}$.
\[
\frac{dn_{\ce{O2}}}{dt} = \frac{12.0\ \mathrm{cm^3.min^{-1}}}{24000\ \mathrm{cm^3.mol^{-1}}} = 5.00 \times 10^{-4}\ \mathrm{mol.min^{-1}}.
\]

\textbf{Step 2: Use stoichiometry to get molar rate of \ce{H2O2} decomposition.}
From the balanced equation: $\ce{2H2O2 -> 2H2O + O2}$, 2 mol \ce{H2O2} produce 1 mol \ce{O2}.
\[
-\frac{dn_{\ce{H2O2}}}{dt} = 2 \times \frac{dn_{\ce{O2}}}{dt} = 2 \times 5.00 \times 10^{-4} = 1.00 \times 10^{-3}\ \mathrm{mol.min^{-1}}.
\]

\textbf{Step 3: Convert to per second.}
\[
-\frac{dn_{\ce{H2O2}}}{dt} = \frac{1.00 \times 10^{-3}}{60} = 1.667 \times 10^{-5}\ \mathrm{mol.s^{-1}}.
\]

\textbf{Step 4: Divide by solution volume to obtain concentration rate.}
Total volume $V = \SI{50}{cm^3} = \SI{0.0500}{dm^3} = \SI{0.0500}{L}$.
\[
v_0 = -\frac{1}{V}\frac{dn_{\ce{H2O2}}}{dt} = \frac{1.667 \times 10^{-5}\ \mathrm{mol.s^{-1}}}{0.0500\ \mathrm{L}} = 3.33 \times 10^{-4}\ \mathrm{mol.L^{-1}.s^{-1}}.
\]

\textbf{Answer:} The initial rate of decomposition of \ce{H2O2} in Run 1 is $\mathbf{3.33 \times 10^{-4}\ \si{mol.L^{-1}.s^{-1}}}$ (3 significant figures).
\end{exoresolu}

\begin{exoresolu}[Task 3: Determination of Orders and Rate Equation]
\textbf{Statement:} Using Runs 1--3, determine the partial order with respect to \ce{H2O2}. Using Runs 2 and 4, determine the partial order with respect to the \ce{MnO2} catalyst. Write the rate equation and state the overall order.

\tcblower
\textbf{Detailed Solution:}

\textbf{Order with respect to \ce{H2O2} (Runs 1, 2, 3):} Mass of \ce{MnO2} constant (\SI{0.20}{g}), $T$ constant.
\begin{itemize}
    \item Run 1 $\to$ Run 2: $[\ce{H2O2}]$ doubles ($0.20 \to \SI{0.40}{mol.L^{-1}}$); initial rate doubles ($12.0 \to \SI{24.0}{cm^3.min^{-1}}$).
    \item Run 1 $\to$ Run 3: $[\ce{H2O2}]$ triples ($0.20 \to \SI{0.60}{mol.L^{-1}}$); initial rate triples ($12.0 \to \SI{36.0}{cm^3.min^{-1}}$).
\end{itemize}
Rate $\propto [\ce{H2O2}]^1$. \cle{Partial order = 1 (first order)} in \ce{H2O2}.

\textbf{Order with respect to \ce{MnO2} (Runs 2 and 4):} $[\ce{H2O2}]$ constant (\SI{0.40}{mol.L^{-1}}), $T$ constant.
\begin{itemize}
    \item Run 2 $\to$ Run 4: mass of \ce{MnO2} doubles ($0.20 \to \SI{0.40}{g}$); initial rate doubles ($24.0 \to \SI{48.0}{cm^3.min^{-1}}$).
\end{itemize}
Rate $\propto (\text{mass of catalyst})^1$. Since the catalyst is a solid, its ``concentration'' is proportional to its mass (or surface area) for a given particle size. \cle{Partial order = 1 (first order)} in catalyst.

\textbf{Rate equation:}
\[
v = k\,[\ce{H2O2}]^1\,[\text{catalyst}]^1 \quad \text{or} \quad v = k'\,[\ce{H2O2}] \quad \text{with } k' = k[\text{catalyst}] \text{ (pseudo-first-order constant)}.
\]

\textbf{Overall order} = $1 + 1 = \mathbf{2}$ (second order overall).

\textbf{Consistency check:} With a fixed amount of catalyst, the reaction behaves as \emph{pseudo-first order} in \ce{H2O2}. This is confirmed by Experiment 4 (constant half-life).
\end{exoresolu}

\begin{exoresolu}[Task 4: Rate Constant from Half-Life Data]
\textbf{Statement:} Calculate the rate constant $k$ at \SI{298}{K} for the \ce{Fe^{3+}}-catalysed reaction from the half-life data of Experiment 4, and check its consistency with the table of Experiment 3.

\tcblower
\textbf{Detailed Solution:}

Experiment 4 shows: $[\ce{H2O2}]$ falls from \SI{0.40}{mol.L^{-1}} to \SI{0.20}{mol.L^{-1}} in \SI{58}{min}, and from \SI{0.20}{mol.L^{-1}} to \SI{0.10}{mol.L^{-1}} in a further \SI{58}{min}.
The half-life $t_{1/2}$ is \textbf{constant} (\SI{58}{min}) independent of concentration. This is the hallmark of a \cle{first-order} reaction in \ce{H2O2}.

For a first-order reaction:
\[
k = \frac{\ln 2}{t_{1/2}}.
\]
Convert $t_{1/2}$ to seconds: $t_{1/2} = 58\ \mathrm{min} \times 60\ \mathrm{s.min^{-1}} = \SI{3480}{s}$.
\[
k = \frac{0.6931}{3480\ \mathrm{s}} = 1.99 \times 10^{-4}\ \mathrm{s^{-1}} \approx \mathbf{2.0 \times 10^{-4}\ \si{s^{-1}}}.
\]

\textbf{Consistency check:} The table in Experiment 3 gives $k(\ce{Fe^{3+}})$ at \SI{298}{K} = $\mathbf{2.0 \times 10^{-4}\ \si{s^{-1}}}$. The two values are in \textbf{exact agreement}, confirming the first-order kinetics and the reliability of the measurements.
\end{exoresolu}

\begin{exoresolu}[Task 5: Arrhenius Plots and Activation Energies]
\textbf{Statement:} Plot (or describe the construction of) the Arrhenius graph $\ln k$ versus $1/T$ for each catalyst, and calculate the activation energy $E_a$ for each.

\tcblower
\textbf{Detailed Solution:}

The Arrhenius equation in linear form:
\[
\ln k = \ln A - \frac{E_a}{R}\,\frac{1}{T}.
\]
A plot of $\ln k$ (y-axis) against $1/T$ (x-axis) yields a straight line of slope $m = -E_a/R$.

\textbf{Step 1: Tabulate $\ln k$ and $1/T$.}

\begin{center}
\begin{tabularx}{\linewidth}{|c|X|X|X|X|}
\hline
\rowcolor{popblueL} $T$ / K & $10^3/T$ / \si{K^{-1}} & $\ln k$ (\ce{MnO2}) & $\ln k$ (\ce{Fe^{3+}}) \\
\hline
298 & 3.356 & $\ln(1.0\times10^{-3}) = -6.908$ & $\ln(2.0\times10^{-4}) = -8.517$ \\
\hline
308 & 3.247 & $\ln(2.2\times10^{-3}) = -6.119$ & $\ln(5.0\times10^{-4}) = -7.601$ \\
\hline
318 & 3.145 & $\ln(4.6\times10^{-3}) = -5.381$ & $\ln(1.2\times10^{-3}) = -6.725$ \\
\hline
\end{tabularx}
\end{center}

\textbf{Step 2: Determine slopes (using the two extreme points for best accuracy).}

\textit{For \ce{MnO2}:}
\[
m_{\ce{MnO2}} = \frac{\ln k_{318} - \ln k_{298}}{1/318 - 1/298} = \frac{-5.381 - (-6.908)}{3.145\times10^{-3} - 3.356\times10^{-3}} = \frac{1.527}{-0.211\times10^{-3}} = -7237\ \mathrm{K}.
\]
\[
E_a(\ce{MnO2}) = -m \times R = 7237 \times 8.314 = \SI{60160}{J.mol^{-1}} \approx \mathbf{60.2\ kJ.mol^{-1}}.
\]

\textit{For \ce{Fe^{3+}}:}
\[
m_{\ce{Fe^{3+}}} = \frac{-6.725 - (-8.517)}{3.145\times10^{-3} - 3.356\times10^{-3}} = \frac{1.792}{-0.211\times10^{-3}} = -8493\ \mathrm{K}.
\]
\[
E_a(\ce{Fe^{3+}}) = 8493 \times 8.314 = \SI{70600}{J.mol^{-1}} \approx \mathbf{70.6\ kJ.mol^{-1}}.
\]

\textbf{Note:} In an examination, you would plot the points on graph paper, draw the best-fit straight line, and measure the slope using a large triangle. The values above are the precise calculations from the given data; graphical determination typically yields $E_a(\ce{MnO2}) \approx \SI{58\text{--}60}{kJ.mol^{-1}}$ and $E_a(\ce{Fe^{3+}}) \approx \SI{68\text{--}71}{kJ.mol^{-1}}$.

\begin{aretenir}[Arrhenius Plot -- Examination Technique]
\begin{itemize}
    \item Always convert $T$ to Kelvin and $k$ to consistent units (\si{s^{-1}} here).
    \item Calculate $1/T$ in $\si{K^{-1}}$ (often $\times 10^3$ for convenient scaling).
    \item Plot $\ln k$ vs $1/T$; the slope is negative.
    \item $E_a = -\text{slope} \times R$. Give $E_a$ in \si{kJ.mol^{-1}} (divide by 1000).
    \item The intercept gives $\ln A$ (pre-exponential factor), but this is rarely asked at A Level.
\end{itemize}
\end{aretenir}
\end{exoresolu}

\begin{exoresolu}[Task 6: Transition-State Theory Interpretation]
\textbf{Statement:} Using transition-state theory, explain why \ce{MnO2} gives a larger rate constant than \ce{Fe^{3+}} at the same temperature, and state what can be said about $\Delta H^{\ddagger}$ for the two catalysed pathways.

\tcblower
\textbf{Detailed Solution:}

In Transition-State Theory (TST), the rate constant is given by the Eyring equation:
\[
k = \frac{k_B T}{h} e^{-\Delta G^{\ddagger}/RT} = \frac{k_B T}{h} e^{\Delta S^{\ddagger}/R} e^{-\Delta H^{\ddagger}/RT},
\]
where $\Delta G^{\ddagger} = \Delta H^{\ddagger} - T\Delta S^{\ddagger}$ is the Gibbs free energy of activation.

\textbf{Comparison at \SI{298}{K}:}
$k(\ce{MnO2}) = 1.0 \times 10^{-3}\ \si{s^{-1}}$; $k(\ce{Fe^{3+}}) = 2.0 \times 10^{-4}\ \si{s^{-1}}$.
$k(\ce{MnO2})$ is \textbf{5 times larger} than $k(\ce{Fe^{3+}})$. Therefore, the \ce{MnO2}-catalysed pathway has a \textbf{lower} $\Delta G^{\ddagger}$ (more negative $\Delta G^{\ddagger}$ means faster reaction).

\textbf{Enthalpy of activation $\Delta H^{\ddagger}$:}
For reactions in solution, $E_a \approx \Delta H^{\ddagger} + RT$ (the $RT$ term is small, $\approx \SI{2.5}{kJ.mol^{-1}}$ at \SI{298}{K}). Thus $E_a$ is a good approximation for $\Delta H^{\ddagger}$.
From Task 5:
\begin{itemize}
    \item $E_a(\ce{MnO2}) \approx \SI{60}{kJ.mol^{-1}} \Rightarrow \Delta H^{\ddagger}_{\ce{MnO2}} \approx \SI{58}{kJ.mol^{-1}}$.
    \item $E_a(\ce{Fe^{3+}}) \approx \SI{71}{kJ.mol^{-1}} \Rightarrow \Delta H^{\ddagger}_{\ce{Fe^{3+}}} \approx \SI{68}{kJ.mol^{-1}}$.
\end{itemize}
$\Delta H^{\ddagger}$ is \textbf{significantly lower (by $\approx \SI{10}{kJ.mol^{-1}}$)} for the \ce{MnO2} pathway.

\textbf{Conclusion:} The \ce{MnO2} surface stabilises the activated complex (transition state) more effectively than the \ce{Fe^{3+}} ions in solution, lowering the enthalpy barrier. The entropy of activation $\Delta S^{\ddagger}$ may also play a role: on a surface, reactants can be pre-adsorbed in a favourable orientation (less negative $\Delta S^{\ddagger}$), further reducing $\Delta G^{\ddagger}$. However, the dominant factor here is the lower $\Delta H^{\ddagger}$.
\end{exoresolu}

\begin{exoresolu}[Task 7: Catalyst Classification, Pros/Cons and Recommendation]
\textbf{Statement:} Classify each catalyst as homogeneous or heterogeneous, give one advantage and one disadvantage of each type in this context, and write a short recommendation to the pharmacist on how to store the \ce{H2O2} stock.

\tcblower
\textbf{Detailed Solution:}

\begin{center}
\begin{tabularx}{\linewidth}{|l|X|X|}
\hline
\rowcolor{popblueL} Catalyst & Classification & Advantage / Disadvantage in this context \\
\hline
\ce{MnO2(s)} & \textbf{Heterogeneous} (solid catalyst, liquid reactants) & \textbf{Adv:} Easily separated by filtration, can be reused indefinitely (cost-effective). \\
& & \textbf{Disadv:} Activity depends on surface area (particle size); can be poisoned by impurities adsorbing on active sites; mass transfer limitations possible. \\
\hline
\ce{Fe^{3+}(aq)} & \textbf{Homogeneous} (same phase -- aqueous) & \textbf{Adv:} Uniform mixing, every catalytic ion is accessible; no mass transfer limitations; often higher selectivity. \\
& & \textbf{Disadv:} Cannot be separated from the product (contaminates the \ce{H2O2} solution); difficult to recover/reuse; here it is less active (higher $E_a$, lower $k$). \\
\hline
\end{tabularx}
\end{center}

\textbf{Recommendation to the Pharmacist:}
\begin{quote}
``Store the hydrogen peroxide stock in a \textbf{cool, dark cupboard} (ideally below \SI{20}{\ensuremath{{}^{\circ}\mathrm{C}}}), in \textbf{tightly closed, opaque (brown glass) bottles}. Low temperature exponentially decreases the rate constant (Arrhenius law); light (especially UV) accelerates photodecomposition. Keep bottles away from dust, rust, and metal ions (Fe, Cu, Mn) which act as homogeneous or heterogeneous catalysts. During the hot season in Yaoundé, consider refrigeration if possible, and avoid leaving bottles open. Replace stock regularly as decomposition is inevitable over time.''
\end{quote}

\begin{savaistu}[Real-World Context]
Commercial \ce{H2O2} solutions contain stabilisers (e.g., phosphoric acid, stannate) that chelate metal ions and slow decomposition. The ``fizzing'' on a wound is due to catalase enzyme in blood/tissues -- a biological catalyst far more efficient than \ce{MnO2} or \ce{Fe^{3+}}!
\end{savaistu}
\end{exoresolu}

\begin{aretenir}[Key Takeaways for the GCE Advanced Level]
\begin{itemize}
    \item \textbf{Gas syringe method:} Ideal for gas-evolving reactions; continuous data; initial rate = tangent at $t=0$.
    \item \textbf{Initial rates method:} Compare rates where only one concentration changes; order $n$ from $(rate_2/rate_1) = ([A]_2/[A]_1)^n$.
    \item \textbf{Half-life method:} Constant $t_{1/2}$ $\Rightarrow$ first order; $t_{1/2} \propto 1/[A]_0$ $\Rightarrow$ second order.
    \item \textbf{Arrhenius plot:} $\ln k$ vs $1/T$ straight line; slope $= -E_a/R$; $E_a$ in \si{kJ.mol^{-1}}.
    \item \textbf{Transition-State Theory:} $k \propto e^{-\Delta G^{\ddagger}/RT}$; lower $\Delta H^{\ddagger}$ (or more positive $\Delta S^{\ddagger}$) $\Rightarrow$ larger $k$.
    \item \textbf{Catalysis:} Homogeneous = same phase (hard to separate); Heterogeneous = different phase (easy to separate, surface area matters).
\end{itemize}
\end{aretenir}

\begin{attention}[Common Pitfalls in Kinetics Questions]
\begin{itemize}
    \item Forgetting to convert time to \textbf{seconds} when calculating $k$ (units of $k$ depend on overall order).
    \item Using $V_m$ at STP (\SI{22.4}{dm^3}) instead of the given molar volume at the experimental temperature.
    \item Confusing \emph{order} with \emph{stoichiometric coefficient} (they are unrelated for complex reactions).
    \item Plotting $k$ vs $T$ instead of $\ln k$ vs $1/T$ for Arrhenius.
    \item Stating that a catalyst ``lowers the activation energy'' without mentioning it provides an \emph{alternative reaction pathway} with a lower $E_a$.
    \item Not specifying that for a heterogeneous catalyst, the rate depends on \emph{surface area} (mass for fixed particle size), not ``concentration''.
\end{itemize}
\end{attention}

\newpage

\coursec{Methods and worked examples}

\courssub{Methods and Worked Examples}

\begin{methode}[Determining Reaction Rates from Concentration--Time Data]
\begin{enumerate}
\item \textbf{Identify the data type:} You may be given a table of concentration versus time, a graph, or a description of an experiment (gas volume, absorbance, conductivity, pressure).
\item \textbf{Plot the graph:} If tabular data are given, plot \textbf{concentration (or a quantity proportional to it: volume, absorbance, pressure) on the $y$-axis} against \textbf{time on the $x$-axis}. Draw a smooth curve through the points.
\item \textbf{Average rate over $\Delta t$:} $\text{Rate}_{\text{avg}} = -\dfrac{\Delta[\text{Reactant}]}{\Delta t} = -\dfrac{[\text{R}]_{t_2} - [\text{R}]_{t_1}}{t_2 - t_1}$ (the minus sign makes the rate positive, since reactant concentration decreases). For products, $\text{Rate}_{\text{avg}} = +\dfrac{\Delta[\text{Product}]}{\Delta t}$.
\item \textbf{Instantaneous rate at time $t$:} Draw a \textbf{tangent} to the curve at time $t$ with a ruler. Calculate its gradient: $\text{Rate}_{\text{inst}} = \dfrac{\text{change in concentration}}{\text{change in time}} = \text{gradient of tangent}$.
\item \textbf{Initial rate ($t=0$):} Draw the tangent at $t=0$. This is the most reliable rate for comparing experiments with different initial concentrations, because at $t=0$ the concentrations are known exactly and no products have accumulated.
\item \textbf{Units:} Always express rate in $\mathrm{mol\,dm^{-3}\,s^{-1}}$. Convert time to seconds if given in minutes, and volumes to $\mathrm{dm^3}$ if given in $\mathrm{cm^3}$.
\item \textbf{Stoichiometry link:} For $\ce{aA + bB -> cC + dD}$,
\[ \text{Rate} = -\frac{1}{a}\frac{d[\ce{A}]}{dt} = -\frac{1}{b}\frac{d[\ce{B}]}{dt} = +\frac{1}{c}\frac{d[\ce{C}]}{dt} = +\frac{1}{d}\frac{d[\ce{D}]}{dt}. \]
\end{enumerate}
\end{methode}

\begin{exoresolu}[Rate Determination from Gas Evolution Data]
In an experiment, $\ce{CaCO3(s)}$ reacts with excess $\ce{HCl(aq)}$: $\ce{CaCO3(s) + 2HCl(aq) -> CaCl2(aq) + H2O(l) + CO2(g)}$. The volume of $\ce{CO2}$ collected at room temperature and pressure (RTP, molar volume $= 24.0\ \mathrm{dm^3\,mol^{-1}}$) is recorded:

\begin{center}
\begin{tabular}{|c|ccccccc|}\hline
Time (s) & 0 & 30 & 60 & 90 & 120 & 150 & 180 \\ \hline
Vol $\ce{CO2}$ (cm$^3$) & 0 & 42 & 72 & 93 & 108 & 117 & 120 \\ \hline
\end{tabular}
\end{center}

\textbf{(a)} Plot the graph of volume against time. \\
\textbf{(b)} Determine the initial rate of reaction in $\mathrm{mol\,dm^{-3}\,s^{-1}}$ if the volume of the reaction mixture is $50.0\ \mathrm{cm^3}$. \\
\textbf{(c)} Determine the instantaneous rate at $t = 60\ \mathrm{s}$. \\
\textbf{(d)} Explain why the rate decreases with time.
\tcblower
\textbf{Detailed Solution}

\textbf{(a) Graph:} Plot Volume ($\mathrm{cm^3}$) on the $y$-axis, Time (s) on the $x$-axis. Points: $(0,0), (30,42), (60,72), (90,93), (120,108), (150,117), (180,120)$. Draw a smooth curve starting steep and levelling off at $120\ \mathrm{cm^3}$ (all the $\ce{CaCO3}$ has reacted).
\begin{center}
\begin{tikzpicture}
\begin{axis}[
width=12cm, height=8cm,
xlabel={Time / s}, ylabel={Volume $\ce{CO2}$ / cm$^3$},
xmin=0, xmax=190, ymin=0, ymax=130,
grid=major, tick style={thick},
]
\addplot[popblue, thick, smooth, mark=*] coordinates {
(0,0) (30,42) (60,72) (90,93) (120,108) (150,117) (180,120)
};
\addplot[poporange, dashed, domain=0:60] {1.4*x};
\addplot[popgreen, dashed, domain=30:90] {0.85*x + 21};
\end{axis}
\end{tikzpicture}
\legende{Volume of $\ce{CO2}$ against time. Orange dashed line: tangent at $t=0$ (initial rate). Green dashed line: tangent at $t=60$ s (instantaneous rate).}
\end{center}

\textbf{(b) Initial Rate:} Gradient of the tangent at $t=0$. Using the first two points as an approximation of the tangent:
\[ \text{Gradient} \approx \frac{42 - 0}{30 - 0} = 1.40\ \mathrm{cm^3\,s^{-1}}. \]
Convert to $\mathrm{dm^3\,s^{-1}}$: $1.40\ \mathrm{cm^3\,s^{-1}} = 1.40 \times 10^{-3}\ \mathrm{dm^3\,s^{-1}}$.
Moles of $\ce{CO2}$ per second $= \dfrac{1.40 \times 10^{-3}}{24.0} = 5.83 \times 10^{-5}\ \mathrm{mol\,s^{-1}}$.
Rate of reaction (with respect to $\ce{CO2}$ formation, per unit volume of mixture):
\[ \text{Rate} = \frac{1}{V_{\text{mixture}}} \times \frac{d n_{\ce{CO2}}}{dt} = \frac{5.83 \times 10^{-5}}{50.0 \times 10^{-3}} = \boxed{1.17 \times 10^{-3}\ \mathrm{mol\,dm^{-3}\,s^{-1}}}. \]
(The rate of disappearance of $\ce{HCl}$ would be $-\frac{1}{2}\frac{d[\ce{HCl}]}{dt}$ equal to this same value, since 2 mol of $\ce{HCl}$ are consumed per 1 mol of $\ce{CO2}$ formed.)

\textbf{(c) Instantaneous Rate at $t=60\ \mathrm{s}$:} Draw the tangent at $t=60\ \mathrm{s}$. Estimate its gradient using two points on the tangent, e.g. at $t=30\ \mathrm{s}$ ($V \approx 46.5\ \mathrm{cm^3}$ on the tangent) and $t=90\ \mathrm{s}$ ($V \approx 97.5\ \mathrm{cm^3}$):
\[ \text{Gradient} \approx \frac{97.5 - 46.5}{90 - 30} = \frac{51}{60} = 0.85\ \mathrm{cm^3\,s^{-1}}. \]
Moles per second $= \dfrac{0.85 \times 10^{-3}}{24.0} = 3.54 \times 10^{-5}\ \mathrm{mol\,s^{-1}}$.
\[ \text{Rate} = \frac{3.54 \times 10^{-5}}{0.0500} = \boxed{7.1 \times 10^{-4}\ \mathrm{mol\,dm^{-3}\,s^{-1}}}. \]

\textbf{(d) Reason for the decrease:} The rate depends on the concentration of $\ce{HCl}$ (the reactant in solution). As the reaction proceeds, $\ce{HCl}$ is consumed and its concentration falls. Since $\text{Rate} = k[\ce{HCl}]^n$, a lower concentration gives a lower rate, so the curve becomes less steep. The rate falls to zero when the $\ce{CaCO3}$ is used up (the acid is in excess).
\end{exoresolu}

\begin{methode}[Titrimetric (Chemical) Method for Rate Determination]
\begin{enumerate}
\item \textbf{Principle:} Samples (aliquots) of the reaction mixture are withdrawn at measured times and the reaction in each sample is \textbf{quenched} (stopped or drastically slowed) by, for example, adding ice-cold water, adding a reagent that removes a catalyst, or neutralising an acid.
\item \textbf{Titration:} The quenched sample is titrated against a standard solution to find the concentration of a reactant remaining or product formed (e.g. $\ce{H+}$, $\ce{I2}$, $\ce{H2O2}$).
\item \textbf{Data processing:} The titre $V_t$ is proportional to the concentration of the titrated species at time $t$. Plot $V_t$ against $t$, or test integrated rate laws (e.g. $\ln V_t$ against $t$ for first order when $V_\infty = 0$).
\item \textbf{Key reactions:}
    \begin{itemize}
    \item \textbf{Acid-catalysed ester hydrolysis:} $\ce{CH3COOCH3 + H2O <=>[H+] CH3COOH + CH3OH}$. Quench by adding the aliquot to ice-cold water; titrate the total acid with standard $\ce{NaOH}$. (The increase in titre above the value at $t=0$ measures the $\ce{CH3COOH}$ formed.)
    \item \textbf{Propanone + iodine:} $\ce{CH3COCH3 + I2 ->[H+] CH3COCH2I + H+ + I-}$. Quench with $\ce{NaHCO3}$ (neutralises the $\ce{H+}$ catalyst); titrate the remaining $\ce{I2}$ with standard $\ce{Na2S2O3}$ using starch indicator.
    \item \textbf{Decomposition of hydrogen peroxide:} $\ce{2H2O2 -> 2H2O + O2}$. Quench by adding the aliquot to cold dilute acid; titrate remaining $\ce{H2O2}$ with standard $\ce{KMnO4}$.
    \end{itemize}
\item \textbf{Advantage / disadvantage:} It gives a direct chemical measure of concentration; but it is discontinuous (sampling errors), relatively slow, and quenching may not be perfectly instantaneous.
\end{enumerate}
\end{methode}

\begin{exoresolu}[Kinetics of the Acid-Catalysed Propanone--Iodine Reaction]
The reaction $\ce{CH3COCH3(aq) + I2(aq) ->[H+] CH3COCH2I(aq) + H+(aq) + I-(aq)}$ is studied with propanone and acid in large excess. $10.0\ \mathrm{cm^3}$ aliquots are withdrawn at intervals, quenched with $\ce{NaHCO3}$, and titrated with $0.0100\ \mathrm{mol\,dm^{-3}}\ \ce{Na2S2O3}$ ($\ce{I2 + 2S2O3^2- -> 2I- + S4O6^2-}$).

\begin{center}
\begin{tabular}{|c|cccccc|}\hline
Time (min) & 0 & 5 & 10 & 15 & 20 & $\infty$ \\ \hline
Titre ($\mathrm{cm^3}$) & 25.0 & 18.4 & 13.5 & 9.9 & 7.3 & 0.0 \\ \hline
\end{tabular}
\end{center}

\textbf{(a)} Show that the reaction is first order with respect to $\ce{I2}$. \\
\textbf{(b)} Calculate the rate constant $k$ (in $\mathrm{s^{-1}}$). \\
\textbf{(c)} Determine the half-life $t_{1/2}$. \\
\textbf{(d)} Why is the order with respect to $\ce{CH3COCH3}$ and $\ce{H+}$ not determined by this experiment?
\tcblower
\textbf{Detailed Solution}

\textbf{(a) First-order proof:} For a first-order reaction, $\ln[\ce{A}]_t = \ln[\ce{A}]_0 - kt$.
Here the titre $V_t \propto [\ce{I2}]_t$, and at $t=\infty$, $V_\infty = 0$ (all the $\ce{I2}$ has reacted). Hence $[\ce{I2}]_t \propto V_t$ and the integrated rate law becomes $\ln V_t = \ln V_0 - kt$.
Calculate $\ln(V_t)$:

\begin{center}
\begin{tabular}{|c|ccccc|}\hline
$t$ (min) & 0 & 5 & 10 & 15 & 20 \\ \hline
$V_t$ ($\mathrm{cm^3}$) & 25.0 & 18.4 & 13.5 & 9.9 & 7.3 \\ \hline
$\ln(V_t)$ & 3.219 & 2.912 & 2.603 & 2.293 & 1.988 \\ \hline
\end{tabular}
\end{center}

A plot of $\ln(V_t)$ against $t$ is a straight line of negative slope, confirming first order. Alternatively, check the constancy of $k = \dfrac{1}{t}\ln\dfrac{V_0}{V_t}$:
\begin{align*}
t=5: &\quad k = \tfrac{1}{5}\ln\tfrac{25.0}{18.4} = 0.0614\ \mathrm{min^{-1}} \\
t=10: &\quad k = \tfrac{1}{10}\ln\tfrac{25.0}{13.5} = 0.0616\ \mathrm{min^{-1}} \\
t=15: &\quad k = \tfrac{1}{15}\ln\tfrac{25.0}{9.9} = 0.0618\ \mathrm{min^{-1}} \\
t=20: &\quad k = \tfrac{1}{20}\ln\tfrac{25.0}{7.3} = 0.0616\ \mathrm{min^{-1}}
\end{align*}
$k$ is constant ($\approx 0.0616\ \mathrm{min^{-1}}$). \textbf{The reaction is first order with respect to $\ce{I2}$.}

\textbf{(b) Rate constant in $\mathrm{s^{-1}}$:}
\[ k = 0.0616\ \mathrm{min^{-1}} = \frac{0.0616}{60}\ \mathrm{s^{-1}} = \boxed{1.03 \times 10^{-3}\ \mathrm{s^{-1}}}. \]

\textbf{(c) Half-life:} For first order, $t_{1/2} = \dfrac{\ln 2}{k} = \dfrac{0.693}{0.0616\ \mathrm{min^{-1}}} = \boxed{11.3\ \mathrm{min}}$ ($\approx 675\ \mathrm{s}$).

\textbf{(d) Reason:} $\ce{CH3COCH3}$ and $\ce{H+}$ are in \textbf{large excess} compared with $\ce{I2}$, so their concentrations remain effectively constant during the reaction. The observed rate law is $\text{Rate} = k'[\ce{I2}]$ where $k' = k[\ce{CH3COCH3}]^x[\ce{H+}]^y$ (a \textbf{pseudo-first-order} constant). To find $x$ and $y$, the experiment must be repeated varying the initial $[\ce{CH3COCH3}]$ and $[\ce{H+}]$ separately while keeping the other reagents in excess.
\end{exoresolu}

\begin{methode}[Determining Order of Reaction: Initial Rates Method]
\begin{enumerate}
\item \textbf{Write the rate law:} $\text{Rate} = k[\ce{A}]^x[\ce{B}]^y$. Overall order $= x+y$.
\item \textbf{Select experiments:} Compare two experiments in which \textbf{only one reactant concentration changes} (all others constant).
\item \textbf{Calculate the ratio:} $\dfrac{\text{Rate}_2}{\text{Rate}_1} = \left(\dfrac{[\ce{A}]_2}{[\ce{A}]_1}\right)^x$.
\item \textbf{Solve for the order $x$:} Take logarithms: $x = \dfrac{\log(\text{Rate}_2/\text{Rate}_1)}{\log([\ce{A}]_2/[\ce{A}]_1)}$. Usually the ratios are simple (doubling the concentration multiplies the rate by 1, 2, 4 or 8), giving integer orders 0, 1, 2 or 3.
\item \textbf{Repeat for the other reactants.}
\item \textbf{Calculate $k$:} Substitute the orders and the data from any one experiment into the rate law. \textbf{State the units of $k$} according to the overall order:
    \begin{itemize}
    \item Order 0: $\mathrm{mol\,dm^{-3}\,s^{-1}}$
    \item Order 1: $\mathrm{s^{-1}}$
    \item Order 2: $\mathrm{mol^{-1}\,dm^3\,s^{-1}}$
    \item Order 3: $\mathrm{mol^{-2}\,dm^6\,s^{-1}}$
    \end{itemize}
\item \textbf{Check:} Use $k$ to predict the rate for another experiment and compare with the measured value.
\end{enumerate}
\end{methode}

\begin{exoresolu}[Initial Rates Analysis: Reaction between NO and O2]
The reaction $\ce{2NO(g) + O2(g) -> 2NO2(g)}$ is studied at constant temperature. Initial rates are measured:

\begin{center}
\begin{tabular}{|c|cc|c|}\hline
Exp. & $[\ce{NO}]_0$ ($\mathrm{mol\,dm^{-3}}$) & $[\ce{O2}]_0$ ($\mathrm{mol\,dm^{-3}}$) & Initial Rate ($\mathrm{mol\,dm^{-3}\,s^{-1}}$) \\ \hline
1 & 0.010 & 0.010 & $2.5 \times 10^{-5}$ \\
2 & 0.020 & 0.010 & $1.0 \times 10^{-4}$ \\
3 & 0.010 & 0.020 & $5.0 \times 10^{-5}$ \\
4 & 0.020 & 0.020 & $2.0 \times 10^{-4}$ \\ \hline
\end{tabular}
\end{center}

\textbf{(a)} Determine the order with respect to $\ce{NO}$ and $\ce{O2}$. Write the rate equation. \\
\textbf{(b)} Calculate the rate constant $k$ with units. \\
\textbf{(c)} Predict the initial rate for Experiment 5: $[\ce{NO}]_0 = 0.030\ \mathrm{mol\,dm^{-3}}$, $[\ce{O2}]_0 = 0.015\ \mathrm{mol\,dm^{-3}}$.
\tcblower
\textbf{Detailed Solution}

\textbf{(a) Order with respect to $\ce{NO}$ (compare Experiments 1 and 2):} $[\ce{O2}]$ is constant ($0.010$); $[\ce{NO}]$ doubles ($0.010 \to 0.020$).
Rate ratio: $\dfrac{1.0 \times 10^{-4}}{2.5 \times 10^{-5}} = 4$.
$\left(\dfrac{0.020}{0.010}\right)^x = 2^x = 4 \implies x = 2$. \textbf{Second order with respect to $\ce{NO}$.}

\textbf{Order with respect to $\ce{O2}$ (compare Experiments 1 and 3):} $[\ce{NO}]$ is constant ($0.010$); $[\ce{O2}]$ doubles ($0.010 \to 0.020$).
Rate ratio: $\dfrac{5.0 \times 10^{-5}}{2.5 \times 10^{-5}} = 2$.
$\left(\dfrac{0.020}{0.010}\right)^y = 2^y = 2 \implies y = 1$. \textbf{First order with respect to $\ce{O2}$.}

\textbf{Rate equation:} $\text{Rate} = k[\ce{NO}]^2[\ce{O2}]$. Overall order $= 2 + 1 = 3$.

\textbf{(b) Rate constant $k$:} Use Experiment 1:
\[ 2.5 \times 10^{-5} = k (0.010)^2 (0.010) = k \times 1.0 \times 10^{-6} \implies k = \frac{2.5 \times 10^{-5}}{1.0 \times 10^{-6}} = 25. \]
Units: $\dfrac{\mathrm{mol\,dm^{-3}\,s^{-1}}}{(\mathrm{mol\,dm^{-3}})^3} = \mathrm{mol^{-2}\,dm^6\,s^{-1}}$.
\[ \boxed{k = 25\ \mathrm{mol^{-2}\,dm^6\,s^{-1}}}. \]
\textbf{Check with Experiment 4:} predicted rate $= 25 \times (0.020)^2 \times (0.020) = 25 \times 8.0 \times 10^{-6} = 2.0 \times 10^{-4}\ \mathrm{mol\,dm^{-3}\,s^{-1}}$. This matches the measured value, confirming the rate law.

\textbf{(c) Prediction for Experiment 5:}
\[ \text{Rate} = 25 \times (0.030)^2 \times (0.015) = 25 \times 9.0 \times 10^{-4} \times 1.5 \times 10^{-2} = \boxed{3.4 \times 10^{-4}\ \mathrm{mol\,dm^{-3}\,s^{-1}}}. \]
\end{exoresolu}

\begin{methode}[Graphical Determination of Order: Integrated Rate Laws and Half-Life]
\begin{enumerate}
\item \textbf{Identify the measured variable:} concentration $[\ce{A}]$, titre $V_t$, pressure $P_t$, absorbance $A_t$, conductivity $\kappa_t$. Ensure it is \textbf{proportional to the concentration} of one species.
\item \textbf{Test for zero order:} Plot $[\ce{A}]_t$ against $t$. A straight line $\implies$ zero order. Slope $= -k$. $t_{1/2} = \dfrac{[\ce{A}]_0}{2k}$ (depends on $[\ce{A}]_0$).
\item \textbf{Test for first order:} Plot $\ln[\ce{A}]_t$ (or $\ln V_t$, or $\ln(V_t - V_\infty)$) against $t$. A straight line $\implies$ first order. Slope $= -k$. $t_{1/2} = \dfrac{\ln 2}{k}$ (\textbf{constant}, independent of $[\ce{A}]_0$).
\item \textbf{Test for second order:} Plot $\dfrac{1}{[\ce{A}]_t}$ against $t$. A straight line $\implies$ second order. Slope $= +k$. $t_{1/2} = \dfrac{1}{k[\ce{A}]_0}$ (halves when $[\ce{A}]_0$ doubles).
\item \textbf{Half-life method (experimental):} Measure $t_{1/2}$ for different initial concentrations $[\ce{A}]_0$:
    \begin{itemize}
    \item $t_{1/2}$ constant $\Rightarrow$ first order;
    \item $t_{1/2} \propto [\ce{A}]_0$ $\Rightarrow$ zero order;
    \item $t_{1/2} \propto 1/[\ce{A}]_0$ $\Rightarrow$ second order.
    \end{itemize}
\item \textbf{Choose the plot giving the best straight line.}
\end{enumerate}
\end{methode}

\begin{exoresolu}[Graphical Method: Decomposition of $\ce{N2O5}$]
The gas-phase decomposition $\ce{2N2O5(g) -> 4NO2(g) + O2(g)}$ is followed by measuring the total pressure $P_{\text{tot}}$ in a constant-volume vessel at $338\ \mathrm{K}$. Initial pressure $P_0 = 100\ \mathrm{kPa}$ (pure $\ce{N2O5}$).

\begin{center}
\begin{tabular}{|c|cccccc|}\hline
Time (s) & 0 & 200 & 400 & 600 & 800 & 1000 \\ \hline
$P_{\text{tot}}$ (kPa) & 100 & 118 & 132 & 143 & 152 & 159 \\ \hline
\end{tabular}
\end{center}

\textbf{(a)} Derive the relationship between $P_{\ce{N2O5}}$ and $P_{\text{tot}}$. \\
\textbf{(b)} Show that the reaction is first order with respect to $\ce{N2O5}$ and determine $k$. \\
\textbf{(c)} Calculate the half-life. \\
\textbf{(d)} What is the partial pressure of $\ce{O2}$ at $t=600\ \mathrm{s}$?
\tcblower
\textbf{Detailed Solution}

\textbf{(a) Stoichiometry and pressure:} Let $x$ = partial pressure of $\ce{N2O5}$ that has decomposed at time $t$. From the equation $\ce{2N2O5 -> 4NO2 + O2}$, 2 mol of $\ce{N2O5}$ give 4 mol of $\ce{NO2}$ and 1 mol of $\ce{O2}$. Hence, in pressure units:
\begin{center}
\begin{tabular}{|l|ccc|}\hline
 & $\ce{N2O5}$ & $\ce{NO2}$ & $\ce{O2}$ \\ \hline
Initial & $100$ & $0$ & $0$ \\
Change & $-x$ & $+2x$ & $+\tfrac{x}{2}$ \\
At time $t$ & $100 - x$ & $2x$ & $\tfrac{x}{2}$ \\ \hline
\end{tabular}
\end{center}
\[ P_{\text{tot}} = (100 - x) + 2x + \frac{x}{2} = 100 + \frac{3x}{2} \implies x = \frac{2}{3}\left(P_{\text{tot}} - 100\right). \]
Therefore $\boxed{P_{\ce{N2O5}} = 100 - \dfrac{2}{3}\left(P_{\text{tot}} - 100\right)}$.

\textbf{(b) First-order test:} Calculate $P_{\ce{N2O5}}$ and $\ln(P_{\ce{N2O5}})$:

\begin{center}
\begin{tabular}{|c|cccccc|}\hline
$t$ (s) & 0 & 200 & 400 & 600 & 800 & 1000 \\ \hline
$P_{\text{tot}}$ (kPa) & 100 & 118 & 132 & 143 & 152 & 159 \\ \hline
$P_{\ce{N2O5}}$ (kPa) & 100 & 88.0 & 78.7 & 71.3 & 65.3 & 60.7 \\ \hline
$\ln(P_{\ce{N2O5}})$ & 4.605 & 4.477 & 4.365 & 4.267 & 4.179 & 4.106 \\ \hline
\end{tabular}
\end{center}
(Example: at $t=200\ \mathrm{s}$, $P_{\ce{N2O5}} = 100 - \frac{2}{3}(118-100) = 100 - 12.0 = 88.0\ \mathrm{kPa}$.)

A plot of $\ln(P_{\ce{N2O5}})$ against $t$ is a straight line, confirming \textbf{first order}. Slope (using $t=0$ and $t=1000\ \mathrm{s}$):
\[ \text{Slope} = \frac{4.106 - 4.605}{1000 - 0} = -4.99 \times 10^{-4}\ \mathrm{s^{-1}} \implies k = -\text{slope} = \boxed{5.0 \times 10^{-4}\ \mathrm{s^{-1}}}. \]

\textbf{(c) Half-life:}
\[ t_{1/2} = \frac{\ln 2}{k} = \frac{0.693}{5.0 \times 10^{-4}} = \boxed{1.39 \times 10^{3}\ \mathrm{s}\ (\approx 23\ \mathrm{min})}. \]

\textbf{(d) Partial pressure of $\ce{O2}$ at $t=600\ \mathrm{s}$:}
\[ P_{\ce{O2}} = \frac{x}{2} = \frac{1}{2}\times\frac{2}{3}\left(P_{\text{tot}} - 100\right) = \frac{P_{\text{tot}} - 100}{3} = \frac{143 - 100}{3} = \boxed{14.3\ \mathrm{kPa}}. \]
\end{exoresolu}

\begin{methode}[Arrhenius Equation: Determining $E_a$ and $A$]
\begin{enumerate}
\item \textbf{Arrhenius equation:} $k = A e^{-E_a/RT}$, where $A$ is the pre-exponential (frequency) factor, $E_a$ the activation energy, $R = 8.314\ \mathrm{J\,mol^{-1}\,K^{-1}}$, $T$ the absolute temperature.
\item \textbf{Linear form:} $\ln k = \ln A - \dfrac{E_a}{R}\cdot\dfrac{1}{T}$.
    \begin{itemize}
    \item Plot $\ln k$ ($y$-axis) against $\dfrac{1}{T}$ ($x$-axis, in $\mathrm{K^{-1}}$).
    \item A straight line confirms Arrhenius behaviour.
    \item Slope $= -\dfrac{E_a}{R} \implies E_a = -\text{slope} \times R$.
    \item Intercept $= \ln A \implies A = e^{\text{intercept}}$.
    \end{itemize}
\item \textbf{Two-point form (only two temperatures available):} $\ln \dfrac{k_2}{k_1} = \dfrac{E_a}{R}\left(\dfrac{1}{T_1} - \dfrac{1}{T_2}\right)$.
\item \textbf{Units:} $E_a$ in $\mathrm{J\,mol^{-1}}$ (convert to $\mathrm{kJ\,mol^{-1}}$ at the end); $T$ in kelvin; $A$ has the same units as $k$.
\end{enumerate}
\end{methode}

\begin{attention}[Common Errors with the Arrhenius Equation]
Using $T$ in $^\circ\mathrm{C}$ instead of kelvin; forgetting to convert $E_a$ from $\mathrm{J}$ to $\mathrm{kJ}$; confusing $\log_{10}$ with $\ln$ (if $\log_{10} k$ is plotted, the slope is $-E_a/2.303R$); quoting $A$ without units.
\end{attention}

\begin{exoresolu}[Activation Energy from Temperature Dependence]
The rate constant $k$ for the decomposition of $\ce{N2O5}$ was measured at different temperatures:

\begin{center}
\begin{tabular}{|c|ccccc|}\hline
$T$ ($^\circ\mathrm{C}$) & 25 & 35 & 45 & 55 & 65 \\ \hline
$k$ ($\mathrm{s^{-1}}$) & $1.7 \times 10^{-5}$ & $6.5 \times 10^{-5}$ & $2.4 \times 10^{-4}$ & $8.2 \times 10^{-4}$ & $2.7 \times 10^{-3}$ \\ \hline
\end{tabular}
\end{center}

\textbf{(a)} Plot the Arrhenius graph and determine $E_a$ (in $\mathrm{kJ\,mol^{-1}}$) and $A$. \\
\textbf{(b)} Calculate $k$ at $50^\circ\mathrm{C}$. \\
\textbf{(c)} Explain the physical significance of $A$ and state why $E_a$ is taken as independent of temperature in this model.
\tcblower
\textbf{Detailed Solution}

\textbf{(a) Data processing:} Convert $T$ to kelvin ($T/\mathrm{K} = T/^\circ\mathrm{C} + 273$). Calculate $1/T$ and $\ln k$:

\begin{center}
\begin{tabular}{|c|ccccc|}\hline
$T$ (K) & 298 & 308 & 318 & 328 & 338 \\ \hline
$10^3/T$ ($\mathrm{K^{-1}}$) & 3.356 & 3.247 & 3.145 & 3.049 & 2.959 \\ \hline
$k$ ($\mathrm{s^{-1}}$) & $1.7\times 10^{-5}$ & $6.5\times 10^{-5}$ & $2.4\times 10^{-4}$ & $8.2\times 10^{-4}$ & $2.7\times 10^{-3}$ \\ \hline
$\ln k$ & -10.98 & -9.64 & -8.33 & -7.11 & -5.91 \\ \hline
\end{tabular}
\end{center}

Plot $\ln k$ against $1/T$ and draw the best straight line. Select two points on the line, far apart:
Point 1: $(3.356 \times 10^{-3},\ -10.98)$; Point 2: $(2.959 \times 10^{-3},\ -5.91)$.
\[ \text{Slope} = \frac{-5.91 - (-10.98)}{(2.959 - 3.356)\times 10^{-3}} = \frac{5.07}{-3.97 \times 10^{-4}} = -1.28 \times 10^{4}\ \mathrm{K}. \]
\[ E_a = -\text{slope} \times R = 1.28 \times 10^{4} \times 8.314 = 1.06 \times 10^{5}\ \mathrm{J\,mol^{-1}} = \boxed{106\ \mathrm{kJ\,mol^{-1}}}. \]
Intercept: using $c = y - mx$ at Point 1:
\[ \ln A = -10.98 - (-1.28 \times 10^{4} \times 3.356 \times 10^{-3}) = -10.98 + 42.96 = 31.98 \implies A = e^{31.98} \approx \boxed{7 \times 10^{13}\ \mathrm{s^{-1}}}. \]
($A$ has the same units as $k$, i.e. $\mathrm{s^{-1}}$ for a first-order reaction.)

\textbf{(b) $k$ at $50^\circ\mathrm{C}$ ($323\ \mathrm{K}$):} $1/T = 3.096 \times 10^{-3}\ \mathrm{K^{-1}}$.
\[ \ln k = 31.98 - (1.28 \times 10^{4})(3.096 \times 10^{-3}) = 31.98 - 39.63 = -7.65 \implies k = e^{-7.65} = \boxed{4.8 \times 10^{-4}\ \mathrm{s^{-1}}}. \]
(Check: this lies between $2.4\times 10^{-4}$ at $45^\circ\mathrm{C}$ and $8.2\times 10^{-4}$ at $55^\circ\mathrm{C}$ — reasonable.)

\textbf{(c) Significance:} $A$ (the pre-exponential factor) is related to the \textbf{collision frequency} and the \textbf{orientation (steric) requirement}: $A = pZ$, where $Z$ is the collision frequency and $p$ the steric factor. It represents the value $k$ would approach at infinitely high temperature, when effectively every collision has enough energy to react. $E_a$ is the \textbf{minimum energy} a colliding pair must possess for reaction to occur — the height of the energy barrier between reactants and the transition state. In the simple Arrhenius model this barrier height is a property of the potential energy surface and is therefore treated as constant (independent of $T$); the temperature dependence of $k$ comes almost entirely from the exponential term $e^{-E_a/RT}$.
\end{exoresolu}

\begin{methode}[Collision Theory and Transition State Theory: Key Ideas]
\begin{enumerate}
\item \textbf{Collision theory:} Rate $= Z \times f \times p$, where
    \begin{itemize}
    \item $Z$: collision frequency (increases with concentration and with $\sqrt{T}$);
    \item $f = e^{-E_a/RT}$: fraction of collisions with energy $\ge E_a$ (from the Maxwell--Boltzmann distribution);
    \item $p$: steric factor (probability that the collision has the correct orientation).
    \end{itemize}
\item \textbf{Maxwell--Boltzmann distribution:} the curve of molecular energies shows that only a small fraction of molecules have $E \ge E_a$; raising the temperature flattens and shifts the curve, greatly increasing this fraction — hence the strong temperature dependence of rates.
\item \textbf{Transition state theory (TST):} Reactants pass through a high-energy \textbf{activated complex} (transition state) at the top of the energy barrier: $\ensuremath{\mathrm{A + B \rightleftharpoons  AB^{\ddagger} \rightarrow  products}}$. The rate depends on the concentration of the activated complex and the frequency with which it decomposes into products.
\item \textbf{Similarities:} both theories require a minimum energy ($E_a$) for reaction; both predict the Arrhenius temperature dependence; both involve an activated, high-energy configuration of atoms.
\item \textbf{Differences:}
    \begin{itemize}
    \item Collision theory treats molecules as hard spheres and applies best to simple gas-phase reactions; the steric factor $p$ is introduced empirically.
    \item TST considers the structure and energy of the activated complex explicitly; it applies to reactions in solution as well as gases, and accounts for orientation and entropy effects through the nature of the transition state rather than an empirical factor.
    \end{itemize}
\end{enumerate}
\end{methode}

\begin{exoresolu}[Collision Theory: Fraction of Effective Collisions]
For a bimolecular gas reaction at $300\ \mathrm{K}$, $E_a = 50.0\ \mathrm{kJ\,mol^{-1}}$. The measured pre-exponential factor is $A = 2.0 \times 10^{11}\ \mathrm{mol^{-1}\,dm^3\,s^{-1}}$, and the collision frequency factor $Z$ (for unit concentrations) is calculated as $4.0 \times 10^{11}\ \mathrm{mol^{-1}\,dm^3\,s^{-1}}$.

\textbf{(a)} Calculate the fraction of collisions with energy $\ge E_a$ at $300\ \mathrm{K}$. \\
\textbf{(b)} Calculate the steric factor $p$ and comment on its value. \\
\textbf{(c)} Explain, in terms of the Maxwell--Boltzmann distribution, why the rate roughly doubles when the temperature rises by about $10\ \mathrm{K}$ near room temperature. \\
\textbf{(d)} Why might the steric factor $p$ be less than 1?
\tcblower
\textbf{Detailed Solution}

\textbf{(a) Fraction $f$ (Maxwell--Boltzmann):}
\[ f = e^{-E_a/RT} = \exp\left(-\frac{50.0 \times 10^{3}}{8.314 \times 300}\right) = \exp(-20.05) = \boxed{2.0 \times 10^{-9}}. \]
Only about 2 collisions in every $10^{9}$ have sufficient energy to react.

\textbf{(b) Steric factor $p$:} Collision theory gives $A = pZ$, so
\[ p = \frac{A}{Z} = \frac{2.0 \times 10^{11}}{4.0 \times 10^{11}} = \boxed{0.50}. \]
\textbf{Comment:} $p = 0.5$ is relatively high (for complex molecules $p$ can be $10^{-1}$ to $10^{-4}$). It indicates a modest orientation requirement, as expected for small, fairly symmetrical reactants. A value below 1 confirms that not every sufficiently energetic collision leads to reaction.

\textbf{(c) Temperature effect:} Raising the temperature broadens the Maxwell--Boltzmann distribution and shifts it to higher energies. Because $f = e^{-E_a/RT}$ depends \textbf{exponentially} on $T$, even a small rise in $T$ produces a large increase in the area under the curve beyond $E_a$. Near room temperature, for typical values of $E_a$ ($\approx 50\ \mathrm{kJ\,mol^{-1}}$), a $10\ \mathrm{K}$ rise approximately doubles $f$, and hence roughly doubles the rate. (The collision frequency $Z$ increases only slightly, as $\sqrt{T}$, and is not the main cause.)

\textbf{(d) Why $p < 1$:} Reacting molecules must collide with the \textbf{correct orientation} for the right bonds to break and form. For example, in $\ce{NO + O3 -> NO2 + O2}$, the nitrogen end of $\ce{NO}$ must strike a terminal oxygen of $\ce{O3}$. Collisions with the wrong geometry, or with energy locked in the wrong vibrational modes, do not lead to products even if $E \ge E_a$.
\end{exoresolu}

\begin{methode}[Catalysis: Homogeneous and Heterogeneous]
\begin{enumerate}
\item \textbf{Definition:} A catalyst increases the rate of a reaction by providing an alternative pathway of \textbf{lower activation energy}. It is not consumed overall, and it does not change $\Delta H$ or the position of equilibrium $K_c$ — it only speeds up the attainment of equilibrium.
\item \textbf{Homogeneous catalysis:} catalyst in the \textbf{same phase} as the reactants (usually aqueous or gas).
    \begin{itemize}
    \item Mechanism: the catalyst reacts to form an intermediate, then is regenerated.
    \item Examples: $\ce{H+}$ in ester hydrolysis; $\ce{Fe^{2+}}$ (or $\ce{Fe^{3+}}$) catalysing $\ce{S2O8^2- + 2I- -> 2SO4^2- + I2}$; $\ce{NO}$ in the lead-chamber process for $\ce{2SO2 + O2 -> 2SO3}$.
    \item The catalyst concentration appears in the rate equation, e.g. $\text{Rate} = k[\ce{H+}][\text{ester}]$.
    \end{itemize}
\item \textbf{Heterogeneous catalysis:} catalyst in a \textbf{different phase} (usually a solid catalyst with gaseous or liquid reactants).
    \begin{itemize}
    \item Steps: \textbf{adsorption} of reactants onto active sites $\to$ reaction on the surface $\to$ \textbf{desorption} of products.
    \item The rate increases with \textbf{surface area} (powder more effective than lumps) and with the number of active sites.
    \item \textbf{Poisoning:} impurities (e.g. $\ce{CO}$, $\ce{H2S}$, sulfur compounds, lead compounds) adsorb strongly and block active sites, destroying catalytic activity.
    \item Examples: $\ce{Fe}$ in the Haber process; $\ce{V2O5}$ in the Contact process; $\ce{Ni}$ or $\ce{Pt}$ in catalytic hydrogenation; $\ce{Pt}/\ce{Rh}$ in catalytic converters.
    \end{itemize}
\item \textbf{Graphical analysis:} on an energy profile diagram, the catalysed route has a lower peak (lower $E_a$) but the same reactant and product energy levels (same $\Delta H$). On an Arrhenius plot, the catalysed reaction gives a line of smaller slope magnitude.
\end{enumerate}
\end{methode}

\begin{exoresolu}[Catalysis: Homogeneous Kinetics and Heterogeneous Poisoning]
\textbf{(a)} The oxidation of $\ce{I-}$ by $\ce{S2O8^2-}$ is catalysed by $\ce{Fe^{2+}}$:
$\ce{S2O8^2- + 2I- -> 2SO4^2- + I2}$ (slow when uncatalysed).
Catalysed mechanism:
Step 1: $\ce{S2O8^2- + 2Fe^{2+} -> 2SO4^2- + 2Fe^{3+}}$ (slow)
Step 2: $\ce{2Fe^{3+} + 2I- -> 2Fe^{2+} + I2}$ (fast)

\textbf{(i)} Identify the catalyst and the intermediate. \\
\textbf{(ii)} Write the overall equation and verify that the catalyst is regenerated. \\
\textbf{(iii)} The rate law is $\text{Rate} = k[\ce{S2O8^2-}][\ce{Fe^{2+}}]$. Explain why the rate depends on $[\ce{Fe^{2+}}]$ but not on $[\ce{I-}]$.

\textbf{(b)} In the Haber process, $\ce{N2 + 3H2 <=> 2NH3}$ ($\ce{Fe}$ catalyst), carbon monoxide acts as a poison.
\textbf{(i)} Explain the mechanism of poisoning. \\
\textbf{(ii)} Sketch the energy profile (energy against reaction coordinate) for the catalysed and uncatalysed reactions, labelling $E_a(\text{cat})$, $E_a(\text{uncat})$ and $\Delta H$.
\tcblower
\textbf{Detailed Solution}

\textbf{(a)(i)} \textbf{Catalyst:} $\ce{Fe^{2+}}$ — it is consumed in Step 1 and regenerated in Step 2. \textbf{Intermediate:} $\ce{Fe^{3+}}$ — it is formed in Step 1 and consumed in Step 2.

\textbf{(a)(ii) Overall equation:} Add Steps 1 and 2:
\[ \ce{S2O8^2- + 2Fe^{2+} + 2Fe^{3+} + 2I- -> 2SO4^2- + 2Fe^{3+} + 2Fe^{2+} + I2} \]
Cancelling $\ce{Fe^{2+}}$ and $\ce{Fe^{3+}}$ from both sides:
\[ \ce{S2O8^2- + 2I- -> 2SO4^2- + I2}. \]
The catalyst $\ce{Fe^{2+}}$ is regenerated. (The uncatalysed reaction is slow because it brings two negative ions, $\ce{S2O8^2-}$ and $\ce{I-}$, together; the catalysed route replaces this by steps between ions of opposite charge, which have lower $E_a$.)

\textbf{(a)(iii) Rate law interpretation:} Step 1 is the slow (rate-determining) step, and it involves $\ce{S2O8^2-}$ and $\ce{Fe^{2+}}$ only. Hence $\text{Rate} = k[\ce{S2O8^2-}][\ce{Fe^{2+}}]$. $\ce{I-}$ reacts in the fast step \textbf{after} the rate-determining step, so changing $[\ce{I-}]$ does not affect the overall rate. Because the catalyst is a reactant in the rate-determining step, its concentration appears in the rate law.

\textbf{(b)(i) Poisoning mechanism:} $\ce{CO}$ chemisorbs \textbf{strongly and essentially irreversibly} onto the active sites of the iron surface — the very sites where $\ce{N2}$ and $\ce{H2}$ must adsorb and dissociate. The blocked sites are no longer available, the effective surface area for the reaction falls, and the rate drops sharply. Because $\ce{CO}$ binds more strongly than $\ce{N2}$ or $\ce{H2}$, even trace amounts cause serious loss of activity; the synthesis gas must therefore be carefully purified.

\textbf{(b)(ii) Energy profile:}
\begin{center}
\begin{tikzpicture}[scale=0.7]
\draw[->] (0,0) -- (10,0) node[right] {Reaction coordinate};
\draw[->] (0,0) -- (0,6) node[above] {Energy};
\draw[poporange, thick] (1,3) .. controls (3,5.5) and (5,5.5) .. (7,1);
\draw[popblue, thick] (1,3) .. controls (3,4) and (5,4) .. (7,1);
\draw[dashed] (1,3) -- (9,3);
\draw[dashed] (7,1) -- (9,1);
\draw[<->, poporange] (4,3) -- (4,5.35) node[midway, right] {$E_a(\text{uncat})$};
\draw[<->, popblue] (2,3) -- (2,3.85) node[midway, right] {$E_a(\text{cat})$};
\draw[<->, popdark] (8.5,1) -- (8.5,3) node[midway, right] {$\Delta H$};
\node at (1,3.4) {Reactants};
\node at (7,0.5) {Products};
\node[poporange] at (6.3,5.6) {Uncatalysed};
\node[popblue] at (6.3,4.1) {Catalysed};
\end{tikzpicture}
\legende{Energy profile for an exothermic reaction with and without a catalyst. The catalyst lowers $E_a$ but leaves $\Delta H$ unchanged.}
\end{center}
\textbf{Key features:} reactants and products are at the same energy levels on both curves, so $\Delta H$ is unchanged; the catalysed pathway has a lower transition-state peak, hence a lower $E_a$ and a faster reaction. (For heterogeneous catalysis the profile is often drawn with two smaller humps separated by a shallow minimum corresponding to the adsorbed intermediate.)
\end{exoresolu}

\begin{methode}[Physical Methods: Colorimetry, Conductimetry, Gas Syringe and Pressure]
\begin{enumerate}
\item \textbf{Colorimetry / spectrophotometry:} based on the Beer--Lambert law, $A = \varepsilon\, b\, c$ (absorbance proportional to concentration).
    \begin{itemize}
    \item Use when \textbf{one species is coloured} (e.g. $\ce{I2}$, $\ce{MnO4-}$, $\ce{Cr2O7^2-}$, $\ce{Br2}$).
    \item Measure absorbance $A_t$ at the wavelength of maximum absorption; $A_t \propto c_t$.
    \item Plot $A_t$ against $t$, or $\ln A_t$ against $t$ to test for first order.
    \item Advantages: continuous, rapid, no sampling, very sensitive.
    \end{itemize}
\item \textbf{Conductimetry:} measures the conductivity $\kappa$ of the mixture, which depends on the concentration and mobility of all ions present.
    \begin{itemize}
    \item Use for reactions in which the \textbf{number or nature of ionic species changes}, e.g. $\ce{CH3COOC2H5 + OH- -> CH3COO- + C2H5OH}$ (conductivity falls because the highly mobile $\ce{OH-}$ is replaced by the less mobile $\ce{CH3COO-}$).
    \item Plot $\kappa_t$ against $t$; for first order, $\ln(\kappa_t - \kappa_\infty)$ against $t$ is linear.
    \item Advantages: continuous, no sampling. Disadvantages: non-specific (all ions contribute) and very sensitive to temperature.
    \end{itemize}
\item \textbf{Gas syringe / pressure measurement:} for reactions involving a gas.
    \begin{itemize}
    \item Gas evolved (e.g. $\ce{CaCO3 + 2HCl}$, $\ce{2H2O2 -> 2H2O + O2}$): collect the gas in a gas syringe; $V_{\text{gas}} \propto n_{\text{gas}}$.
    \item Change in moles of gas at constant $V$ and $T$ (e.g. $\ce{2N2O5 -> 4NO2 + O2}$): $P \propto n_{\text{total}}$; use the stoichiometry to relate $P_{\text{tot}}$ to the reactant pressure.
    \end{itemize}
\item \textbf{Dilatometry:} follows the small volume change of a liquid reaction mixture in a dilatometer (a flask with a capillary stem); the height of the meniscus is proportional to the extent of reaction. Mentioned for completeness.
\end{enumerate}
\end{methode}

\begin{exoresolu}[Conductimetric Study: Hydrolysis of Ethyl Ethanoate with NaOH]
$\ce{CH3COOC2H5(aq) + OH-(aq) -> CH3COO-(aq) + C2H5OH(aq)}$, followed by conductivity measurements.
Molar conductivities at the working temperature: $\lambda_{\ce{OH-}} = 198$, $\lambda_{\ce{CH3COO-}} = 41$, $\lambda_{\ce{Na+}} = 50\ \mathrm{S\,cm^2\,mol^{-1}}$; the ester and ethanol are non-ionic.
Initial $[\ce{OH-}]_0 = [\text{ester}]_0 = 0.010\ \mathrm{mol\,dm^{-3}}$; $\ce{Na+}$ is a spectator.

\textbf{(a)} Explain why the conductivity decreases during the reaction. \\
\textbf{(b)} Derive an expression for the conductivity $\kappa_t$ at time $t$ in terms of $x$, the concentration that has reacted, and show that $\kappa_t - \kappa_\infty \propto [\ce{OH-}]_t$. \\
\textbf{(c)} A student claims that, because the reaction is second order, a plot of $\ln(\kappa_t - \kappa_\infty)$ against $t$ should be linear. Discuss this claim and state the correct linear plot. \\
\textbf{(d)} If the reaction were carried out with $\ce{OH-}$ in large excess, what order would be observed? Write the observed rate law and the appropriate linear plot.
\tcblower
\textbf{Detailed Solution}

\textbf{(a) Conductivity change:} Conductivity depends on the ions present and their mobilities. $\ce{Na+}$ is a spectator and its concentration is constant. During the reaction, highly mobile $\ce{OH-}$ ions ($\lambda = 198$) are consumed and replaced by much less mobile $\ce{CH3COO-}$ ions ($\lambda = 41$). The total conductivity therefore \textbf{decreases} as the reaction proceeds.

\textbf{(b) Expression for $\kappa_t$:} Let $x$ = concentration reacted ($\mathrm{mol\,dm^{-3}}$). Then
$[\ce{OH-}]_t = 0.010 - x$; $[\ce{CH3COO-}]_t = x$; $[\ce{Na+}]_t = 0.010$ (constant).
\[ \kappa_t = 10^{-3}\left[ (0.010)(50) + (0.010 - x)(198) + (x)(41) \right] = 10^{-3}\left[ 2.48 - 157x \right]\ \mathrm{S\,cm^{-1}}. \]
(The factor $10^{-3}$ converts $\mathrm{mol\,dm^{-3}}$ to $\mathrm{mol\,cm^{-3}}$.)
At $t=0$ ($x=0$): $\kappa_0 = 2.48 \times 10^{-3}\ \mathrm{S\,cm^{-1}}$.
At $t=\infty$ ($x=0.010$): $\kappa_\infty = 10^{-3}[2.48 - 1.57] = 0.91 \times 10^{-3}\ \mathrm{S\,cm^{-1}}$.
Then
\[ \kappa_t - \kappa_\infty = 10^{-3}\left[ (2.48 - 157x) - 0.91 \right] = 157 \times 10^{-3}(0.010 - x) = 157 \times 10^{-3}\,[\ce{OH-}]_t. \]
\textbf{Hence $\kappa_t - \kappa_\infty$ is directly proportional to $[\ce{OH-}]_t$.}

\textbf{(c) Discussion of the claim:} The claim is \textbf{incorrect}. With equal initial concentrations, $[\ce{OH-}] = [\text{ester}]$ throughout, so
\[ -\frac{d[\ce{OH-}]}{dt} = k[\ce{OH-}][\text{ester}] = k[\ce{OH-}]^2, \]
a \textbf{second-order} rate law, whose integrated form is
\[ \frac{1}{[\ce{OH-}]_t} = \frac{1}{[\ce{OH-}]_0} + kt. \]
Since $[\ce{OH-}]_t \propto \kappa_t - \kappa_\infty$, the correct linear plot is
\[ \frac{1}{\kappa_t - \kappa_\infty} \text{ against } t \quad (\text{straight line, slope} \propto k). \]
A plot of $\ln(\kappa_t - \kappa_\infty)$ against $t$ is the test for \textbf{first} order and would \emph{not} be linear here. (Moral: the logarithmic plot is not universal — match the plot to the order being tested.)

\textbf{(d) Excess $\ce{OH-}$ (pseudo-first order):} If $[\ce{OH-}]_0 \gg [\text{ester}]_0$, then $[\ce{OH-}]$ remains effectively constant and
\[ \text{Rate} = k[\ce{OH-}][\text{ester}] = k_{\text{obs}}[\text{ester}], \quad \text{where } k_{\text{obs}} = k[\ce{OH-}]_0. \]
The observed order is \textbf{first order} (with respect to the ester). In this case the conductivity change now tracks the ester concentration, and a plot of $\ln(\kappa_t - \kappa_\infty)$ against $t$ \textbf{is} linear, with slope $= -k_{\text{obs}}$.
\end{exoresolu}

\newpage

\coursec{Exercises}

\courssub{I check my knowledge}

\begin{exercice}[Multiple choice questions]
For each question, choose the single correct answer and justify your choice in one sentence.
\begin{enumerate}
\item The unit of the rate of a chemical reaction is:
\begin{enumerate}
\item $\mathrm{mol\,dm^{-3}}$
\item $\mathrm{mol\,dm^{-3}\,s^{-1}}$
\item $\mathrm{dm^{3}\,mol^{-1}\,s^{-1}}$
\item $\mathrm{s^{-1}}$
\end{enumerate}
\item For a reaction which is first order in \ce{A} and second order in \ce{B}, the unit of the rate constant $k$ is:
\begin{enumerate}
\item $\mathrm{s^{-1}}$
\item $\mathrm{dm^{3}\,mol^{-1}\,s^{-1}}$
\item $\mathrm{dm^{6}\,mol^{-2}\,s^{-1}}$
\item $\mathrm{mol\,dm^{-3}\,s^{-1}}$
\end{enumerate}
\item The half-life of a first order reaction:
\begin{enumerate}
\item depends on the initial concentration of the reactant
\item is independent of the initial concentration of the reactant
\item doubles when the temperature doubles
\item is always equal to the rate constant
\end{enumerate}
\item A catalyst increases the rate of a reaction because it:
\begin{enumerate}
\item shifts the position of equilibrium towards the products
\item increases the average kinetic energy of the particles
\item provides an alternative reaction path of lower activation energy
\item increases the enthalpy change of the reaction
\end{enumerate}
\end{enumerate}
\end{exercice}

\begin{exercice}[True or false]
State whether each statement is true or false, and justify your answer briefly.
\begin{enumerate}
\item The rate of a reaction always increases when the concentration of a reactant increases.
\item The instantaneous rate of a reaction at time $t$ is obtained from the gradient of the tangent to the concentration--time curve at that time.
\item In the titrimetric method of following a reaction, samples withdrawn from the mixture must be quenched before titration.
\item According to collision theory, every collision between reactant particles leads to the formation of products.
\item The Maxwell--Boltzmann distribution shows that raising the temperature increases the proportion of particles with energy greater than the activation energy.
\end{enumerate}
\end{exercice}

\begin{exercice}[Definitions]
Define each of the following terms precisely, as you would in the GCE Advanced Level examination:
\begin{enumerate}
\item rate of reaction;
\item order of reaction with respect to a reactant (partial order);
\item overall order of reaction;
\item rate constant;
\item half-life of a reaction;
\item activation energy;
\item activated complex (transition state).
\end{enumerate}
\end{exercice}

\begin{exercice}[Direct calculations]
\begin{enumerate}
\item In the decomposition of hydrogen peroxide, $\ce{2H2O2(aq) -> 2H2O(l) + O2(g)}$, the concentration of \ce{H2O2} falls from $\SI{0.80}{mol\,dm^{-3}}$ to $\SI{0.50}{mol\,dm^{-3}}$ in $\SI{100}{s}$. Calculate the average rate of disappearance of \ce{H2O2} over this interval.
\item A first order reaction has a rate constant $k = \SI{2.31e-3}{s^{-1}}$. Calculate its half-life.
\item For the reaction $\ce{A + B -> products}$, the rate equation is $\text{rate} = k[\ce{A}][\ce{B}]$. Given $k = \SI{4.0e-2}{dm^{3}\,mol^{-1}\,s^{-1}}$, $[\ce{A}] = \SI{0.20}{mol\,dm^{-3}}$ and $[\ce{B}] = \SI{0.10}{mol\,dm^{-3}}$, calculate the rate of the reaction.
\item The rate constant of a reaction is $\SI{1.5e-4}{dm^{6}\,mol^{-2}\,s^{-1}}$. Deduce the overall order of the reaction.
\end{enumerate}
\end{exercice}

\courssub{I practise}

\begin{exercice}[Average, instantaneous and initial rates]
The decomposition of dinitrogen pentoxide in tetrachloromethane,
\[
\ce{2N2O5 -> 4NO2 + O2},
\]
was followed by measuring the concentration of \ce{N2O5} at different times:
\begin{center}
\begin{tabular}{|l|c|c|c|c|c|c|c|}
\hline
\rowcolor{popblueL} $t$ / s & 0 & 200 & 400 & 600 & 800 & 1000 & 1200 \\
\hline $[\ce{N2O5}]$ / $\mathrm{mol\,dm^{-3}}$ & 1.00 & 0.78 & 0.61 & 0.47 & 0.37 & 0.29 & 0.22 \\
\hline
\end{tabular}
\end{center}
\begin{enumerate}
\item Plot the graph of $[\ce{N2O5}]$ against time.
\item Calculate the average rate of disappearance of \ce{N2O5} between $t = 0$ s and $t = 400$ s, and between $t = 400$ s and $t = 800$ s. Comment on the values obtained.
\item By drawing tangents, estimate the instantaneous rate at $t = 0$ s (initial rate) and at $t = 600$ s.
\item Show that the ratio of the two instantaneous rates is approximately equal to the ratio of the corresponding concentrations, and state what this suggests about the order of the reaction.
\end{enumerate}
\end{exercice}

\begin{exercice}[Initial rate method]
The reaction $\ce{A + 2B -> products}$ was studied at $\SI{298}{K}$ and the following initial rates were measured:
\begin{center}
\begin{tabular}{|c|c|c|c|}
\hline
\rowcolor{popblueL} Experiment & $[\ce{A}]_0$ / $\mathrm{mol\,dm^{-3}}$ & $[\ce{B}]_0$ / $\mathrm{mol\,dm^{-3}}$ & Initial rate / $\mathrm{mol\,dm^{-3}\,s^{-1}}$ \\
\hline 1 & 0.10 & 0.10 & $2.0\times 10^{-4}$ \\
\hline 2 & 0.20 & 0.10 & $4.0\times 10^{-4}$ \\
\hline 3 & 0.10 & 0.20 & $1.6\times 10^{-3}$ \\
\hline 4 & 0.20 & 0.20 & $3.2\times 10^{-3}$ \\
\hline
\end{tabular}
\end{center}
\begin{enumerate}
\item Using the ratio method, determine the order of the reaction with respect to \ce{A} and with respect to \ce{B}.
\item Write the rate equation and state the overall order.
\item Calculate the value of the rate constant $k$, giving its units.
\item Predict the initial rate when $[\ce{A}]_0 = \SI{0.30}{mol\,dm^{-3}}$ and $[\ce{B}]_0 = \SI{0.15}{mol\,dm^{-3}}$.
\end{enumerate}
\end{exercice}

\begin{attention}[Check your data before applying the ratio method]
In an initial-rate problem, first verify that the experiments are mutually consistent: the rate predicted for any row from the rate law must match the measured value. If doubling a concentration multiplies the rate by $r$, the partial order $x$ satisfies $2^{x} = r$ — and $x$ should normally be a small whole number (0, 1 or 2).
\end{attention}

\begin{exercice}[Arrhenius plot]
The rate constant $k$ of a first order reaction was measured at several temperatures:
\begin{center}
\begin{tabular}{|l|c|c|c|c|c|}
\hline
\rowcolor{popblueL} $T$ / K & 600 & 650 & 700 & 750 & 800 \\
\hline $k$ / $\mathrm{s^{-1}}$ & $2.8\times 10^{-4}$ & $1.6\times 10^{-3}$ & $7.5\times 10^{-3}$ & $3.0\times 10^{-2}$ & $1.0\times 10^{-1}$ \\
\hline
\end{tabular}
\end{center}
\begin{enumerate}
\item Starting from the Arrhenius equation $k = A\,e^{-E_a/(RT)}$, show that a graph of $\ln k$ against $\dfrac{1}{T}$ should be a straight line, and give the expressions of its gradient and its intercept.
\item Complete a table of values of $\dfrac{1}{T}$ and $\ln k$, then plot the Arrhenius graph.
\item From the gradient, determine the activation energy $E_a$ in $\mathrm{kJ\,mol^{-1}}$ ($R = \SI{8.31}{J\,K^{-1}\,mol^{-1}}$).
\item Determine the pre-exponential factor $A$ and comment on the linearity of the graph.
\end{enumerate}
\end{exercice}

\begin{exercice}[Conductimetric following of an ester hydrolysis]
The alkaline hydrolysis of an ester,
\[
\ce{CH3COOC2H5 + OH- -> CH3COO- + C2H5OH},
\]
is followed by measuring the conductivity $\sigma$ of the mixture as a function of time, the ester and sodium hydroxide being mixed in equal initial concentrations $c_0$.
\begin{enumerate}
\item Explain qualitatively why the conductivity of the mixture decreases as the reaction proceeds.
\item Sketch the expected shape of the curve $\sigma = f(t)$.
\item The reaction is found to be first order in each reactant. Write the rate equation and show that, since the two initial concentrations are equal, the integrated rate law reduces to
\[
\frac{1}{c} = \frac{1}{c_0} + kt.
\]
\item Given that $\sigma$ is proportional to the concentration of the conducting ions, explain how a suitable linear plot allows $k$ to be determined.
\item In one experiment with $c_0 = \SI{0.050}{mol\,dm^{-3}}$, the plot gave a gradient of $\SI{0.12}{dm^{3}\,mol^{-1}\,min^{-1}}$. Deduce the value of $k$ and state its units.
\end{enumerate}
\end{exercice}

\courssub{I prepare for the GCE}

\begin{exercice}[Kinetics of the oxidation of iodide ions (GCE type)]
The reaction between hydrogen peroxide and iodide ions in acidic solution,
\[
\ce{H2O2(aq) + 2I-(aq) + 2H+(aq) -> I2(aq) + 2H2O(l)},
\]
is studied at $\SI{298}{K}$. The iodine liberated is determined by withdrawing $\SI{10.0}{cm^{3}}$ samples at known times, quenching them in ice-cold water, and titrating with standard sodium thiosulfate solution:
\[
\ce{I2(aq) + 2S2O3^{2-}(aq) -> 2I-(aq) + S4O6^{2-}(aq)}.
\]
\begin{enumerate}
\item Explain why the samples are quenched before titration, and state two other methods by which quenching could be achieved. \hfill (3 marks)
\item Name the indicator used in the thiosulfate titration and state the colour change at the end point. \hfill (2 marks)
\item The following titres of $\SI{0.0100}{mol\,dm^{-3}}$ thiosulfate were obtained:
\begin{center}
\begin{tabular}{|l|c|c|c|c|c|c|}
\hline
\rowcolor{popblueL} $t$ / min & 0 & 5 & 10 & 15 & 20 & 30 \\
\hline $V(\ce{S2O3^{2-}})$ / $\mathrm{cm^{3}}$ & 0.0 & 4.2 & 7.6 & 10.4 & 12.6 & 15.4 \\
\hline
\end{tabular}
\end{center}
Calculate the concentration of \ce{I2} in the reaction mixture at each time and plot $[\ce{I2}]$ against $t$. \hfill (5 marks)
\item From your graph, determine the initial rate of formation of \ce{I2} and the instantaneous rate at $t = 15$ min. \hfill (4 marks)
\item Separate initial-rate experiments gave: doubling $[\ce{H2O2}]$ doubles the rate; doubling $[\ce{I-}]$ doubles the rate; the rate is independent of $[\ce{H+}]$. Deduce the rate equation, the overall order, and the units of $k$. \hfill (4 marks)
\item Suggest a physical method, other than titrimetry, by which this reaction could be followed, and justify your choice. \hfill (2 marks)
\end{enumerate}
\end{exercice}

\begin{exercice}[First order kinetics and half-life (GCE type)]
The gas-phase decomposition of ethanal,
\[
\ce{CH3CHO(g) -> CH4(g) + CO(g)},
\]
is followed at constant volume by measuring the total pressure in the vessel.
\begin{enumerate}
\item Explain why the total pressure can be used to follow the progress of this reaction. \hfill (2 marks)
\item The reaction is first order. Starting from the rate equation, derive the integrated rate law
\[
\ln\frac{[\ce{CH3CHO}]_0}{[\ce{CH3CHO}]} = kt,
\]
and hence show that the half-life is $t_{1/2} = \dfrac{\ln 2}{k}$. \hfill (5 marks)
\item At $\SI{800}{K}$, the pressure measurements show that the concentration of ethanal falls to one quarter of its initial value after $\SI{420}{s}$. Calculate the rate constant $k$ at this temperature. \hfill (3 marks)
\item At $\SI{700}{K}$, the rate constant is $\SI{4.1e-4}{s^{-1}}$. Using the Arrhenius equation, calculate the activation energy of the reaction in $\mathrm{kJ\,mol^{-1}}$ ($R = \SI{8.31}{J\,K^{-1}\,mol^{-1}}$). \hfill (4 marks)
\item Sketch, on the same axes, the Maxwell--Boltzmann distribution of molecular energies at $\SI{700}{K}$ and at $\SI{800}{K}$, mark the activation energy, and use your sketch to explain why the reaction is faster at the higher temperature. \hfill (4 marks)
\item State two reasons, according to collision theory, why most collisions between ethanal molecules do not lead to reaction. \hfill (2 marks)
\end{enumerate}
\end{exercice}

\begin{exercice}[Energy profiles and catalysis (GCE essay type)]
\begin{enumerate}
\item Sketch, on the same diagram, the potential energy profile for an exothermic reaction occurring (i) without a catalyst and (ii) with a catalyst. Label clearly the activation energy $E_a$ of each path, the enthalpy change of reaction $\Delta H$, and the activated complex. \hfill (5 marks)
\item Using transition state theory, describe how the activated complex is formed and explain the role it plays in the conversion of reactants into products. \hfill (4 marks)
\item State two similarities and two differences between collision theory and transition state theory. \hfill (4 marks)
\item Distinguish between homogeneous and heterogeneous catalysis, giving one named industrial example of each, with a balanced equation for the reaction catalysed. \hfill (4 marks)
\item Explain, with reference to the steps of adsorption, reaction and desorption, how a heterogeneous catalyst such as iron in the Haber process lowers the activation energy, and explain why the catalyst can become poisoned. \hfill (3 marks)
\end{enumerate}
\end{exercice}



\newpage

\coursec{Solutions to the exercises}

\coursec{Corrected Exercises: Reaction Kinetics}

\begin{corrige}[Exercise 1 — Multiple choice questions]
\begin{enumerate}
\item \textbf{Answer (b).} A rate is a change of concentration per unit time, so its unit is $\mathrm{mol\,dm^{-3}\,s^{-1}}$; option (a) is a concentration, (d) is the unit of a first order rate constant, and (c) is a second order $k$ unit.
\item \textbf{Answer (c).} The overall order is $1 + 2 = 3$; from $k = \dfrac{\text{rate}}{[\ce{A}][\ce{B}]^{2}}$ the units are $\dfrac{\mathrm{mol\,dm^{-3}\,s^{-1}}}{(\mathrm{mol\,dm^{-3}})^{3}} = \mathrm{dm^{6}\,mol^{-2}\,s^{-1}}$.
\item \textbf{Answer (b).} Since $t_{1/2} = \dfrac{\ln 2}{k}$ contains no concentration term, the half-life of a first order reaction is independent of the initial concentration.
\item \textbf{Answer (c).} A catalyst provides an alternative reaction path of lower activation energy; it changes neither $\Delta H$ nor the position of equilibrium, and it does not raise the temperature of the particles.
\end{enumerate}
\end{corrige}

\begin{corrige}[Exercise 2 — True or false]
\begin{enumerate}
\item \textbf{False.} This is only true when the partial order with respect to that reactant is not zero; for a zero order reactant, the rate is independent of its concentration.
\item \textbf{True.} The instantaneous rate is $\left|\dfrac{d[\ce{X}]}{dt}\right|$, which is precisely the gradient of the tangent drawn to the concentration--time curve at time $t$.
\item \textbf{True.} Quenching (rapid cooling, dilution, or chemical removal of a reactant) stops the reaction in the sample so that the titre reflects the composition at the exact moment of withdrawal.
\item \textbf{False.} Only \cle{effective collisions} — those with energy at least equal to the activation energy \emph{and} with the correct orientation — lead to products.
\item \textbf{True.} Raising the temperature flattens and shifts the Maxwell--Boltzmann distribution towards higher energies, so the area under the curve beyond $E_a$ (the fraction of particles able to react) increases.
\end{enumerate}
\end{corrige}

\begin{corrige}[Exercise 3 — Definitions]
\begin{enumerate}
\item \textbf{Rate of reaction:} the change in concentration of a reactant or product per unit time, e.g.\ $\text{rate} = -\dfrac{d[\ce{A}]}{dt}$ for a reactant \ce{A}; unit $\mathrm{mol\,dm^{-3}\,s^{-1}}$.
\item \textbf{Partial order:} the power to which the concentration of that reactant is raised in the experimentally determined rate equation.
\item \textbf{Overall order:} the sum of the partial orders of all the species appearing in the rate equation.
\item \textbf{Rate constant:} the proportionality constant $k$ in the rate equation; numerically equal to the rate when every reactant concentration is $\SI{1}{mol\,dm^{-3}}$, and constant at a given temperature.
\item \textbf{Half-life:} the time taken for the concentration of a reactant to fall to half of its initial value.
\item \textbf{Activation energy:} the minimum energy that colliding particles must possess for a collision to result in reaction.
\item \textbf{Activated complex (transition state):} the unstable, high-energy arrangement of atoms at the maximum of the energy profile, in which old bonds are partially broken and new bonds partially formed; it can decompose either to products or back to reactants.
\end{enumerate}
\end{corrige}

\begin{corrige}[Exercise 4 — Direct calculations]
\begin{enumerate}
\item Average rate $= -\dfrac{\Delta[\ce{H2O2}]}{\Delta t} = \dfrac{0.80 - 0.50}{100} = \SI{3.0e-3}{mol\,dm^{-3}\,s^{-1}}$.
\item $t_{1/2} = \dfrac{\ln 2}{k} = \dfrac{0.693}{2.31\times 10^{-3}} = \SI{300}{s}$.
\item Rate $= k[\ce{A}][\ce{B}] = 4.0\times 10^{-2} \times 0.20 \times 0.10 = \SI{8.0e-4}{mol\,dm^{-3}\,s^{-1}}$.
\item For overall order $n$, the units of $k$ are $\mathrm{dm^{3(n-1)}\,mol^{-(n-1)}\,s^{-1}}$. Here $n - 1 = 2$, so $n = 3$: the reaction is \textbf{third order overall}.
\end{enumerate}
\end{corrige}

\begin{corrige}[Exercise 5 — Average, instantaneous and initial rates]
\begin{enumerate}
\item The graph of $[\ce{N2O5}]$ against $t$ is a smooth decay curve, steep at the start and flattening as the concentration falls.
\item Between 0 and 400 s: rate $= \dfrac{1.00 - 0.61}{400} = \SI{9.8e-4}{mol\,dm^{-3}\,s^{-1}}$. Between 400 and 800 s: rate $= \dfrac{0.61 - 0.37}{400} = \SI{6.0e-4}{mol\,dm^{-3}\,s^{-1}}$. The average rate decreases with time because the concentration of reactant — and hence the frequency of effective collisions — decreases as the reaction proceeds.
\item Tangent at $t = 0$: gradient $\approx \SI{1.0e-3}{mol\,dm^{-3}\,s^{-1}}$ (initial rate). Tangent at $t = 600$ s: gradient $\approx \SI{4.7e-4}{mol\,dm^{-3}\,s^{-1}}$. (Values within $\pm 10\%$ are accepted, depending on the tangent drawn.)
\item $\dfrac{\text{rate}_{0}}{\text{rate}_{600}} \approx \dfrac{1.0\times 10^{-3}}{4.7\times 10^{-4}} \approx 2.1$ and $\dfrac{[\ce{N2O5}]_{0}}{[\ce{N2O5}]_{600}} = \dfrac{1.00}{0.47} \approx 2.1$. Since rate $\propto [\ce{N2O5}]$, the reaction is \textbf{first order} in \ce{N2O5}.
\end{enumerate}
\end{corrige}

\begin{corrige}[Exercise 6 — Initial rate method]
\textbf{Note on the data.} Comparing experiments 1 and 3 ($[\ce{A}]$ constant, $[\ce{B}]$ doubled): the rate is multiplied by $\dfrac{1.6\times 10^{-3}}{2.0\times 10^{-4}} = 8 = 2^{3}$, which would give order 3 in \ce{B}; but experiments 2 and 4 then give $\dfrac{3.2\times 10^{-3}}{4.0\times 10^{-4}} = 8$ for doubling \emph{both} concentrations, which is inconsistent ($2^{1}\times 2^{3} = 16 \neq 8$). The table as printed is therefore internally inconsistent. The intended, consistent data are: experiment 3, rate $= 8.0\times 10^{-4}$; experiment 4, rate $= 1.6\times 10^{-3}\ \mathrm{mol\,dm^{-3}\,s^{-1}}$. These corrected values are used below.
\begin{enumerate}
\item Experiments 1 and 2 ($[\ce{B}]$ constant): doubling $[\ce{A}]$ doubles the rate, so $2^{x} = 2$, giving \textbf{order 1 in \ce{A}}. Experiments 1 and 3 ($[\ce{A}]$ constant): doubling $[\ce{B}]$ multiplies the rate by $\dfrac{8.0\times 10^{-4}}{2.0\times 10^{-4}} = 4 = 2^{2}$, so \textbf{order 2 in \ce{B}}. Check with experiment 4: $2^{1}\times 2^{2} = 8$ and $\dfrac{1.6\times 10^{-3}}{2.0\times 10^{-4}} = 8$ — consistent.
\item Rate equation: $\text{rate} = k[\ce{A}][\ce{B}]^{2}$; overall order $= 1 + 2 = 3$.
\item $k = \dfrac{\text{rate}}{[\ce{A}][\ce{B}]^{2}} = \dfrac{2.0\times 10^{-4}}{0.10 \times (0.10)^{2}} = \SI{0.20}{dm^{6}\,mol^{-2}\,s^{-1}}$ (third order units).
\item Rate $= 0.20 \times 0.30 \times (0.15)^{2} = \SI{1.35e-3}{mol\,dm^{-3}\,s^{-1}}$.
\end{enumerate}
\textbf{Examiner's advice:} always verify that every row of the table satisfies the rate law you propose; if doubling a concentration multiplies the rate by $r$, the partial order $x$ obeys $2^{x} = r$, and $x$ should be a small whole number (0, 1 or 2).
\end{corrige}

\begin{corrige}[Exercise 7 — Arrhenius plot]
\begin{enumerate}
\item Taking natural logarithms of $k = A\,e^{-E_a/(RT)}$:
\[ \ln k = \ln A - \frac{E_a}{R}\cdot\frac{1}{T}. \]
This is of the form $y = mx + c$ with $y = \ln k$ and $x = \dfrac{1}{T}$: a straight line of gradient $-\dfrac{E_a}{R}$ and intercept $\ln A$.
\item
\begin{center}
\begin{tabular}{|c|c|c|c|c|c|}
\hline
\rowcolor{popblueL} $\dfrac{1}{T}$ / $\mathrm{K^{-1}}$ ($\times 10^{3}$) & 1.667 & 1.538 & 1.429 & 1.333 & 1.250 \\
\hline $\ln k$ & $-8.18$ & $-6.44$ & $-4.89$ & $-3.51$ & $-2.30$ \\
\hline
\end{tabular}
\end{center}
The five points lie on a good straight line of negative gradient.
\item Gradient $= \dfrac{-2.30 - (-8.18)}{(1.250 - 1.667)\times 10^{-3}} = \dfrac{5.88}{-4.17\times 10^{-4}} \approx \SI{-1.41e4}{K}$. Hence
\[ E_a = 1.41\times 10^{4} \times 8.31 \approx \SI{1.17e5}{J\,mol^{-1}} = \SI{117}{kJ\,mol^{-1}}. \]
\item Intercept: $\ln A \approx -8.18 + 1.41\times 10^{4} \times 1.667\times 10^{-3} \approx 15.3$, so $A \approx e^{15.3} \approx \SI{4.4e6}{s^{-1}}$. The excellent linearity confirms that the reaction obeys the Arrhenius equation over this temperature range, i.e.\ $E_a$ is effectively constant.
\end{enumerate}
\end{corrige}

\begin{corrige}[Exercise 8 — Conductimetric following of an ester hydrolysis]
\begin{enumerate}
\item The highly mobile \ce{OH-} ions (very high molar ionic conductivity) are consumed and replaced by \ce{CH3COO-} ions, which conduct much less; the ester and ethanol are non-ionic. The total ionic conductance therefore falls as the reaction proceeds.
\item $\sigma$ starts high, decreases rapidly at first, then more and more slowly, and levels off at a non-zero final value $\sigma_{\infty}$ due to the \ce{CH3COO-} and \ce{Na+} ions remaining in solution.
\item Rate $= k[\text{ester}][\ce{OH-}]$. Since the initial concentrations are equal and the reactants are consumed in a $1:1$ ratio, $[\text{ester}] = [\ce{OH-}] = c$ at all times, so
\[ -\frac{dc}{dt} = kc^{2} \quad\Longrightarrow\quad -\frac{dc}{c^{2}} = k\,dt. \]
Integrating between $(c_0, 0)$ and $(c, t)$: $\dfrac{1}{c} - \dfrac{1}{c_0} = kt$, i.e.
\[ \frac{1}{c} = \frac{1}{c_0} + kt, \]
the integrated second order rate law.
\item Since $\sigma \propto c$ (the concentration of conducting ions), $\dfrac{1}{\sigma} \propto \dfrac{1}{c}$. A plot of $\dfrac{1}{\sigma}$ against $t$ is therefore a straight line whose gradient is proportional to $k$; calibration with solutions of known concentration converts the gradient into $k$.
\item The gradient of the $\dfrac{1}{c}$ versus $t$ plot \emph{is} $k$, so $k = \SI{0.12}{dm^{3}\,mol^{-1}\,min^{-1}}$ (second order units), equivalently $\SI{2.0e-3}{dm^{3}\,mol^{-1}\,s^{-1}}$.
\end{enumerate}
\end{corrige}

\begin{corrige}[Exercise 9 — Kinetics of the oxidation of iodide ions (GCE type)]
\begin{enumerate}
\item Quenching stops (or slows to a negligible rate) the reaction in the withdrawn sample, so the titre measures the \ce{I2} present at the exact time of sampling. \textbf{(1)} Other quenching methods: dilution with a large volume of ice-cold water \textbf{(1)}; chemical quenching, e.g.\ neutralising the acid or adding a reagent that rapidly consumes one reactant \textbf{(1)}.
\item Indicator: freshly prepared starch solution, added near the end point \textbf{(1)}; colour change from blue-black to colourless \textbf{(1)}.
\item $n(\ce{S2O3^{2-}}) = 0.0100 \times \dfrac{V}{1000}$; from the stoichiometry, $n(\ce{I2}) = \dfrac{1}{2}n(\ce{S2O3^{2-}})$ in the $\SI{10.0}{cm^{3}}$ sample, so $[\ce{I2}] = \dfrac{n(\ce{I2})}{0.0100\ \mathrm{dm^{3}}}$. \textbf{(2 for method)}
\begin{center}
\begin{tabular}{|c|c|c|c|c|c|c|}
\hline
\rowcolor{popblueL} $t$ / min & 0 & 5 & 10 & 15 & 20 & 30 \\
\hline $[\ce{I2}]$ / $\mathrm{mol\,dm^{-3}}$ ($\times 10^{3}$) & 0.0 & 2.1 & 3.8 & 5.2 & 6.3 & 7.7 \\
\hline
\end{tabular}
\end{center}
\textbf{(2 for the values)} The graph is a curve rising steeply at first and flattening as the reactants are consumed. \textbf{(1 for the plot)}
\item Initial rate (tangent at the origin) $\approx \dfrac{2.1\times 10^{-3}}{5} \approx \SI{4.2e-4}{mol\,dm^{-3}\,min^{-1}}$ \textbf{(2)}; tangent at $t = 15$ min gives $\approx \SI{2.4e-4}{mol\,dm^{-3}\,min^{-1}}$ \textbf{(2)}. (Values consistent with the candidate's own tangents are accepted.)
\item First order in \ce{H2O2}, first order in \ce{I-}, zero order in \ce{H+} \textbf{(2)}: $\text{rate} = k[\ce{H2O2}][\ce{I-}]$; overall order $= 2$ \textbf{(1)}; units of $k$: $\mathrm{dm^{3}\,mol^{-1}\,s^{-1}}$ \textbf{(1)}.
\item Colorimetry \textbf{(1)}: iodine is coloured (yellow-brown, or blue-black with starch), so the absorbance at a suitable wavelength is proportional to $[\ce{I2}]$ and can be recorded continuously without sampling \textbf{(1)}.
\end{enumerate}
\textbf{Examiner's expectations:} correct use of the $1:2$ stoichiometry between \ce{I2} and \ce{S2O3^{2-}}; axes labelled with quantities and units; tangents drawn with a ruler; units quoted for every rate.
\end{corrige}

\begin{corrige}[Exercise 10 — First order kinetics and half-life (GCE type)]
\begin{enumerate}
\item One mole of gaseous reactant gives two moles of gaseous products, so the number of moles of gas — and hence the total pressure at constant volume and temperature — increases as the reaction proceeds; the pressure rise is proportional to the extent of reaction. \textbf{(2)}
\item Rate $= -\dfrac{d[\ce{CH3CHO}]}{dt} = k[\ce{CH3CHO}]$ \textbf{(1)}. Separating variables: $-\dfrac{d[\ce{CH3CHO}]}{[\ce{CH3CHO}]} = k\,dt$ \textbf{(1)}. Integrating from $[\ce{CH3CHO}]_0$ at $t = 0$ to $[\ce{CH3CHO}]$ at time $t$:
\[ \ln[\ce{CH3CHO}]_0 - \ln[\ce{CH3CHO}] = kt \quad\Longrightarrow\quad \ln\frac{[\ce{CH3CHO}]_0}{[\ce{CH3CHO}]} = kt \quad \textbf{(1)}. \]
At $t = t_{1/2}$, $[\ce{CH3CHO}] = \dfrac{[\ce{CH3CHO}]_0}{2}$ \textbf{(1)}, so $\ln 2 = k\,t_{1/2}$ and $t_{1/2} = \dfrac{\ln 2}{k}$ \textbf{(1)}.
\item Falling to one quarter means two half-lives: $t_{1/2} = \SI{210}{s}$ \textbf{(1)}, so $k = \dfrac{0.693}{210} = \SI{3.3e-3}{s^{-1}}$ \textbf{(2)}.
\item $\ln\dfrac{k_{800}}{k_{700}} = \dfrac{E_a}{R}\left(\dfrac{1}{700} - \dfrac{1}{800}\right)$ \textbf{(1)}. $\ln\dfrac{3.3\times 10^{-3}}{4.1\times 10^{-4}} = \ln 8.05 = 2.09$ \textbf{(1)}; $\dfrac{1}{700} - \dfrac{1}{800} = 1.79\times 10^{-4}\ \mathrm{K^{-1}}$ \textbf{(1)}. Hence
\[ E_a = \frac{2.09 \times 8.31}{1.79\times 10^{-4}} \approx \SI{9.7e4}{J\,mol^{-1}} = \SI{97}{kJ\,mol^{-1}} \quad \textbf{(1)}. \]
\item Both curves start at the origin, rise to a maximum and fall asymptotically; the $\SI{800}{K}$ curve is lower, flatter and shifted towards higher energies \textbf{(2)}. With $E_a$ marked on the energy axis, the area under the curve beyond $E_a$ — the fraction of molecules energetic enough to react on collision — is markedly larger at $\SI{800}{K}$, so the frequency of effective collisions, and hence the rate, is greater \textbf{(2)}.
\item (i) Most colliding molecules have kinetic energy less than the activation energy \textbf{(1)}; (ii) many collisions occur with an unsuitable orientation for bond breaking and bond making \textbf{(1)}.
\end{enumerate}
\textbf{Examiner's expectations:} a clear derivation with correct limits of integration; the Arrhenius ratio form used correctly with temperatures in kelvin; Maxwell--Boltzmann curves that both start at the origin and never touch the axis again.
\end{corrige}

\begin{corrige}[Exercise 11 — Energy profiles and catalysis (GCE essay type)]
\begin{enumerate}
\item The diagram shows potential energy (vertical axis) against reaction coordinate: reactants at one level, products at a lower level ($\Delta H < 0$) \textbf{(1)}. The uncatalysed path rises to a single high peak (the activated complex) with activation energy $E_a$ \textbf{(2)}; the catalysed path passes through one or more lower peaks with activation energy $E_a' < E_a$ \textbf{(1)}; $\Delta H$ is identical for both paths \textbf{(1)}.
\item As reactant particles approach with sufficient energy, old bonds begin to break while new bonds begin to form; at the top of the energy barrier the atoms adopt a partially bonded, unstable configuration called the activated complex (transition state) \textbf{(2)}. It may either fall apart back into reactants or proceed forward to products; the rate of reaction is governed by the concentration of activated complexes and the frequency with which they cross the barrier \textbf{(2)}.
\item \emph{Similarities} \textbf{(2)}: both require a minimum energy (the activation energy) for reaction; both explain the increase of rate with temperature through the greater number of species able to cross the energy barrier. \emph{Differences} \textbf{(2)}: collision theory treats reaction as occurring at the instant of an effective collision and uses the Maxwell--Boltzmann distribution, whereas transition state theory postulates a definite, short-lived intermediate species (the activated complex) at the energy maximum; transition state theory accounts for orientation through the structure of the complex (entropy of activation), while collision theory introduces orientation only through a steric factor.
\item \emph{Homogeneous catalysis:} catalyst and reactants in the same phase \textbf{(1)} — e.g.\ the acid-catalysed hydrolysis of an ester, $\ce{CH3COOC2H5 + H2O ->[H+] CH3COOH + C2H5OH}$ \textbf{(1)}. \emph{Heterogeneous catalysis:} catalyst in a different phase from the reactants \textbf{(1)} — e.g.\ the Haber process, $\ce{N2(g) + 3H2(g) ->[Fe] 2NH3(g)}$, solid iron catalysing gaseous reactants \textbf{(1)}. (Also accepted: \ce{V2O5} in the Contact process, $\ce{2SO2(g) + O2(g) -> 2SO3(g)}$; Ni in the hydrogenation of alkenes.)
\item \emph{Adsorption:} \ce{N2} and \ce{H2} molecules are adsorbed onto the iron surface, which weakens the strong \ce{N#N} and \ce{H-H} bonds; \emph{reaction:} the adsorbed atoms are held close together in favourable orientations and combine stepwise to \ce{NH3} along a path of lower activation energy; \emph{desorption:} the ammonia molecules leave the surface, freeing the active sites \textbf{(2)}. \emph{Poisoning:} impurities such as sulfur compounds or carbon monoxide adsorb strongly and irreversibly on the active sites, blocking them and destroying the catalytic activity \textbf{(1)}.
\end{enumerate}
\textbf{Examiner's expectations:} a clearly labelled energy profile with both paths and an unchanged $\Delta H$; precise vocabulary (adsorption, active sites, desorption, poison); balanced equations with the catalyst written over the arrow.
\end{corrige}

\newpage

\coursec{Summary sheet}

\begin{popfigure}
\begin{tikzpicture}[
  mindmap, concept color=popblue,
  root concept/.append style={font=\bfseries\small, text=white, minimum size=3.2cm},
  level 1 concept/.append style={sibling angle=60, level distance=4.6cm, minimum size=2.6cm, font=\scriptsize\bfseries, text=white},
  every node/.append style={align=center}]
\node[concept, root concept, align=center]{Reaction\\Kinetics}
  child[grow=90,  concept color=poporange] { node[concept] {Foundations} 
    [clockwise from=150, level 2 concept/.append style={sibling angle=45, level distance=3.2cm, minimum size=2.1cm, font=\tiny, text=white}]
    child { node[concept] {Rate $v=-\frac{1}{\nu_i}\frac{d[i]}{dt}$} }
    child { node[concept] {Factors: conc., $T$, surface, catalyst, light} }
  }
  child[grow=30,  concept color=popgreen] { node[concept, align=center]{Rate from\\curves}
    [clockwise from=90, level 2 concept/.append style={sibling angle=45, level distance=3.2cm, minimum size=2.1cm, font=\tiny, text=white}]
    child { node[concept] {Average rate} }
    child { node[concept] {Tangent $\to$ instantaneous} }
    child { node[concept] {Initial rate at $t=0$} }
  }
  child[grow=-30, concept color=poppurple] { node[concept, align=center]{Measuring\\rates}
    [clockwise from=30, level 2 concept/.append style={sibling angle=45, level distance=3.2cm, minimum size=2.1cm, font=\tiny, text=white}]
    child { node[concept] {Titration (quench)} }
    child { node[concept] {Colorimetry (Beer-Lambert)} }
    child { node[concept] {Conductimetry} }
    child { node[concept] {Gas syringe, pressure} }
    child { node[concept] {Dilatometry} }
  }
  child[grow=-90, concept color=popdark] { node[concept] {Rate law}
    [clockwise from=-30, level 2 concept/.append style={sibling angle=45, level distance=3.2cm, minimum size=2.1cm, font=\tiny, text=white}]
    child { node[concept] {$v=k[\text{A}]^m[\text{B}]^n$} }
    child { node[concept] {Order, half-life $t_{1/2}$} }
    child { node[concept] {Initial-rate method} }
    child { node[concept] {Graphical method} }
  }
  child[grow=-150, concept color=popgold] { node[concept] {Theories}
    [clockwise from=-90, level 2 concept/.append style={sibling angle=45, level distance=3.2cm, minimum size=2.1cm, font=\tiny, text=white}]
    child { node[concept] {Collision theory} }
    child { node[concept] {Arrhenius: $k=Ae^{-E_a/RT}$} }
    child { node[concept] {Transition state} }
    child { node[concept] {Maxwell-Boltzmann} }
  }
  child[grow=150, concept color=popink] { node[concept] {Catalysis}
    [clockwise from=210, level 2 concept/.append style={sibling angle=45, level distance=3.2cm, minimum size=2.1cm, font=\tiny, text=white}]
    child { node[concept] {Homogeneous} }
    child { node[concept] {Heterogeneous} }
    child { node[concept] {Lowers $E_a$} }
  };
\end{tikzpicture}
\legende{Mind map of the chapter \emph{Reaction Kinetics}.}
\end{popfigure}

\begin{aretenir}[Reaction Kinetics — Summary Sheet]

\textbf{1. Key definitions}
\begin{itemize}
\item \cle{Rate of reaction} $v$: for a reaction $\ensuremath{\mathrm{\nu_A A + \nu_B B \rightarrow  \nu_C C + \nu_D D}}$, 
$v = -\dfrac{1}{\nu_A}\dfrac{d[\ce{A}]}{dt} = -\dfrac{1}{\nu_B}\dfrac{d[\ce{B}]}{dt} = +\dfrac{1}{\nu_C}\dfrac{d[\ce{C}]}{dt} = +\dfrac{1}{\nu_D}\dfrac{d[\ce{D}]}{dt}$.
Usual unit: $\mathrm{mol\,dm^{-3}\,s^{-1}}$.
\item \cle{Average rate}: gradient of the chord between two points of a concentration–time curve; 
\cle{instantaneous rate}: gradient of the tangent at time $t$; 
\cle{initial rate}: instantaneous rate at $t=0$.
\item \cle{Rate law}: experimentally determined equation $v = k[\ce{A}]^{m}[\ce{B}]^{n}$.
\item \cle{Rate constant} $k$: proportionality constant, dependent only on temperature (and catalyst).
\item \cle{Order} with respect to A $= m$; \cle{overall order} $= m+n$. Orders are experimental, \emph{not} stoichiometric coefficients; they can be fractional or negative.
\item \cle{Half-life} $t_{1/2}$: time for the concentration of a reactant to fall to half its initial value.
\item \cle{Activation energy} $E_a$: minimum energy colliding particles must possess to react.
\item \cle{Catalyst}: substance that increases the rate by providing a path of lower $E_a$, without being consumed.
\end{itemize}

\textbf{2. Laws and formulas to know}
\begin{itemize}
\item Rate law: $v = k[\ce{A}]^{m}[\ce{B}]^{n}$; units of $k$: $\mathrm{mol^{1-(m+n)}\,dm^{3(m+n)-3}\,s^{-1}}$.
\item First order: $[\ce{A}]_t = [\ce{A}]_0 e^{-kt}$; $\ln[\ce{A}]_t = \ln[\ce{A}]_0 - kt$; $t_{1/2} = \dfrac{\ln 2}{k} \approx \dfrac{0.693}{k}$ (independent of $[\ce{A}]_0$).
\item Zero order: $[\ce{A}]_t = [\ce{A}]_0 - kt$ (straight line of $[\ce{A}]$ vs $t$); $t_{1/2} = \dfrac{[\ce{A}]_0}{2k}$.
\item Second order (single reactant): $\dfrac{1}{[\ce{A}]_t} = \dfrac{1}{[\ce{A}]_0} + kt$; $t_{1/2} = \dfrac{1}{k[\ce{A}]_0}$.
\item Arrhenius equation: $k = A\,e^{-E_a/(RT)}$, i.e.\ $\ln k = \ln A - \dfrac{E_a}{RT}$; plot of $\ln k$ vs $\dfrac{1}{T}$ gives slope $-\dfrac{E_a}{R}$.
\item Two-temperature form: $\ln\dfrac{k_2}{k_1} = \dfrac{E_a}{R}\left(\dfrac{1}{T_1} - \dfrac{1}{T_2}\right)$.
\end{itemize}

\textbf{3. Methods of measuring rates}
\begin{itemize}
\item \cle{Chemical (titrimetric)}: withdraw samples at known times, \emph{quench} (cool, dilute or remove catalyst), then titrate. Example: \ce{H2O2} decomposition followed by titration with \ce{KMnO4}.
\item \cle{Colorimetry}: absorbance proportional to concentration (Beer–Lambert law $A=\varepsilon cl$); suited to coloured species such as \ce{I2} or \ce{MnO4-}.
\item \cle{Conductimetry}: follows changes in ionic concentration (e.g.\ ester hydrolysis producing \ce{H+}); \cle{dilatometry}: volume change in a dilatometer.
\item \cle{Gas syringe / pressure}: volume of gas evolved vs time (e.g.\ \ce{CaCO3 + HCl}); for gases at constant volume, pressure is proportional to amount of gas.
\item \cle{Initial-rate method}: measure $v_0$ for several experiments varying one concentration at a time; compare ratios to find each partial order.
\item \cle{Graphical method}: test linearity of $[\ce{A}]$ vs $t$ (order 0), $\ln[\ce{A}]$ vs $t$ (order 1), $\dfrac{1}{[\ce{A}]}$ vs $t$ (order 2); constant half-life $\Rightarrow$ first order.
\end{itemize}

\textbf{4. Theories}
\begin{itemize}
\item \cle{Collision theory}: reaction requires collisions with energy $\geq E_a$ \emph{and} correct orientation; rate increases with concentration, temperature (Maxwell–Boltzmann distribution shifts), surface area.
\item \cle{Transition-state theory}: reactants pass through a high-energy \cle{activated complex} at the top of the energy profile; $\Delta H = E_a(\text{forward}) - E_a(\text{reverse})$ for elementary steps.
\end{itemize}

\textbf{5. Catalysis}
\begin{itemize}
\item \cle{Homogeneous}: catalyst and reactants in the same phase (e.g.\ \ce{Fe^{2+}} catalysing \ce{S2O8^{2-}}/\ce{I-}).
\item \cle{Heterogeneous}: different phases; reactants adsorb onto the surface, react, then desorb (e.g.\ Fe in the Haber process, Ni in hydrogenation).
\item A catalyst lowers $E_a$, is unchanged at the end, and does \emph{not} shift the position of equilibrium.
\end{itemize}

\textbf{6. Common mistakes}
\begin{itemize}
\item Confusing order with stoichiometric coefficients — orders come only from experiment.
\item Forgetting that units of $k$ depend on the overall order.
\item Using $t_{1/2} = \ln 2/k$ for reactions that are not first order.
\item Forgetting to quench samples in titrimetric methods (reaction continues!).
\item Saying a catalyst "increases the yield" — it only increases the rate.
\item Mixing units: convert $\mathrm{cm^3}$ to $\mathrm{dm^3}$, minutes to seconds when required.
\end{itemize}

\textbf{7. GCE tips}
\begin{itemize}
\item Always draw the tangent carefully and show the gradient triangle when reading rates from curves.
\item In initial-rate problems, tabulate experiments and divide rates where one concentration is constant.
\item State units for every quantity; give $k$ with its correct units.
\item For Arrhenius plots, quote the slope $= -E_a/R$ and compute $E_a$ in $\mathrm{kJ\,mol^{-1}}$.
\end{itemize}
\end{aretenir}

\findecours

\end{document}
